\documentclass[a4paper,11pt]{ltjsarticle}
\usepackage[margin=22mm,headheight=16pt]{geometry}
\usepackage{luatexja-fontspec}
\setmainfont{Latin Modern Roman}
\setsansfont{Latin Modern Sans}
\setmonofont{Latin Modern Mono}[Scale=.88]
\setmainjfont{HaranoAjiMincho-Regular.otf}[BoldFont=HaranoAjiMincho-Bold.otf]
\setsansjfont{HaranoAjiGothic-Regular.otf}[BoldFont=HaranoAjiGothic-Bold.otf]
\usepackage{amsmath,amssymb,bm,mathtools,fix-cm}
\usepackage{booktabs,array,longtable,tabularx}
\usepackage{enumitem}
\usepackage[dvipsnames]{xcolor}
\usepackage{xurl,hyperref}
\usepackage{fancyhdr}
\hypersetup{colorlinks=true,linkcolor=MidnightBlue,citecolor=MidnightBlue,urlcolor=MidnightBlue,
 pdftitle={Class 5 / Class 6 Landau--Lifshitz 系の量子--古典 Lax 対応：研究経緯を含む完全ノート},
 pdfauthor={研究用完全計算ノート}}
\pagestyle{fancy}\fancyhf{}
\fancyhead[L]{\small Class 5 / Class 6 量子--古典 Lax 対応：完全ノート}
\fancyhead[R]{\small 2026年9月13日}\fancyfoot[C]{\thepage}
\fancypagestyle{plain}{\fancyhf{}\fancyfoot[C]{\thepage}\renewcommand{\headrulewidth}{0pt}}
\setlength{\parskip}{.32em}
\setlength{\parindent}{1em}
\setlength{\emergencystretch}{2.5em}
\allowdisplaybreaks[2]
\numberwithin{equation}{section}
\setcounter{tocdepth}{2}
\newcommand{\ii}{\mathrm i}
\newcommand{\one}[1]{\mathbf 1_{#1}}
\newcommand{\Id}{\mathbf 1}
\newcommand{\tr}{\operatorname{tr}}
\newcommand{\End}{\operatorname{End}}
\newcommand{\Svec}{\boldsymbol S}
\newcommand{\bs}{\boldsymbol S}
\newcommand{\vvec}{\boldsymbol v}
\newcommand{\Qmap}{\mathcal Q_{\lambda}}
\newcommand{\comm}[2]{[#1,#2]}
\newcommand{\down}{\downarrow}
\newcommand{\eps}{\epsilon}
\newcommand{\Cstar}{\mathcal C_{\star}}
\newcommand{\vdown}{V^{\downarrow}}
\newcommand{\Pup}{P_{\uparrow}}
\newcommand{\Eplus}{E_{+}}
\newcolumntype{L}[1]{>{\raggedright\arraybackslash}p{#1}}
\newcommand{\status}[1]{\par\noindent\textbf{出典区分：}#1\par}
\newenvironment{historymemo}{\begin{quote}\small\noindent\textbf{研究経緯メモ.}\ }{\end{quote}}
\newenvironment{cautionnote}{\begin{quote}\small\noindent\textbf{注意.}\ }{\end{quote}}
\title{\vspace{-1.2em}\bfseries 対角化不可能な Class 5 / Class 6 Landau--Lifshitz 系の\\量子--古典 Lax 対応\\[2mm]
\large 研究経緯を含む完全計算ノート\\[1mm]
\normalsize ML Lax 探索から off-shell certificate、既知 hierarchy 同定、有限格子 time-Lax 極限、coherent-state 作用素積補正、古典 $r$-matrix 照合まで}
\author{研究用完全計算ノート}
\date{2026年9月13日}

\begin{document}
\maketitle
\thispagestyle{plain}

\begin{abstract}
本ノートは、対角化不可能な可積分 spin chain の Class 5 および Class 6 と、その Landau--Lifshitz 型 continuum limit について、研究開始から最終的な量子--古典 Lax 対応の閉鎖までに行った計算と判断を、外部の作業報告を参照しなくても追える形で保存するための研究ノートである。論文原稿のように最短の論理だけを残すのではなく、最初に何を未解決問題だと考えたか、機械学習による探索がどこで役立ちどこで不十分だったか、fake Lax を排除するため判定条件をどう強めたか、既知 hierarchy の発見により新規性評価をどう修正したか、coherent-state lower symbol の扱いで一度循環的な議論になった箇所をどう独立再導出したか、最後に古典 $r$-matrix / monodromy route とどう照合したかを記録する。

Class 5 では、continuum 運動方程式から時間 Lax 成分を構成し、曲率が運動方程式残差の可逆な線形像になることを off shell に示す。さらに量子 Jordanian Drinfeld twist、$R/L$-operator、有限格子の時間 Lax 作用素、coherent-state 作用素積、continuum $U,V$、古典 $r$-matrix、保存量 $Q_2,Q_3,Q_4$、null-like warped $SL(2)$ Landau--Lifshitz hierarchy との局所辞書、境界 twist を一つの規約体系で接続する。量子有限格子から得た時間 Lax 行列は、独立な古典 $r$-matrix / monodromy 構成から得た行列と、曲率を変えない単位行列成分を除いて一致する。

Class 6 では、三次 nilpotent anisotropy に対する行列平方根型 Lax pair と off-shell certificate を構成し、量子 $R_6$ から有限格子の時間 Lax 作用素を経て continuum $U,V$ を独立に回収する。Class 5 では共有格子点の coherent-state 補正が全微分へ簡単化するが、Class 6 では交換子を含む adjoint covariant derivative が必要になる。この差を正例と failure boundary として用いる。量子 Class 6 が eleven-vertex model であること自体は既知なので、新しい模型や新しい Lax hierarchy の発見とは主張しない。代わりに recent continuum Class 6 と既知 eleven-vertex Landau--Lifshitz hierarchy の局所的な Poisson / EOM / $r$-matrix / Lax 辞書を整理する。

本ノートの目的は次段階の論文草稿に必要な圧縮前の母体を作ることである。したがって、文献入力、規約変換、本ノートでの導出、独立検証、研究上の評価を明示的に区別する。
\end{abstract}

\clearpage
\begingroup\small\setlength{\parskip}{0pt}\tableofcontents\endgroup
\clearpage

\part{研究課題と研究経緯}

\section{このノートの目的と読み方}
\label{sec:purpose}
本ノートは論文原稿ではない。論文では最終結果に不要な探索の失敗、規約の揺れ、source audit、反証計算は通常削るが、ここではそれらも保存する。後に論文草稿を作る際には、本ノートを圧縮し、同じ結論への重複経路を削り、最終的な論理だけを残すことを想定する。

本ノートは二つの並べ方を併用する。本文の数式は、最終的に正しいと確認された\textbf{論理依存}に沿って並べる。一方、各節末の「研究経緯メモ」では、実際の\textbf{発見の時系列}を保存する。これにより「後知恵で最初から答えを知っていた」ような記述を避ける。

式や主張の出典は次の五区分に分ける。
\begin{enumerate}[leftmargin=2.5em]
\item \textbf{文献入力：} 先行論文から採用する $R$-matrix、Hamiltonian、運動方程式、既知 Lax pair など。
\item \textbf{文献式の規約変換：} 文献式を本ノートの基底、時間、spectral parameter、Poisson bracket へ写したもの。
\item \textbf{本ノートで導出：} 文献入力から途中式を明示して再導出した関係。priority claim を意味しない。
\item \textbf{独立検証：} 発見時の helper や既知の古典時間成分を入力せず、別経路・別実装・有限行列計算で再確認した結果。
\item \textbf{研究上の評価：} 文献監査を踏まえた新規性、進歩性、global sector、終了条件についての判断。
\end{enumerate}

未知の記号は初出時に物理的意味とともに宣言する。例えば、量子格子の連続極限で用いる小さい距離を\textbf{格子間隔 $\epsilon$}、monodromy hierarchy を展開する変数を\textbf{生成スペクトルパラメータ $\zeta$}と呼ぶ。同じ文字が別用途で現れる場合は、その節の冒頭で再定義する。

\section{研究開始時の問題設定：Class 5 と Class 6 は何か}
\status{文献入力と研究開始時の評価}
Class 5 と Class 6 は、de Leeuw--Pribytok--Ryan による二状態最近接量子可積分 spin chain の分類に現れる模型である\cite{DPR2019}。ここで「対角化不可能」とは、対象の複素行列が複素数体上で対角化可能でないことをいう。de Leeuw--Fontanella--Nieto Garc\'ia は、この二つを XXX spin chain の対角化不可能な可積分変形として取り上げ、spin coherent state を用いた continuum limit から Landau--Lifshitz 型の場の理論を得た\cite{DFN2026}。

Class 5 の continuum spin 場を三成分複素ベクトル $\boldsymbol S(x,t)$ とし、運動学的拘束 $\boldsymbol S^{\mathsf T}\boldsymbol S=1$ を課す。Class 5 の変形 coupling を $\alpha_5$ と書き、複素反対称行列
\begin{equation}
M=\begin{pmatrix}0&0&-1\\0&0&\ii\\1&-\ii&0\end{pmatrix},\qquad M^{\mathsf T}=-M
\label{eq:intro:M5}
\end{equation}
を用いると、Hamiltonian と運動方程式は
\begin{align}
H_5&=\int dx\left[\frac12\boldsymbol S_x^{\mathsf T}\boldsymbol S_x+
\frac{\alpha_5}{2}\boldsymbol S^{\mathsf T}M\boldsymbol S_x\right],
\label{eq:intro:H5}\\
\boldsymbol S_t&=\boldsymbol S\times(\boldsymbol S_{xx}-\alpha_5M\boldsymbol S_x)
\label{eq:intro:E5}
\end{align}
と書ける\cite{DFN2026}。

Class 6 では、後で使う計算基底に入る前の文献基底の異方性行列を
\begin{equation}
N_{\rm paper}=\begin{pmatrix}0&0&1\\0&0&-\ii\\1&-\ii&0\end{pmatrix}
\label{eq:intro:N6paper}
\end{equation}
と書く。文献は continuum action、Hamiltonian、運動方程式、boost による高次保存量を与えている\cite{DFN2026}。量子 Class 6 が eleven-vertex model であること自体は Corcoran--de Leeuw に明記されており\cite{CdL2024}、eleven-vertex $R$-matrix に対応する Landau--Lifshitz 型 field theory は Levin--Olshanetsky--Zotov により既に研究されている\cite{LOZ2014}。

研究開始時に Class 5 を第一対象とした理由は、2026 年論文が量子 $L$-operator の coherent-state limit から\textbf{空間 Lax 成分}と古典 $r$-matrix を得ながら、運動方程式を再現する\textbf{時間 Lax 成分}（同論文の呼称では companion matrix）を構成できなかったと明記していたためである\cite{DFN2026}。可積分性そのものは量子 Yang--Baxter structure と boost で生成される保存量 hierarchy により別経路で支持されていた。従って「存在するはずの構造を逆問題として探す」という制御された設定だった。

\begin{historymemo}
最初の研究目標は単純だった。「先行研究で未構成だった Class 5/6 の Lax pair を ML と記号計算で見つければよい」と考えた。ただし、Lax pair は gauge transformation、場の再定義、spectral reparametrization による非一意性が大きいので、式が新しく見えること自体を新規性とはみなさない、という警戒は当初から置いた。
\end{historymemo}

\section{研究の時系列：何を考え、何を修正したか}
\label{sec:chronology}
本節では計算の論理順ではなく、研究上の意思決定を時系列に並べる。

\begin{longtable}{L{24mm}L{118mm}}
\toprule
段階 & 起きたことと、その時点での判断\\
\midrule\endhead
出発点 & Krippendorf--L\"ust--Syvaeri の ML Lax search を参考に、Class 5/6 の continuum EOM から Lax 構造を探索する方針を採った\cite{KLS2021}。\\
判定条件の強化 & on-shell flatness loss だけでは恒等的に flat な接続や EOM の一部だけを載せる接続を排除できないと分かった。曲率を EOM 残差の線形像として off shell に因数分解し、係数写像の rank まで見る方針へ移った。\\
Class 5 の解析化 & ML が回収した係数を、null vector と nilpotent matrix を使う局所 ansatz に落とし、解析 Lax pair と可逆な EOM coefficient matrix を得た。\\
Class 6 の解析化 & 三次 nilpotent anisotropy に対して、有限項で打ち切られる行列平方根から Lax pair が得られることを認識した。\\
最初の新規性修正 & 文献監査により、Class 5 は Jordanian / null-like warped Landau--Lifshitz hierarchy、Class 6 は eleven-vertex hierarchy と深く関連することが分かった。「新しい Lax pair」自体を主結果にする方針を撤回した。\\
量子--古典 bridge へ & Class 5 の recent quantum chain から既知 Jordanian hierarchy までの $R/L$、有限格子、continuum $U,V$、$r$-matrix、保存量、境界条件の辞書を主候補に置いた。Class 6 は recent continuum model と eleven-vertex field theory の辞書として位置付け直した。\\
coherent symbol の問題 & 量子有限格子の時間 Lax 成分を continuum 化すると、$\langle AB\rangle$ を $\langle A\rangle\langle B\rangle$ へ因子化した処方では有限差が残ることが分かった。最初は正しい古典 $V$ との差を $\Delta V$ と定義したため、説明として循環的だった。\\
独立再導出 & Class 5 と Class 6 について、既知の古典 $V$ を入力せず、共有格子点上の作用素積を先に評価して補正を導出し直した。Class 6 では Class 5 の全微分型補正が破れ、交換子項が必要なことが分かった。\\
文献位置づけの再修正 & coherent-state symbol の非乗法性、quantum spin-chain の連続極限、古典 $r$-matrix から $V$ を作る方法は既知であることを確認した。主張を「一般原理の発見」から、Class 5/6 の有限格子 time operator から continuum time Lax を直接閉じる具体的 bridge へ限定した。\\
最後の閉鎖 & Class 5 について、古典 $r$-matrix / monodromy から独立に時間 Lax 行列を生成し、量子有限格子 route と場に依存しない単位行列成分を除いて一致することを確認した。\\
現在 & soft particle production の構造的起源は興味深い follow-up だが、現在の quantum--classical Lax bridge を閉じるための必須条件から外した。\\
\bottomrule
\end{longtable}


\part{機械学習探索の完全記録}

\section{この part の目的：機械学習による発見経路の記録}
\label{sec:ml-purpose}
\status{KLS は文献入力。PCM、$S^2$、$T^{1,1}$、Class 5、Class 6 の探索経緯と数値結果は本研究の記録}
本 part は、Class 5 / Class 6 の Lax 構造に到達した機械学習（machine learning; ML）の発見経路を保存する。KLS は Lax equation / zero-curvature equation を最適化問題として扱い、Heisenberg model や principal chiral model で既知構造を回収した\cite{KLS2021}。field theory の探索では、対称性を ansatz に組み込んでいる。

本研究でも ML を最終証明には用いない。ML が担った役割は、有限個の局所 feature の中から必要な tensor structure と係数関係を選び、解析 ansatz を狭めることである。最終的な成立条件は、曲率と運動方程式の off-shell 因数分解、係数行列の可逆性、量子 $R/L$ 構造、Poisson bracket、既知 hierarchy との照合によって別に確認する。

人間が与えた情報と学習で決めた情報を混同しないため、本 part では次の区別を用いる。
\begin{longtable}{L{39mm}L{103mm}}
\toprule
区分 & 内容\\
\midrule\endhead
人間が与えた物理構造 & 模型の EOM、spin constraint、既知の対称性、Class 5 の null direction、Class 6 の nilpotency、許す空間微分次数、有限 feature library。\\
ML が選んだもの & feature の係数、不要 feature の抑制、spectral parameter に沿う係数関係、複数 run で再現される低 residual branch。\\
解析計算へ移したもの & 学習係数の簡単な函数形、曲率の $F=QE$ 因数分解、$Q$ の determinant / inverse、nilpotent matrix square root。\\
ML では決めないもの & 新規性、既知 hierarchy との同一性、量子 $R$-matrix の正しさ、大域的 boundary sector。これらは解析計算と文献監査で判定する。\\
\bottomrule
\end{longtable}

\section{方法開発の前史：PCM、$S^2$、$T^{1,1}$}
\label{sec:ml-prehistory}
Class 5 / Class 6 の探索以前に、二次元非線形 sigma model を用いて ML Lax search の成立条件を検討した。$SU(2)$ principal chiral model (PCM) と対称 coset $S^2=SU(2)/U(1)$ では、局所 current data を入力とする最適化から既知の連続 spectral-parameter family を回収できた。この段階で、単一の Lax connection を一つ当てるよりも、学習された係数の集合が一次元の solution manifold を形成することを確認する方が、spectral parameter family の検出として重要だと認識した。

次に、非対称 coset $T^{1,1}$ の二次元 sigma model では、複数 run で再現する低い flatness loss と compact な block-diagonal candidate が得られた。しかし解析的に曲率を EOM residual と Maurer--Cartan 恒等式へ分解すると、algebraic remainder を消す条件が EOM に掛かる写像も同時に消してしまうことが分かった。EOM residual を $E$、coset 部分に作用する二つの線形写像を $\Phi_\pm$ とすると、曲率中の EOM coefficient は $\frac12(\Phi_- - \Phi_+)$ である。一方、学習された branch を exact に flat にする algebraic cancellation は $\Phi_+=\Phi_-$ を要求し、従って
\begin{equation}
\frac12(\Phi_- - \Phi_+)E=0
\end{equation}
となる。これは flatness が EOM を符号化していないことを意味する。二次元 field theory では fake Lax だった一方、point-particle reduction では genuine mechanical Lax pair が得られた。

この失敗から、低い flatness loss だけを integrability の certificate にしてはいけないことが明確になった。Class 5 / Class 6 では、候補 connection の発見と、その候補が EOM 全体を忠実に符号化することの証明を明示的に分離する方針を採った。

\section{on-shell flatness から off-shell $F=QE$ certificate へ}
\label{sec:ml-offshell}
原コードでは、Class 5 は off-shell の残差を学習し、Class 6 は EOM から時間微分を与えた on-shell 曲率を学習した。両模型の最終的な解析検証には off-shell の因数分解を用いた。本節の loss と独立な時間微分の sampling は Class 5 の学習方法を説明する。

Lax 曲率の独立成分を並べた数値ベクトルを $\boldsymbol f_{\rm curv}$、独立な EOM residual を並べた数値ベクトルを $\boldsymbol e_{\rm EOM}$ とする。空間 Lax ansatz の係数から定まる行列を $Q_{\rm EOM}^{(\mathrm{fit})}$ と書き、基本 loss を
\begin{equation}
\mathcal L_{\rm off}
=
\left\langle
\left\|
\boldsymbol f_{\rm curv}
-Q_{\rm EOM}^{(\mathrm{fit})}\boldsymbol e_{\rm EOM}
\right\|^2
\right\rangle
\label{eq:ml-off-loss}
\end{equation}
とした。ここで $\boldsymbol f_{\rm curv}$ と $\boldsymbol e_{\rm EOM}$ は、局所 tensor / matrix component を学習用に並べたベクトルである。$Q_{\rm EOM}^{(\mathrm{fit})}$ は空間 ansatz と係数を共有し、独立な学習対象にはしていない。

sampling では on-shell trajectory を必要としない。ある一点における spin field $\boldsymbol S$ とその局所微分を生成し、
\begin{equation}
\boldsymbol S^{\mathsf T}\boldsymbol S=1,
\qquad
\boldsymbol S^{\mathsf T}\boldsymbol S_x=0,
\qquad
\boldsymbol S^{\mathsf T}\boldsymbol S_{xx}=-\boldsymbol S_x^{\mathsf T}\boldsymbol S_x
\label{eq:ml-kin-constraints}
\end{equation}
のような運動学的拘束だけを課す。時間微分 $\boldsymbol S_t$ は EOM から計算せず独立に与える。従って loss は EOM を満たす超曲面の外でも評価され、on-shell でだけ起きる cancellation と局所恒等式を区別できる。

さらに、$Q_{\rm EOM}^{(\mathrm{fit})}$ が EOM の独立成分数と同じ column rank を持つことを確認する。最終解析では
\begin{equation}
\boldsymbol f_{\rm curv}=Q_{\rm EOM}\boldsymbol e_{\rm EOM},
\qquad
\operatorname{rank}Q_{\rm EOM}=\dim\boldsymbol e_{\rm EOM}
\label{eq:ml-rank-certificate}
\end{equation}
まで示す。正方行列なら determinant または inverse を明示する。この条件により、恒等的に flat な connection と、EOM の一部しか含まない connection を genuine な Lax representation から区別する。

$S^2=SU(2)/U(1)$ の練習計算では、coset current の $U(1)$ 成分を $J_\pm^{(0)}$、coset 成分を $J_\pm^{(1)}$ として
\begin{equation}
L_+=aJ_+^{(0)}+bJ_+^{(1)},
\qquad
L_-=cJ_-^{(0)}+dJ_-^{(1)}
\end{equation}
を入れると、EOM に乗らない項を消す条件が
\begin{equation}
a=c=1,
\qquad
bd=1
\end{equation}
となる。$b=\lambda$、$d=\lambda^{-1}$ とすれば spectral family を得る。一方、$\lambda=1$ では $L_\pm=J_\pm$ となり Maurer--Cartan 恒等式だけで off shell に flat である。この例も、単一点の flatness と EOM の faithful encoding を分ける必要性を示す。

\section{Class 5：symmetry-informed feature library と係数回収}
\label{sec:ml-class5}
Class 5 の解析式を得る前の探索では、場の再定義後に現れる complex null direction を利用した。null vector を
\begin{equation}
\boldsymbol m_5=(-\ii,-1,0)^{\mathsf T},
\qquad
\boldsymbol m_5^{\mathsf T}\boldsymbol m_5=0
\end{equation}
とし、field-redefined spin を $\boldsymbol S_5^{\rm fr}$ と書く。Class 5 coupling $\alpha_5$ の半分を $c=\alpha_5/2$、continuum spectral parameter を $z_5$ とする。局所 scalar feature を
\begin{equation}
\chi:=\boldsymbol m_5^{\mathsf T}\boldsymbol S_5^{\rm fr},
\qquad
\xi:=\boldsymbol m_5^{\mathsf T}\boldsymbol S_{5,x}^{\rm fr},
\qquad
\varrho:=\boldsymbol m_5^{\mathsf T}
(\boldsymbol S_5^{\rm fr}\times\boldsymbol S_{5,x}^{\rm fr})
\label{eq:ml5-scalars}
\end{equation}
と定義した。

spatial component には
\begin{equation}
\boldsymbol U_{5}^{\rm fr}
=z_5\boldsymbol S_5^{\rm fr}+\beta_5\chi\boldsymbol m_5
\label{eq:ml5-U}
\end{equation}
を用いた。temporal component の feature library には、最終的に必要だった
\begin{equation}
\boldsymbol S_5^{\rm fr}\times\boldsymbol S_{5,x}^{\rm fr},
\qquad
\varrho\boldsymbol m_5,
\qquad
\boldsymbol S_5^{\rm fr},
\qquad
\chi\boldsymbol m_5
\label{eq:ml5-physical-features}
\end{equation}
に加え、不要な候補として
\begin{equation}
\begin{gathered}
\boldsymbol S_{5,x}^{\rm fr},\quad
\chi\boldsymbol S_{5,x}^{\rm fr},\quad
\xi\boldsymbol S_5^{\rm fr},\quad
\boldsymbol m_5\times\boldsymbol S_{5,x}^{\rm fr},\\
\chi(\boldsymbol m_5\times\boldsymbol S_5^{\rm fr}),\quad
\varrho\boldsymbol S_5^{\rm fr},\quad
\boldsymbol m_5,\quad
\chi^2\boldsymbol m_5
\end{gathered}
\label{eq:ml5-distractors}
\end{equation}
の八項を含めた。時間成分の十二係数と $\beta_5$ の計十三個を実数として学習した。null direction、外積、一次空間微分という物理的構造は library の入力である。

EOM residual を $\boldsymbol E_5^{\rm fr}$、曲率の三成分表示を $\boldsymbol F_5^{\rm fr}$ とし、$K_5:=\boldsymbol m_5\boldsymbol m_5^{\mathsf T}$ とすると、主要 loss は
\begin{equation}
\mathcal L_5
=
\left\langle
\left\|
\boldsymbol F_5^{\rm fr}
-\left(z_5\one{3}+\beta_5K_5\right)
\boldsymbol E_5^{\rm fr}
\right\|^2
\right\rangle
+\lambda_{\rm sparse}\mathcal R_{\rm sparse}
\label{eq:ML5loss}
\end{equation}
である。sparsity penalty の重みを $\lambda_{\rm sparse}>0$ とし、各 feature の学習係数を $\theta_A$ と書いて
\begin{equation}
\mathcal R_{\rm sparse}
:=
\sqrt{\beta_5^2+10^{-20}}
+\sum_{A=1}^{12}\sqrt{\theta_A^2+10^{-20}}
\label{eq:ml5-sparse}
\end{equation}
を用いた。罰則は十三係数すべてに一様に掛けた。$10^{-20}$ は絶対値の微分不可能点を滑らかにするための数値 regularization である。

複数の $(c,z_5)$ に対する九回の独立 run で、学習された主要係数は
\begin{equation}
\beta_5=\frac{c^2}{2z_5},
\qquad
\theta_{\boldsymbol S_5^{\rm fr}\times\boldsymbol S_{5,x}^{\rm fr}}=z_5,
\qquad
\theta_{\varrho\boldsymbol m_5}=\frac{c^2}{2z_5},
\qquad
\theta_{\boldsymbol S_5^{\rm fr}}=-z_5^2,
\qquad
\theta_{\chi\boldsymbol m_5}=\frac{c^2}{2}
\label{eq:ml5-relations}
\end{equation}
という簡単な函数形へ整理された。保存済み集計を表\ref{tab:ml5-summary}に示す。
\begin{table}[htbp]
\centering
\caption{Class 5 の symmetry-informed ML による係数回収。validation は学習に用いていない独立標本上の値。}
\label{tab:ml5-summary}
\begin{tabular}{lr}
\toprule
量 & 保存済み結果\\
\midrule
独立 run 数 & 9\\
$\max|\beta_5-c^2/(2z_5)|$ & $1.04\times10^{-9}$\\
physical coefficient の最大誤差 & $1.15\times10^{-7}$\\
validation MSE の最大値 & $7.99\times10^{-15}$\\
validation MSE の中央値 & $8.86\times10^{-20}$\\
不要 feature の最大係数 & $7.17\times10^{-8}$\\
\bottomrule
\end{tabular}
\end{table}

この係数関係を得た後、曲率を解析的に展開して coefficient equations を解き、式\eqref{eq:ml5-relations}が厳密な解であることを確認した。ML は、過剰な feature library から解析すべき branch を選ぶ役割を担った。

当時の保存記録では、集計結果を \path{results/class5_ml_aggregate.json}、再学習 script を \path{code/ml/class5_ml_recovery.py} に保存していた。これらのファイルがなくても本ノートの物理的定義は閉じるよう、式と数値集計を本文にも保存する。

\section{Class 6：nilpotent matrix polynomial library と平方根構造}
\label{sec:ml-class6}
Class 6 では、計算基底での異方性行列を
\begin{equation}
N=
\begin{pmatrix}
0&0&-1\\
0&0&-\ii\\
-1&-\ii&0
\end{pmatrix},
\qquad
N^3=0
\label{eq:ml6-N}
\end{equation}
とする。これは原コードの異方性行列に定数回転 $\operatorname{diag}(1,-1,-1)$ を施した表示であり、同じ回転を spin にも施す。三次 nilpotency を用い、行列 polynomial の有限 library
\begin{equation}
\left\{\one{3},N,N^2\right\}
\end{equation}
の係数を学習した。三つの係数行列を
\begin{align}
A_6^{\rm ML}&=a_0\one{3}+a_1N+a_2N^2,\\
B_6^{\rm ML}&=b_0\one{3}+b_1N+b_2N^2,\\
C_6^{\rm ML}&=c_0\one{3}+c_1N+c_2N^2
\label{eq:ml6-polynomials}
\end{align}
とし、候補 Lax pair を
\begin{equation}
\boldsymbol U_6=A_6^{\rm ML}\boldsymbol S,
\qquad
\boldsymbol V_6=B_6^{\rm ML}\boldsymbol S
+C_6^{\rm ML}(\boldsymbol S\times\boldsymbol S_x)
\label{eq:ml6-UV}
\end{equation}
とした。ここで $a_0$ は連続 family を parametrise する自由係数であり、後に spectral parameter と同定される。各 run では $a_0$ を固定し、残りの八係数を実数として学習した。この探索での時間を $\tau$ とすると、EOM は
\begin{equation}
\partial_\tau\boldsymbol S
=\boldsymbol S\times(\partial_x^2\boldsymbol S-\alpha N\boldsymbol S)
\end{equation}
である。曲率
$\boldsymbol F_6=\partial_\tau\boldsymbol U_6-\partial_x\boldsymbol V_6
+\boldsymbol U_6\times\boldsymbol V_6$
の評価時にはこの EOM から時間微分を与え、三成分の絶対値二乗の標本平均を最小化した。学習時の loss は on-shell flatness である。

六回の独立 run から
\begin{align}
a_1=c_1&=\frac{\alpha}{2a_0},
&
a_2=c_2&=-\frac{\alpha^2}{8a_0^3},\\
b_0&=-a_0^2,
&
b_1&=\frac{\alpha}{2},
&
b_2&=-\frac{3\alpha^2}{8a_0^2},
&
c_0&=a_0
\label{eq:ml6-relations}
\end{align}
が回収された。保存済み validation RMS の最大値は $1.64\times10^{-14}$、解析係数との差の最大値は $4.88\times10^{-15}$ だった。

この係数列を見て $A_6^{\rm ML}$ の二乗を計算すると、$N^3=0$ により
\begin{equation}
\left(A_6^{\rm ML}\right)^2=a_0^2\one{3}+\alpha N
\label{eq:ml6-square}
\end{equation}
となる。従って
\begin{equation}
A_6^{\rm ML}
=a_0\one{3}
+\frac{\alpha}{2a_0}N
-\frac{\alpha^2}{8a_0^3}N^2
\label{eq:ml6-root}
\end{equation}
は、nilpotent matrix square root の二項展開が有限項で打ち切られたものだと分かった。これは ML の係数列を解析構造へ蒸留した段階であり、その後の最終証明では $A_6^2=a_0^2\one{3}+\alpha N$ と曲率の off-shell 因数分解を直接用いる。

当時の保存記録では、集計結果を \path{results/class6_ml_summary.json}、再学習 script を \path{code/ml/class6_ml_recovery.py} に保存していた。

\section{数値 validation、解析式への移行、文献監査}
\label{sec:ml-validation}
ML run の loss が小さいことは、最終的な証明ではない。Class 5 / Class 6 では、候補係数を得た後に次の順で解析計算へ移した。
\begin{enumerate}[leftmargin=2.5em]
\item 学習係数を簡単な $c,z_5,\alpha,a_0$ の函数へ整理する。
\item 曲率を symbolic に再計算し、EOM residual と独立 tensor structure の係数方程式を解く。
\item $F=Q_{\rm EOM}E$ を off shell に示し、$Q_{\rm EOM}$ の determinant または inverse を求める。
\item 量子 $R/L$ data、finite-lattice zero-curvature、Poisson bracket、classical $r$-matrix、保存量との整合性を別実装で確認する。
\item 既知 hierarchy を検索し、候補が既知構造の別表示でないかを監査する。
\end{enumerate}

最終 verification でも on-shell trajectory は用いず、式\eqref{eq:ml-kin-constraints}だけを満たす局所 jet を生成し、$\boldsymbol S_t$ を EOM と独立にサンプルした。代表的な保存済み誤差には、Class 5 full $gl(2)$ off-shell factorization の RMS $1.95\times10^{-16}$、Class 6 off-shell Lax factorization の最大 residual $7.22\times10^{-15}$ 以下などがある。これらは数値安定性の確認であり、解析恒等式を置き換えるものではない。

文献監査により、Class 5 は Jordanian / null-like warped Landau--Lifshitz hierarchy、Class 6 は eleven-vertex Landau--Lifshitz hierarchy と対応することが分かった。この時点で「ML が新しい Lax pair を発見した」という解釈を撤回した。ML の成功は、既知 hierarchy を知らない段階から正しい局所 branch と係数構造へ到達し、その後の quantum--classical dictionary と source audit の入口を与えたことにある。

\begin{historymemo}
$T^{1,1}$ の fake Lax を通じて、低 loss の候補が EOM を符号化するかを別に検証する必要性を認識した。Class 5 では off-shell 残差を探索段階に組み込み、Class 6 では on-shell の探索で得た係数列を nilpotent square root として整理した後に off-shell 因数分解を証明した。ML は局所構造の選択と解析すべき係数関係の抽出に有効だった。
\end{historymemo}

\section{原コードの復元と論文規約への変換}
\label{sec:ml-restored-conventions}
2026 年 9 月に、当時の Class 5 の学習 script と Class 5 / Class 6 の保存 package を再取得した。添付されたファイルと Library 内の原ファイルはバイト単位で一致した。リポジトリ初期化時に作った最小二乗法による再構成版は、Class 5 の候補を八項に減らし疎性罰則を省いており、Class 6 の学習も off-shell に変更していた。今回、実行対象を原コードに基づくものへ置き換え、再構成版の集計は由来が分かる形で archive に移した。

原コード中の場を $\boldsymbol y_5(x,\tau)$、$\boldsymbol y_6(x,\tau)$ とし、論文の場と時間を $\boldsymbol S(x,t)$ とする。定数回転を
\begin{equation}
R_5=\operatorname{diag}(-1,1,-1),\qquad
R_6=\operatorname{diag}(1,-1,-1)
\end{equation}
と定める。論文基底での Class 5 の null vector は
$\boldsymbol m_5^{\rm paper}=(\ii,-1,0)^{\mathsf T}=R_5\boldsymbol m_5$
であり、$M_5^{\rm paper}\boldsymbol v
:=\boldsymbol m_5^{\rm paper}\times\boldsymbol v$ と定義する。
場とパラメータの対応は
\begin{align}
\boldsymbol S(x,t)
&=e^{-\alpha_5xM_5^{\rm paper}/2}R_5\boldsymbol y_5(x,-2t),
&c&=\frac{\alpha_5}{2},
&z_5&=\frac{\ii}{2\lambda},
\label{eq:ml5-paper-dictionary}\\
\boldsymbol S(x,t)
&=R_6\boldsymbol y_6(x,-2t),
&\alpha&=\alpha_6,
&a_0&=\frac{\ii}{2\lambda}
\label{eq:ml6-paper-dictionary}
\end{align}
である。第一式では Class 5、第二式では Class 6 の場を表す。
Class 6 の変換後の異方性行列は式\eqref{eq:ml6-N}である。
Class 5 の原コードの空間微分は、論文変数では
$\mathcal D_x:=\partial_x+(\alpha_5/2)M_5^{\rm paper}$
に対応する。さらに補助空間のベクトル表示に
$O_5:=e^{\alpha_5\lambda M_5^{\rm paper}}$
を作用させると、論文の空間 Lax 行列と一致する。
この回転は複素直交行列であり、$O_5^{\mathsf T}O_5=\one{3}$、$\det O_5=1$ を満たす。

論文の曲率のベクトル成分を $\boldsymbol{\mathcal F}^{(m)}$、EOM residual を $\boldsymbol E^{(m)}$ と書く。
Class 5 の空間係数行列は
\begin{equation}
\mathsf Q_5
=O_5\left(\frac{\ii}{2\lambda}\one{3}
+\beta_5\boldsymbol m_5^{\rm paper}(\boldsymbol m_5^{\rm paper})^{\mathsf T}\right)
\end{equation}
となる。元の正のノルムで評価した loss を保つため、Class 5 では
$-\tfrac12O_5^{-1}(\boldsymbol{\mathcal F}^{(5)}-\mathsf Q_5\boldsymbol E^{(5)})$、
Class 6 では $-\tfrac12\boldsymbol{\mathcal F}^{(6)}|_{\boldsymbol E^{(6)}=0}$
を数値残差に用いた。時間の変換に由来する $-1/2$ と、Class 5 の $O_5^{-1}$ を含めることで、任意の学習係数における元の loss とその勾配が保存される。
行列への写像を $\iota(\boldsymbol v):=-\ii v^i\sigma^i/2$ とし、論文で固定された scalar part も加えると、回収された $U_5,V_5,U_6,V_6$ は論文の行列要素と厳密に一致した。scalar part に追加の学習係数は導入していない。

sampling、初期化、Adam と L-BFGS の設定も原コードから維持した。Class 5 は学習 768 点、検証 4096 点、Adam 600 step、L-BFGS 最大 50 iteration で九回、Class 6 は学習 768 点、検証 3072 点、Adam 1300 step、L-BFGS 最大 120 iteration で六回実行した。Class 5 の疎性罰則の重みは Adam で $10^{-9}$、L-BFGS で $5\times10^{-11}$ である。

Python 3.12、PyTorch 2.10.0 の CPU 実行で、Class 5 の physical coefficient の最大誤差は $3.61\times10^{-7}$、不要項の最大係数は $8.10\times10^{-8}$、validation MSE の最大値は $4.94\times10^{-14}$ だった。Class 6 の最大係数誤差は $4.33\times10^{-12}$、validation RMS の最大値は $1.12\times10^{-11}$ だった。これらは今回の再実行値であり、表\ref{tab:ml5-summary}などの当時の保存値はそのまま保持した。規約変換、任意係数での loss と勾配の一致、解析解での論文行列との一致を検査する実行可能な検証を追加し、二十五項目すべてが通過した。


さらに、Class 6 の標準実行を off-shell の loss に拡張した。
元の on-shell 実行も選択可能であり、ansatz、空間標本、最適化設定は維持した。
$\boldsymbol u_6=\mathsf A_6\boldsymbol S$ の $\mathsf A_6$ は場に依存しないので、
$\boldsymbol{\mathcal F}^{(6)}-\mathsf A_6\boldsymbol E^{(6)}$
中の $\mathsf A_6\partial_t\boldsymbol S$ は相殺する。
従って、この off-shell 残差は任意の学習係数で on-shell 曲率と一致する。
独立な接ベクトルとして時間微分を生成し、原コードとの loss と勾配の一致を検査した。
この変更により探索する係数条件は増えない。
off-shell の六回の再実行では最大係数誤差が $1.12\times10^{-11}$、
validation RMS の最大値が $2.92\times10^{-11}$ となった。
論文補遺ではこの off-shell 設定と結果を採用し、元の on-shell loss との等価性も示した。

\part{共通する数学と量子格子の準備}

\section{固定規約：spin、行列表示、off shell}
Class 5/6 では場と coupling が複素になるので、三成分ベクトル $\boldsymbol x,\boldsymbol y$ の内積は Hermitian 内積ではなく複素双線形な transpose contraction
\begin{equation}
\boldsymbol x\cdot\boldsymbol y:=\boldsymbol x^{\mathsf T}\boldsymbol y
\label{eq:dot}
\end{equation}
とする。外積は Levi--Civita symbol を $\varepsilon_{123}=+1$ として複素線形に拡張する。後で格子間隔に $\epsilon$ を使うため、Levi--Civita symbol は必ず $\varepsilon_{abc}$ と書き分ける。

spin 場は
\begin{equation}
\boldsymbol S^{\mathsf T}\boldsymbol S=1,
\qquad \boldsymbol S^{\mathsf T}\boldsymbol S_x=0,
\qquad \boldsymbol S^{\mathsf T}\boldsymbol S_{xx}=-\boldsymbol S_x^{\mathsf T}\boldsymbol S_x
\label{eq:kin}
\end{equation}
を満たす。式\eqref{eq:kin}は運動学的拘束であり、これを使った恒等式は本ノートでは off shell と呼ぶ。

Pauli 行列を $\sigma_a$ とし、三成分ベクトルから traceless $2\times2$ 行列への写像を
\begin{equation}
\iota(\boldsymbol x):=-\frac{\ii}{2}\boldsymbol x\cdot\boldsymbol\sigma
\label{eq:iota}
\end{equation}
と定義する。この規約では
\begin{equation}
[\iota(\boldsymbol x),\iota(\boldsymbol y)]=\iota(\boldsymbol x\times\boldsymbol y).
\label{eq:iotacomm}
\end{equation}
空間 Lax 行列を $U$、時間 Lax 行列を $V$ と呼び、その曲率を
\begin{equation}
F_{tx}:=\partial_tU-\partial_xV+[U,V]
\label{eq:curv}
\end{equation}
と定義する。

\section{対称 anisotropy を持つ Landau--Lifshitz 方程式の共通恒等式}
\status{一般 anisotropic LL の可積分構造は既知。本節は Class 5/6 に共通な代数を自足的に再導出する}
複素対称な $3\times3$ 行列を\textbf{異方性行列 $\mathcal J$} と呼び、$\mathcal J^{\mathsf T}=\mathcal J$ とする。運動方程式とその残差を
\begin{equation}
\boldsymbol S_t=\boldsymbol S\times(\boldsymbol S_{xx}-\mathcal J\boldsymbol S),\qquad
\boldsymbol E_{\mathcal J}:=\boldsymbol S_t-\boldsymbol S\times(\boldsymbol S_{xx}-\mathcal J\boldsymbol S)
\label{eq:genericLL}
\end{equation}
と定義する。

spectral parameter を $\lambda$ とし、対称行列 $A_\lambda$ を
\begin{equation}
A_\lambda^{\mathsf T}=A_\lambda,\qquad A_\lambda^2=\lambda^2\one3+\mathcal J
\label{eq:Aroot}
\end{equation}
で選ぶ。任意の $3\times3$ 行列 $A$ に対する余因子転置、すなわち adjugate matrix を
\begin{equation}
C_A:=A^2-(\tr A)A+\frac12\left[(\tr A)^2-\tr(A^2)\right]\one3
\label{eq:adj}
\end{equation}
と書く。Cayley--Hamilton より $AC_A=C_AA=(\det A)\one3$ である。時間成分に使う行列を $B_\lambda:=-C_{A_\lambda}$ と定義する。

ベクトル表示の Lax pair を
\begin{equation}
\boldsymbol U_\lambda=A_\lambda\boldsymbol S,
\qquad
\boldsymbol V_\lambda=B_\lambda\boldsymbol S+A_\lambda(\boldsymbol S\times\boldsymbol S_x)
\label{eq:genLax}
\end{equation}
と置く。恒等式
\begin{equation}
(A\boldsymbol x)\times(A\boldsymbol y)=C_A^{\mathsf T}(\boldsymbol x\times\boldsymbol y)
\label{eq:crossadj}
\end{equation}
と spin constraint を用いると、曲率は
\begin{equation}
\partial_t\boldsymbol U_\lambda-\partial_x\boldsymbol V_\lambda+\boldsymbol U_\lambda\times\boldsymbol V_\lambda
=A_\lambda\boldsymbol E_{\mathcal J}
\label{eq:genericFactor}
\end{equation}
と off shell に因数分解する。従って $\det A_\lambda\neq0$ なら zero curvature と EOM は同値である。

Class 5/6 では $\mathcal J$ が nilpotent 行列の有限次数多項式になるので、平方根 $A_\lambda$ の二項展開が有限項で打ち切られる。ML が見つけた係数は、後から見るとこの有限行列平方根の係数だった。

\section{量子格子 Lax 方程式、Sutherland relation、時間 Lax 作用素}
格子 site $n$ の局所量子空間を $\mathcal H_n$、全量子空間を $\mathcal H_q=\bigotimes_n\mathcal H_n$、transfer matrix を作るための補助空間を $\mathcal V_a$ とする。局所 Lax 作用素を $L_{a,n}(u)\in\End(\mathcal V_a\otimes\mathcal H_n)$ と書き、$u$ を\textbf{量子スペクトルパラメータ}と呼ぶ。

regular $R$-matrix は $R(0)=P$ を満たすとする。局所 Hamiltonian density を
\begin{equation}
h_{n-1,n}:=P_{n-1,n}\,\partial_uR_{n-1,n}(u)\big|_{u=0}
\label{eq:hdef}
\end{equation}
と定義する。Yang--Baxter equation を regular point で微分すると Sutherland relation
\begin{equation}
[h_{n-1,n},L_{a,n}L_{a,n-1}]
=L_{a,n}\partial_uL_{a,n-1}-(\partial_uL_{a,n})L_{a,n-1}
\label{eq:SuthMain}
\end{equation}
を得る。量子格子の Heisenberg 時間を $t_{\rm lat}$ とし、$\partial_{t_{\rm lat}}X=\ii[H_{\rm lat},X]$ とする。一つの局所時間 Lax 作用素は
\begin{equation}
\mathcal A_{a,n}=\ii h_{n-1,n}-\ii L_{a,n}^{-1}h_{n-1,n}L_{a,n}-\ii L_{a,n}^{-1}\partial_uL_{a,n}
\label{eq:Afinite}
\end{equation}
と取れ、
\begin{equation}
\ii[H_{\rm lat},L_{a,n}]
=\mathcal A_{a,n+1}L_{a,n}-L_{a,n}\mathcal A_{a,n}
\label{eq:discZC}
\end{equation}
を満たす。

この有限格子の時間成分は、Class 5 原論文で continuum 側に見つからなかった companion matrix と概念的に同じ役割を量子格子上で担う。今回の量子--古典 bridge の中心は、式\eqref{eq:discZC}を coherent state で連続極限へ移し、古典 $V$ を得ることである。

\section{coherent-state lower symbol と共有格子点の作用素積}
各 site の規格化 coherent state を $|\mathfrak c_n\rangle\in\mathcal H_n$、全量子空間の product coherent state を $|\mathfrak C\rangle=\bigotimes_n|\mathfrak c_n\rangle$ とする。補助空間を持つ量子作用素 $X_a$ の\textbf{lower symbol} を
\begin{equation}
X_a^\downarrow:=(\one{\mathcal V_a}\otimes\langle\mathfrak C|)X_a(\one{\mathcal V_a}\otimes|\mathfrak C\rangle)
\label{eq:lowerMain}
\end{equation}
と定義する。結果は補助空間上の行列として残る。

lower symbol は作用素積の準同型ではない。二つの作用素について
\begin{equation}
(X_a^\downarrow\star Y_a^\downarrow):=(X_aY_a)^\downarrow
\label{eq:starMain}
\end{equation}
と star product を定義すれば、一般に
\begin{equation}
X_a^\downarrow\star Y_a^\downarrow\neq X_a^\downarrow Y_a^\downarrow.
\label{eq:nonmult}
\end{equation}
この事実自体は coherent-state / Berezin symbol の標準的な性質であり、新規性ではない\cite{SuzukiTsuchiya2017,Schlichenmaier2010}。

重要なのは、式\eqref{eq:discZC}の積 $\mathcal A_{n+1}L_n$ と $L_n\mathcal A_n$ では、時間 Lax 作用素と空間 Lax 作用素が\textbf{同じ物理格子点}を共有することである。異なる site の product state expectation は因子化できるが、共有 site の作用素積は先に掛けてから expectation value を取る必要がある。Class 5/6 では、この connected な一 site 二次モーメントが、格子間隔 $\epsilon\to0$ の scaling と組み合わさって有限の時間 Lax 補正として残る。


\part{Class 5：continuum Lax、量子 Jordanian chain、量子--古典 bridge}

\section{Class 5 の nilpotent 構造と場の再定義}
Class 5 の行列 $M$ は式\eqref{eq:intro:M5}で定義した。複素 null vector を
\begin{equation}
\boldsymbol m_5:=(-\ii,-1,0)^{\mathsf T}
\label{eq:m5}
\end{equation}
と定義すると、
\begin{equation}
\boldsymbol m_5^{\mathsf T}\boldsymbol m_5=0,
\qquad M\boldsymbol x=\boldsymbol m_5\times\boldsymbol x,
\qquad M^2=\boldsymbol m_5\boldsymbol m_5^{\mathsf T}=:K_5,
\qquad K_5^2=0,
\label{eq:Mnil}
\end{equation}
従って $M^3=0$ である。$K_5$ は rank-one nilpotent matrix であり projector ではない。

Class 5 の変形 coupling $\alpha_5$ の半分を
\begin{equation}
c:=\frac{\alpha_5}{2},\qquad k_5:=c^2=\frac{\alpha_5^2}{4}
\label{eq:c5}
\end{equation}
と定義する。$x$ 依存の場の再定義
\begin{equation}
\boldsymbol S=e^{cxM}\boldsymbol S_5^{\rm fr}
\label{eq:fieldredef5}
\end{equation}
で新しい spin 場 $\boldsymbol S_5^{\rm fr}$ を定義する。上付き ${\rm fr}$ は field redefinition 後という意味で、複素共役や微分を表さない。$M^{\mathsf T}=-M$ なので spin constraint は保たれる。

新しい場で運動方程式は
\begin{equation}
\boldsymbol S_{5,t}^{\rm fr}
=\boldsymbol S_5^{\rm fr}\times\left[
\boldsymbol S_{5,xx}^{\rm fr}-c^2\boldsymbol m_5(\boldsymbol m_5^{\mathsf T}\boldsymbol S_5^{\rm fr})
\right]
\label{eq:E5fr}
\end{equation}
となる。null direction への射影を
\begin{equation}
\chi:=\boldsymbol m_5^{\mathsf T}\boldsymbol S_5^{\rm fr}
\label{eq:chi}
\end{equation}
と書けば、再定義後の Hamiltonian は
\begin{equation}
H_5^{\rm fr}=\int dx\left[
\frac12(\boldsymbol S_{5,x}^{\rm fr})^{\mathsf T}\boldsymbol S_{5,x}^{\rm fr}
+\frac{k_5}{2}\chi^2\right].
\label{eq:H5fr}
\end{equation}

\section{Class 5 continuum Lax pair と off-shell certificate}
Class 5 の continuum spectral parameter を $z_5$ と呼ぶ。さらに、null direction に沿う一階微分を
\begin{equation}
\xi:=\boldsymbol m_5^{\mathsf T}\boldsymbol S_{5,x}^{\rm fr},
\qquad
\varrho:=\boldsymbol m_5^{\mathsf T}(\boldsymbol S_5^{\rm fr}\times\boldsymbol S_{5,x}^{\rm fr})
\label{eq:xirho}
\end{equation}
と定義する。

局所 ansatz
\begin{align}
\boldsymbol U_5^{\rm fr}&=a_5\boldsymbol S_5^{\rm fr}+b_5\chi\boldsymbol m_5,
\\
\boldsymbol V_5^{\rm fr}&=a_5\boldsymbol S_5^{\rm fr}\times\boldsymbol S_{5,x}^{\rm fr}
+b_5\varrho\boldsymbol m_5+d_5\boldsymbol S_5^{\rm fr}+e_5\chi\boldsymbol m_5
\label{eq:ansatz5}
\end{align}
を曲率へ代入し、運動方程式に乗らない tensor structure を独立に消すと
\begin{equation}
d_5=-a_5^2,\qquad e_5=a_5b_5,\qquad 2a_5b_5=c^2
\label{eq:coeff5}
\end{equation}
を得る。$a_5=z_5\neq0$ と選べば
\begin{equation}
b_5=\frac{c^2}{2z_5},\qquad d_5=-z_5^2,\qquad e_5=\frac{c^2}{2}.
\end{equation}
従って
\begin{align}
\boldsymbol U_5^{\rm fr}(z_5)&=z_5\boldsymbol S_5^{\rm fr}
+\frac{c^2}{2z_5}\chi\boldsymbol m_5,
\label{eq:U5fr}\\
\boldsymbol V_5^{\rm fr}(z_5)&=z_5\boldsymbol S_5^{\rm fr}\times\boldsymbol S_{5,x}^{\rm fr}
+\frac{c^2}{2z_5}\varrho\boldsymbol m_5-z_5^2\boldsymbol S_5^{\rm fr}
+\frac{c^2}{2}\chi\boldsymbol m_5.
\label{eq:V5fr}
\end{align}

元の Class 5 変数で運動方程式残差を
\begin{equation}
\boldsymbol E_5:=\boldsymbol S_t-\boldsymbol S\times(\boldsymbol S_{xx}-\alpha_5M\boldsymbol S_x)
\label{eq:E5res}
\end{equation}
と定義する。行列曲率は
\begin{equation}
F_{tx}(z_5)=\iota\!\left(Q_{{\rm EOM},5}(z_5)\boldsymbol E_5\right),
\qquad
Q_{{\rm EOM},5}(z_5)=z_5\one3+\frac{c^2}{2z_5}K_5
\label{eq:QE5}
\end{equation}
と off shell に因数分解する。$K_5^2=0$ より
\begin{equation}
Q_{{\rm EOM},5}^{-1}=\frac1{z_5}\left(\one3-\frac{c^2}{2z_5^2}K_5\right),
\qquad
\det Q_{{\rm EOM},5}=z_5^3.
\label{eq:QEinv5}
\end{equation}
従って generic な $z_5\neq0$ で flatness と Class 5 EOM は同値である。

\begin{historymemo}
この Lax pair が見つかった時点では「Class 5 の missing Lax pair を発見した」という見方が有力だった。しかし後の source audit で、Jordanian/null-like Landau--Lifshitz hierarchy に既知 Lax structure が存在することを確認し、主張を Lax pair 自体から量子--古典 bridge へ移した。
\end{historymemo}

\section{Class 5 の解析式と ML 発見経路の対応}
Class 5 の ML 探索の feature library、loss、九回の独立 run、validation 結果は第\ref{sec:ml-class5}節にまとめた。そこで学習された
\begin{equation}
\beta_5=\frac{c^2}{2z_5}
\end{equation}
などの係数関係は、本 part で解析的に得た Lax pair と off-shell coefficient matrix の係数に一致する。従って ML は解析証明を置き換えるのではなく、式\eqref{eq:QE5}--\eqref{eq:QEinv5}へ至る局所 ansatz を絞り込んだ発見経路として位置付ける。

\section{量子 Jordanian Drinfeld twist と generic $L$-operator の source audit}
\status{Class 5 が XXX の Jordanian twist であることは文献入力。generic $L$ は twist definition から本ノートで再導出}
Class 5 量子鎖は XXX spin chain の Jordanian Drinfeld twist として記述される\cite{DFN2026}。量子 spectral parameter を $u$ とし、fundamental XXX $R$-matrix を
\begin{equation}
R_{\rm XXX}(u)=2u\one4+P
\label{eq:Rxxx}
\end{equation}
とする。twist parameter の一つを $a_2$ とし、$a_3=0$ の表示を使う。twist coefficient を
\begin{equation}
\vartheta:=\frac{a_2}{2}
\end{equation}
と定義する。physical spin generator を $\mathfrak s^a$ とし、$\mathfrak s^\pm=\mathfrak s^x\pm\ii\mathfrak s^y$ とする。Drinfeld twist $\mathcal F^{(5)}$ を定義式から作用させることで generic physical representation の $L$-operator を再導出した。

この source audit で重要だったのは、公開された generic $L$ 表式を本ノートの通常の operator product 規約で文字通り実装したものと、twist definition から直接得た表式が一致しなかったことである。direct-twist branch は spin-$\tfrac12$ で fundamental $R_5$ に戻り、$RLL$ relation を満たした。一方、公開表式の literal implementation は本ノートの規約では nonzero residual を残した。

\begin{cautionnote}
ここから「公開式は誤植である」と断定しない。未記載の operator ordering、基底、規格化の可能性を排除していない。安全な記録は「本ノートの明示した規約では direct twist が内部整合し、printed expression の literal implementation とは一致しなかった」である。
\end{cautionnote}

この監査は、ML の continuum branch が direct twist と整合する方を選んでいたことを後から確認する役割も持った。

\section{Class 5 の量子空間 Lax 作用素から continuum 空間 Lax 行列へ}
ここからは独立な spin-$\tfrac12$ 表示を使う。二空間の permutation 作用素 $P$ と nilpotent 作用素 $K$ を
\begin{equation}
P=\begin{pmatrix}1&0&0&0\\0&0&1&0\\0&1&0&0\\0&0&0&1\end{pmatrix},\qquad
K=\begin{pmatrix}0&1&-1&0\\0&0&0&0\\0&0&0&0\\0&0&0&0\end{pmatrix}
\label{eq:C5PK}
\end{equation}
と定義する。この表示では $P^2=\one4$、$K^2=0$、$PK=K$、$KP=-K$ である。量子変形パラメータを $a$ と書くと
\begin{equation}
R(u;a)=2u\one4+P+2auK,
\qquad
h=2P+2aK.
\label{eq:C5Rind}
\end{equation}

連続極限で\textbf{格子間隔 $\epsilon$}、\textbf{連続変形 coupling $\alpha$}、\textbf{連続スペクトルパラメータ $\lambda$}、\textbf{連続時間 $T$} を
\begin{equation}
a=\epsilon\alpha,\qquad u=\frac{\lambda}{\epsilon},\qquad x_n=n\epsilon,
\qquad T=\epsilon^2t_{\rm lat}
\label{eq:C5scaling}
\end{equation}
で定義する。この $\alpha,\lambda$ は本節の量子--連続照合用規約であり、前節の $\alpha_5,z_5$ とは定数変換を介して対応する。

単位行列に近づく規格化 Lax 作用素を
\begin{equation}
\widehat L:=\frac{L}{2u}=\one4+\epsilon\ell,
\qquad
\ell=\frac{P}{2\lambda}+\alpha K
\label{eq:C5Lhat}
\end{equation}
と定義する。ここでは厳密に $\ell^2=\one4/(4\lambda^2)$ が成り立つ。

spin の複素成分を
\begin{equation}
p:=S^1+\ii S^2,\qquad q:=S^1-\ii S^2,\qquad z:=S^3,
\qquad pq+z^2=1
\label{eq:C5pqz}
\end{equation}
と定義する。\textbf{$p$ は運動量ではなく spin の複素成分}である。lower symbol による空間 Lax 行列 $U:=\ell^\downarrow$ は
\begin{equation}
U=\frac1{4\lambda}
\begin{pmatrix}
1+z+2\alpha\lambda p&q-2\alpha\lambda(1+z)\\
p&1-z
\end{pmatrix}.
\label{eq:C5Uquant}
\end{equation}
場依存 trace を含む full $\mathfrak{gl}(2)$ connection であり、
\begin{equation}
\tr U=\frac1{2\lambda}+\frac\alpha2p,
\qquad \det U=\frac{\alpha p}{4\lambda}
\label{eq:C5trace}
\end{equation}
である。traceless part だけでなく scalar component も量子 $R$-matrix / classical $r$-matrix の照合では保持する。

\section{Class 5 の保存量 hierarchy と量子 transfer hierarchy}
Class 5 の continuum monodromy から Riccati expansion を行うと $Q_2,Q_3,Q_4$ を取り出せる。Riccati 用の定数基底変換後の spin を $\boldsymbol S_5^{\rm Ric}$ と書き、Riccati spectral parameter を $\zeta_5:=\ii/z_5$ と定義する。stereographic coordinates をこの小節だけ $w,v$ と呼ぶ。

展開係数は
\begin{equation}
\Gamma=w+\zeta_5\gamma_1+\zeta_5^2\gamma_2+\cdots,
\qquad
\gamma_1=-w_x,
\end{equation}
\begin{equation}
\gamma_2=w_{xx}-\frac{v w_x^2}{1+wv}-\frac{k_5w^3}{1+wv}
\end{equation}
と得られる。対応する局所密度 $j_r$ の積分を $I_{r,5}$ と書くと、境界項を除いて
\begin{equation}
I_{1,5}=-\frac12Q_{2,5},\qquad
I_{2,5}=\frac{\ii}{4}Q_{3,5},\qquad
I_{3,5}=\frac1{12}Q_{4,5}
\label{eq:charges5}
\end{equation}
が得られた。$Q_{4,5}$ の場の再定義後の局所密度は
\begin{align}
\mathcal Q_{4,5}^{\rm fr}={}&3\boldsymbol S_{xx}^{\mathsf T}\boldsymbol S_{xx}
-\frac{15}{4}(\boldsymbol S_x^{\mathsf T}\boldsymbol S_x)^2
+3k_5(\boldsymbol m_5^{\mathsf T}\boldsymbol S_x)^2\notag\\
&-\frac92k_5(\boldsymbol m_5^{\mathsf T}\boldsymbol S)^2(\boldsymbol S_x^{\mathsf T}\boldsymbol S_x)
-\frac34k_5^2(\boldsymbol m_5^{\mathsf T}\boldsymbol S)^4.
\label{eq:Q45}
\end{align}
ここで式を短くするため、この式だけ $\boldsymbol S=\boldsymbol S_5^{\rm fr}$ と書いた。

有限周期鎖の transfer matrix の logarithmic derivative から得る量子保存量との比較では、代表的な有限鎖で $Q_4^{\rm tr}-Q_4^{\rm local}$ の最大誤差が $3.59\times10^{-13}$、保存量間の commutator norm は $10^{-12}$ 程度以下だった。局所密度の差が全微分の場合、積分量の同一視には周期境界条件など境界項を消す条件が必要である。

\section{Class 5 と null-like warped $SL(2)$ Landau--Lifshitz hierarchy}
\status{Kameyama--Yoshida の Lax hierarchy は文献入力。以下は Class 5 変数との明示的規約辞書}
Kameyama--Yoshida の null-like warped $SL(2)$ Landau--Lifshitz model は Jordanian twist と関連する既知の Lax structure を持つ\cite{KY2014}。従って Class 5 の Lax pair 自体を新規とは主張しない。

Class 5 の場の再定義後の spin を $\boldsymbol S_5^{\rm fr}$ とし、比較先の orbit variable を $\boldsymbol n$ と呼ぶ。複素線形写像
\begin{equation}
\boldsymbol n=C_5\boldsymbol S_5^{\rm fr},
\qquad C_5=\operatorname{diag}(1,\ii,-\ii)
\label{eq:C5KYmap}
\end{equation}
を使う。KY 側の metric を $\eta=\operatorname{diag}(-1,1,1)$ とすると
\begin{equation}
C_5^{\mathsf T}\eta C_5=-\one3.
\end{equation}
従って compact complex sphere と noncompact real slice は同じ real theory ではないが、複素化した orbit 上では局所的な Hamiltonian / Poisson / Lax structure を比較できる。

KY 側の定数を null deformation $C_K$、正規化長さ $L_K$、coupling $\lambda_s$ と呼ぶ。本ノートとの coupling dictionary は
\begin{equation}
\alpha_5^2=\frac{32\pi^2L_K^2C_K}{\lambda_s}.
\label{eq:C5KYcoupling}
\end{equation}
Hamiltonian、Poisson tensor、symplectic form、時間は
\begin{equation}
H_K=-\frac{\lambda_s}{16\pi^2L_K}H_5^{\rm fr},
\qquad
\Pi_K=\frac2{L_K}\Pi_5^{\rm fr},
\qquad
\omega_K=\frac{L_K}{2}\omega_5^{\rm fr},
\label{eq:C5KYPoisson}
\end{equation}
\begin{equation}
t_K=\frac{8\pi^2L_K^2}{\lambda_s}t_5
\label{eq:C5KYtime}
\end{equation}
と対応する。適切な spectral parameter 再規格化と定数共役 $G_5=\operatorname{diag}(\ii,1)$ を入れると、空間・時間 Lax 行列も同時に対応する。

大域的には注意が必要である。元の Class 5 場が空間周期 $\ell_x$ で periodic でも、式\eqref{eq:fieldredef5}の再定義後は
\begin{equation}
\boldsymbol S_5^{\rm fr}(x+\ell_x)=e^{-\alpha_5\ell_xM/2}\boldsymbol S_5^{\rm fr}(x)
\label{eq:C5twistBC}
\end{equation}
となる。従って安全な結論は「局所的な複素化 Hamiltonian / Poisson / Lax hierarchy が対応する」であり、「同じ実模型」「同じ global boundary sector」とは言わない。

\begin{historymemo}
この照合で「Class 5 に新しい Lax pair」という見方は退いた。しかし研究は弱くなっただけではなかった。recent quantum Class 5 chain から既知 Jordanian hierarchy まで、量子 $R/L$、finite lattice、continuum、Poisson、保存量、boundary twist を一続きにする bridge という別の主題が見えた。
\end{historymemo}

\section{最初の coherent-symbol 補正はなぜ論理的に不十分だったか}
量子有限格子の時間 Lax 作用素を lower symbol へ移すと、既知の古典時間行列と単純には一致しなかった。初期段階では
\begin{equation}
\Delta V:=V_{\rm cont}-V^\downarrow
\label{eq:circularDelta}
\end{equation}
と差を定義し、その差が star-product correction を吸収することを確認した。この計算は恒等式の検証としては正しいが、\textbf{量子側から時間成分を導いた}という主張には循環がある。なぜなら $V_{\rm cont}$ を先に知っているからである。

この問題を解くため、後の独立監査では式\eqref{eq:discZC}に現れる $\mathcal A L$ と $L\mathcal A$ の共有格子点作用素積を、古典 $V$ を使わず直接 lower symbol へ縮約した。その結果として $\Delta V$ を後からではなく量子演算子積から構成し直した。

\begin{historymemo}
ここは研究の重要な自己修正である。「正しい答えとの差を補正と呼ぶ」だけでは mechanism の導出にならない。以後、補正の導出順序を $R\to h,\mathcal A\to(\mathcal AL)^\downarrow,(L\mathcal A)^\downarrow\to\mathcal C_{\rm site}\to\Delta V$ に固定した。
\end{historymemo}

\section{Class 5：共有格子点の作用素積から時間 Lax 補正を独立導出}
式\eqref{eq:C5Lhat}の一次作用素係数 $\ell$ の lower symbol が空間 Lax 行列 $U$ である。作用素二乗を取ってから lower symbol を取る場合と、lower symbol を取った行列を二乗する場合の差を\textbf{二次モーメントの symbol 差 $\Xi_5$} と定義する：
\begin{equation}
\Xi_5:=(\ell^2)^\downarrow-(\ell^\downarrow)^2
=\frac{\one2}{4\lambda^2}-U^2.
\label{eq:Xi5}
\end{equation}

有限格子 zero-curvature equation の共有 site を先に作用素として掛け、格子間隔二次の continuum coefficient を取り出すと、因子化処方で失われる項は
\begin{equation}
\mathcal C_{\rm site}^{(5)}=-2\ii\,\partial_x(U^2)=2\ii\,\partial_x\Xi_5.
\label{eq:Csite5}
\end{equation}
これは既知の古典時間 Lax 行列を入力せず得られる。

通常の zero-curvature equation に戻すための局所補正を
\begin{equation}
\Delta V_5=-2\ii U^2+\frac{\ii}{4\lambda^2}\one2
\label{eq:DeltaV5}
\end{equation}
と取ると $[\Delta V_5,U]=0$ であり、
\begin{equation}
\partial_x\Delta V_5+[\Delta V_5,U]=\mathcal C_{\rm site}^{(5)}.
\end{equation}
従って補正後の時間 Lax 行列
\begin{equation}
V_5=V_5^\downarrow+\Delta V_5
\label{eq:V5corr}
\end{equation}
が通常の continuum zero-curvature equation を満たす。

spin-$\tfrac12$ の一 site 二次モーメントは、spin coherent state の単位ベクトルを $\boldsymbol S$ として
\begin{equation}
\langle\widehat X_a\widehat X_b\rangle-\langle\widehat X_a\rangle\langle\widehat X_b\rangle
=\delta_{ab}-S^aS^b+\ii\varepsilon_{abc}S^c
\label{eq:spinhalf2nd}
\end{equation}
という形に整理できる。一般 spin-$j$ の affine operator class では同様の connected moment が $1/(2j)$ に比例するので、Class 5 の同じ affine class に限れば補正も $1/(2j)$ に比例する。任意の高次 operator ansatz へ一般化した定理ではない。

\section{Class 5：Hamiltonian EOM との独立な off-shell 照合と XXX control}
量子 Hamiltonian の coherent-state continuum limit から、独立に古典 Hamiltonian density
\begin{equation}
\mathcal H=-\frac12(p_xq_x+z_x^2)
+\frac\alpha2\left[(1+z)p_x-pz_x\right]
\label{eq:C5Hind}
\end{equation}
を得る。本節の Poisson bracket 規約では
\begin{equation}
\{S^a(x),S^b(y)\}=2\varepsilon_{abc}S^c(x)\delta(x-y).
\label{eq:C5PB2}
\end{equation}
運動方程式残差を $E_p,E_q,E_z$ と書くと、補正後の Lax 曲率は
\begin{equation}
F_{Tx}=\frac1{4\lambda}
\begin{pmatrix}
E_z+2\alpha\lambda E_p&E_q-2\alpha\lambda E_z\\
E_p&-E_z
\end{pmatrix}.
\label{eq:C5OffInd}
\end{equation}
従って
\begin{equation}
E_p=4\lambda F_{21},\qquad
E_z=-4\lambda F_{22},\qquad
E_q=4\lambda(F_{12}-2\alpha\lambda F_{22})
\label{eq:C5InverseInd}
\end{equation}
で曲率から EOM 残差を逆に復元できる。これは on-shell flatness より強い確認である。

変形を消す $\alpha=0$ の XXX / isotropic Landau--Lifshitz 極限でも、共有 site の作用素積補正は残る。従ってこの現象は Jordanian deformation だけの特殊事故ではない。

有限格子間隔 $\epsilon$ で直接作用素積を保持した近似は $\epsilon^2$ 程度で正しい continuum target へ収束した一方、factorized lower-symbol product は有限の差へ近づいた。ここから「原著者が誤った因子化をした」とは言わない。本ノートで明示的に定義した factorized procedure が失敗することだけを主張する。

\section{Class 5：古典 $r$-matrix / monodromy route との最終照合}
\status{Semenov--Tian--Shansky 型の生成公式は既知。本節の Class 5 への適用と量子 route との比較を独立検証}
ここでは二種類のスペクトルパラメータを区別する。最終的な Lax pair に残るものを\textbf{Lax スペクトルパラメータ $\mu$}、transfer / monodromy hierarchy を小さい値で展開して Hamiltonian flow を選ぶものを\textbf{生成スペクトルパラメータ $\zeta$} と呼ぶ。

spin の複素成分 $p,q,z$ は式\eqref{eq:C5pqz}と同じである。rank-one coherent-state projector を
\begin{equation}
\rho=\frac12\begin{pmatrix}1+z&q\\p&1-z\end{pmatrix},
\qquad \rho^2=\rho,\qquad\tr\rho=1
\label{eq:rhoSTS}
\end{equation}
と定義する。Class 5 空間 Lax 行列は
\begin{equation}
U(x,\mu)=\frac{\rho}{2\mu}+\alpha K,
\qquad
K:=\frac12\begin{pmatrix}p&-(1+z)\\0&0\end{pmatrix}.
\label{eq:USTS}
\end{equation}

上向き基底への projector を $P_\uparrow:=\operatorname{diag}(1,0)$、raising matrix を $E_+:=\begin{psmallmatrix}0&1\\0&0\end{psmallmatrix}$ と定義し、Jordanian tensor を
\begin{equation}
C_{12}:=P_\uparrow\otimes E_+-E_+\otimes P_\uparrow
\label{eq:C12}
\end{equation}
とする。本節の規約の古典 $r$-matrix は
\begin{equation}
r_{12}(\zeta,\mu)=\frac{\ii}{2(\zeta-\mu)}P_{12}+\ii\alpha C_{12}.
\label{eq:rSTS}
\end{equation}
これは前節までの別 gauge / normalization で書いた $r$-matrix と同じ式そのものではないので、無条件に成分比較しない。重要なのはこの節の $U$ と組にして線形 Poisson algebra を満たすことである。

monodromy を
\begin{equation}
T(x,y;\zeta)=\mathcal P\exp\!\left(\int_y^xU(s,\zeta)ds\right),
\qquad t(\zeta)=\tr T(L,0;\zeta)
\label{eq:monodromySTS}
\end{equation}
と定義する。標準的な $r$-matrix construction は
\begin{equation}
\mathbb V_b(x;\zeta,\mu)=t(\zeta)^{-1}\tr_a\!\left[
T_a(L,x;\zeta)r_{ab}(\zeta,\mu)T_a(x,0;\zeta)
\right]
\label{eq:STSformula}
\end{equation}
で時間側の generating matrix を与える\cite{DoikouKaraiskos2011}。

小さい生成スペクトルパラメータ $\zeta$ における局所 projector を
\begin{equation}
\Pi(x;\zeta)=\rho+\zeta\Pi_1+O(\zeta^2),
\qquad \Pi^2=\Pi,
\qquad \partial_x\Pi=[U,\Pi]
\label{eq:PiexpMain}
\end{equation}
と置くと、一次係数は
\begin{equation}
\Pi_1=2[\rho,\rho_x]-2\alpha[\rho,[K,\rho]].
\label{eq:Pi1Main}
\end{equation}
この一次係数が今回の Class 5 Hamiltonian flow に対応する。Poisson bracket の引数順による符号を合わせた physical time matrix を $V_r$、量子 finite-lattice route から得た時間行列を $V_q$ と呼ぶと、独立計算の結論は
\begin{equation}
V_r(x,\mu)-V_q(x,\mu)=\frac{\ii}{4\mu^2}\one2.
\label{eq:VrVq}
\end{equation}
右辺は場にも空間座標にも依存しない単位行列成分なので、空間微分も $U$ との交換子も零であり、両者は同じ zero-curvature equation を生成する。

この照合は局所的な $\zeta\to0$ の形式的漸近展開に基づく。$\exp(-c/\zeta)$ のような全てのべき級数で見えない monodromy 情報まで解析した、という主張ではない。

\begin{historymemo}
これが現在の Class 5 quantum--classical Lax bridge の計算上の終了点である。有限格子量子 time operator、独立 Hamiltonian/EOM、古典 $r$-matrix/monodromy の三経路が同じ $V$ へ閉じた。ここから別の模型を追加することを、現在の仕事を完成させる必須条件にはしない。
\end{historymemo}


\part{Class 6：nilpotent square-root Lax と eleven-vertex dictionary}

\section{Class 6 continuum 方程式と三次 nilpotent anisotropy}
Class 6 の文献基底の spin を $\boldsymbol S_{\rm paper}$ とし、反射行列
\begin{equation}
R_x:=\operatorname{diag}(-1,1,1)
\end{equation}
で計算基底の spin を $\boldsymbol S:=R_x\boldsymbol S_{\rm paper}$ と定義する。文献時間 $t_{\rm paper}$ に対し本節の連続時間を $T:=t_{\rm paper}/2$ とする。

計算基底の異方性行列は
\begin{equation}
N:=R_xN_{\rm paper}R_x
=\begin{pmatrix}0&0&-1\\0&0&-\ii\\-1&-\ii&0\end{pmatrix},
\qquad N^3=0.
\label{eq:N6}
\end{equation}
continuum coupling を $\alpha$ とし、運動方程式残差を
\begin{equation}
\boldsymbol E_6:=\boldsymbol S_T+2\boldsymbol S\times(\boldsymbol S_{xx}-\alpha N\boldsymbol S)
\label{eq:E6res}
\end{equation}
と定義する。

\section{Class 6 の有限行列平方根型 Lax pair と off-shell certificate}
Class 6 の square-root 表示の spectral parameter を $\zeta_6$ と呼ぶ。三次 nilpotency $N^3=0$ により、対称行列
\begin{equation}
A_{\zeta_6}=\zeta_6\one3+\frac{\alpha}{2\zeta_6}N
-\frac{\alpha^2}{8\zeta_6^3}N^2
\label{eq:A6}
\end{equation}
は厳密に
\begin{equation}
A_{\zeta_6}^2=\zeta_6^2\one3+\alpha N
\end{equation}
を満たす。時間成分用の行列は
\begin{equation}
B_{\zeta_6}=-\zeta_6^2\one3+\frac\alpha2N-
\frac{3\alpha^2}{8\zeta_6^2}N^2.
\label{eq:B6}
\end{equation}
ベクトル Lax pair を
\begin{equation}
\boldsymbol U_6=A_{\zeta_6}\boldsymbol S,
\qquad
\boldsymbol V_6=-2\left[B_{\zeta_6}\boldsymbol S+A_{\zeta_6}(\boldsymbol S\times\boldsymbol S_x)\right]
\label{eq:Lax6}
\end{equation}
と取ると、共通恒等式の時間規格化を合わせて
\begin{equation}
F_{Tx}=\iota(A_{\zeta_6}\boldsymbol E_6).
\label{eq:QE6}
\end{equation}
$\det A_{\zeta_6}=\zeta_6^3$ なので generic な $\zeta_6\neq0$ で flatness と EOM は同値である。

Class 6 の ML 探索の詳細は第\ref{sec:ml-class6}節にまとめた。有限 library $\{\one{3},N,N^2\}$ の係数として回収された関係は、式\eqref{eq:A6}の有限行列平方根と一致する。ML はこの square-root structure を認識する入口を与えたが、最終証明は式\eqref{eq:A6}--\eqref{eq:QE6}の解析恒等式である。

\section{Class 6 量子 $R$-matrix と continuum 空間 Lax 行列}
spin-$\tfrac12$ の量子表示で、二空間の permutation 作用素 $P$ と変形行列 $K$ を
\begin{equation}
K=\begin{pmatrix}
0&1&1&0\\0&0&0&-1\\0&0&0&-1\\0&0&0&0
\end{pmatrix}
\label{eq:C6K}
\end{equation}
とする。$P^2=\one4$、$PK=KP=K$、$K^3=0$ である。量子 deformation parameter を $a$、量子 spectral parameter を $u$ とすると
\begin{equation}
R(u;a)=2u\one4+P+au(1+2u)K+
\frac{a^2u^2(1+2u)^2}{2}K^2.
\label{eq:C6R}
\end{equation}
局所 Hamiltonian density は $h=2P+aK$ であり、代数的 inverse relation
\begin{equation}
R(u;a)R(-u;a)=(1-4u^2)\one4
\end{equation}
を満たす。

Class 6 では\textbf{格子間隔 $\epsilon$}に対して量子 deformation が二次で scaling する：
\begin{equation}
a=\epsilon^2\alpha,
\qquad u=\frac{\lambda}{\epsilon},
\qquad T=\epsilon^2t_{\rm lat}.
\label{eq:C6scaling}
\end{equation}
規格化 Lax 作用素 $\widehat L=L/(2u)$ は
\begin{equation}
\widehat L=\one4+\epsilon\ell_1+\epsilon^2\ell_2+\epsilon^3\ell_3
\label{eq:C6Lseries}
\end{equation}
という厳密な有限多項式になり、
\begin{align}
\ell_1&=\frac{P}{2\lambda}+\alpha\lambda K+\alpha^2\lambda^3K^2,\\
\ell_2&=\frac\alpha2K+\alpha^2\lambda^2K^2,\\
\ell_3&=\frac{\alpha^2\lambda}{4}K^2.
\label{eq:C6ells}
\end{align}
ここで重要な恒等式は
\begin{equation}
\ell_1^2=\frac{\one4}{4\lambda^2}+2\ell_2.
\label{eq:C6l1sq}
\end{equation}
Class 5 と違い、一次係数の二乗が二次係数を含む。

spin の複素成分を $S^\pm:=S^1\pm\ii S^2$ と定義すると、lower symbol から得る空間 Lax 行列は
\begin{equation}
U=\frac1{4\lambda}
\begin{pmatrix}
1+S^3+2\alpha\lambda^2S^+&
S^-+4\alpha\lambda^2S^3-4\alpha^2\lambda^4S^+\\
S^+&1-S^3-2\alpha\lambda^2S^+
\end{pmatrix}.
\label{eq:C6U}
\end{equation}
traceless part を $U_0:=U-\one2/(4\lambda)$ と書く。

\section{Class 6：共有格子点補正と Class 5 からの failure boundary}
Class 6 でも、共有 site の作用素積を古典時間成分を使わず直接評価する。二次モーメントの symbol 差を
\begin{equation}
\Xi_6:=(\ell_1^2)^\downarrow-(\ell_1^\downarrow)^2
=\frac{\one2}{4\lambda^2}+2\ell_2^\downarrow-U^2
\label{eq:Xi6}
\end{equation}
と定義する。空間 Lax 行列 $U$ に対する\textbf{adjoint covariant derivative} を
\begin{equation}
D_x^{(U)}X:=\partial_xX+[X,U]
\label{eq:Dadj}
\end{equation}
と定義すると、共有格子点補正は
\begin{equation}
\mathcal C_{\rm site}^{(6)}=2\ii D_x^{(U)}\Xi_6
\label{eq:Csite6}
\end{equation}
と書ける。局所時間 Lax 補正は
\begin{equation}
\Delta V_6=2\ii\Xi_6-\frac{\ii}{4\lambda^2}\one2
=4\ii\ell_2^\downarrow-2\ii U^2+\frac{\ii}{4\lambda^2}\one2
\label{eq:DeltaV6}
\end{equation}
と取れ、$\mathcal C_{\rm site}^{(6)}=\partial_x\Delta V_6+[\Delta V_6,U]$ である。

量子時間作用素の lower symbol だけから得る行列を $V_6^\downarrow$ と呼ぶと、補正後は
\begin{equation}
V_6=V_6^\downarrow+\Delta V_6
=-2\mathcal Q_\lambda(\boldsymbol S\times\boldsymbol S_x)-\ii\partial_\lambda U_0.
\label{eq:C6Vfinal}
\end{equation}
ここで古典ベクトル $\boldsymbol v$ を補助空間行列へ写す線形写像を
\begin{equation}
\mathcal Q_\lambda(\boldsymbol v):=\frac1{4\lambda}
\begin{pmatrix}
v^3+2\alpha\lambda^2v^+&v^-+4\alpha\lambda^2v^3-4\alpha^2\lambda^4v^+\\
v^+&-v^3-2\alpha\lambda^2v^+
\end{pmatrix}
\label{eq:Qmap6}
\end{equation}
と定義した。$v^\pm=v^1\pm\ii v^2$ である。

Class 6 が重要なのは、Class 5 の純粋な全微分型補正をそのまま一般化できないことを示す点である。例えば空間的に一様な瞬間配置 $\boldsymbol S=(0,0,1)^{\mathsf T}$、$\boldsymbol S_x=0$ でも、
\begin{equation}
\mathcal C_{\rm site}^{(6)}=[\Delta V_6,U]
=\begin{pmatrix}0&-\ii\alpha/\lambda\\0&0\end{pmatrix}
\label{eq:C6uniform}
\end{equation}
と nonzero である。従って correction の本質は「全微分」ではなく adjoint covariant derivative である。

\section{Class 6：Hamiltonian EOM と off-shell 因数分解の独立検証}
量子 Hamiltonian の continuum limit から、本節の規約の Hamiltonian density は
\begin{equation}
\mathcal H=-\frac12\boldsymbol S_x^{\mathsf T}\boldsymbol S_x+\alpha S^+S^3
=-\frac12\boldsymbol S_x^{\mathsf T}\boldsymbol S_x-\frac\alpha2\boldsymbol S^{\mathsf T}N\boldsymbol S
\label{eq:C6Hind}
\end{equation}
と得られる。Poisson bracket は
\begin{equation}
\{S^a(x),S^b(y)\}=2\varepsilon_{abc}S^c(x)\delta(x-y)
\end{equation}
とする。

運動方程式残差 $\boldsymbol E=(E^1,E^2,E^3)$ の複素成分を $E^\pm=E^1\pm\ii E^2$ とすると、補正後の曲率は
\begin{equation}
F_{Tx}=\mathcal Q_\lambda(\boldsymbol E)
\label{eq:C6Off}
\end{equation}
であり、逆に
\begin{align}
E^+&=4\lambda F_{21},\\
E^3&=4\lambda(F_{11}-2\alpha\lambda^2F_{21}),\\
E^-&=4\lambda(F_{12}-4\alpha\lambda^2F_{11}+12\alpha^2\lambda^4F_{21})
\label{eq:C6inv}
\end{align}
と復元できる。係数写像の determinant は $-1/(64\lambda^3)$ であり、generic な $\lambda\neq0$ で可逆である。

独立な symbolic audit では 40 個の等式が全て厳密に零となった。有限格子の discrete zero-curvature equation は 32 個の複素 parameter case で最大相対誤差 $5.02\times10^{-15}$、EOM を課さない局所 jet による factorization は 128 case で最大相対誤差 $1.18\times10^{-15}$ だった。有限格子間隔の比較では、leading continuum target と next-order target の収束次数が異なるため、同じ「誤差次数」として混同しない。

\section{Class 6 と eleven-vertex hierarchy の局所辞書}
\status{Class 6=eleven-vertex の量子同定と LOZ hierarchy は既知。以下は recent continuum variables との明示的規約辞書}
量子段階では Class 6 が eleven-vertex model であること自体は既知である\cite{CdL2024}。Class 6 $R$-matrix は適切な parameter map と定数 gauge transformation により LOZ の eleven-vertex $R$-matrix へ写る\cite{LOZ2014}。従って「Class 6 が実は eleven-vertex だった」という主張はしない。

continuum 側では、Class 6 の計算基底 spin $\boldsymbol S$ から $2\times2$ orbit matrix を
\begin{equation}
\Sigma:=\beta
\begin{pmatrix}
-S^3&S^1+\ii S^2\\S^1-\ii S^2&S^3
\end{pmatrix},
\qquad \beta^2=\frac\alpha8
\label{eq:Sigma6}
\end{equation}
と定義する。$\alpha\neq0$ では $\beta$ の平方根 branch を固定する。undeformed point $\alpha=0$ はこの明示写像へ直接代入せず、$\alpha\to0$ の連続極限として扱う。spin constraint より
\begin{equation}
\Sigma^2=\beta^2\one2.
\end{equation}

Class 6 bracket を orbit matrix へ写すと
\begin{equation}
\{\Sigma_{ij}(x),\Sigma_{kl}(y)\}_6
=\ii\beta(\Sigma_{kj}\delta_{il}-\Sigma_{il}\delta_{kj})\delta(x-y).
\label{eq:SigmaPB}
\end{equation}
従って Poisson tensor と symplectic form は定数倍で対応する。

LOZ flow の線形 map を、本ノートの generic anisotropy $\mathcal J$ と区別して
\begin{equation}
\mathcal J_{11}(\Sigma):=-
\begin{pmatrix}\Sigma_{12}&0\\\Sigma_{11}-\Sigma_{22}&-\Sigma_{12}\end{pmatrix}
\label{eq:J11}
\end{equation}
と書く。本ノートの zero-curvature / Poisson 規約に合わせた比較用 flow density を
\begin{equation}
\mathcal H_{11}^{\rm flow}:=\frac1{2\alpha}\tr(\Sigma_x^2)-\frac12\tr[\Sigma\mathcal J_{11}(\Sigma)]
\label{eq:H11flow}
\end{equation}
と定義すると、Class 6 density と $\mathcal H_6=4\mathcal H_{11}^{\rm flow}$ という成分恒等式が得られる。時間は
\begin{equation}
t_{11}=\frac{\ii\alpha}{\beta}T
\label{eq:t11}
\end{equation}
と再規格化される。spectral parameter も定数倍で写り、Poisson bracket、EOM、classical $r$-matrix、空間・時間 Lax 行列が同じ局所 dictionary の中で整合する。

\begin{cautionnote}
LOZ 論文が公開した Hamiltonian density そのものと Class 6 Hamiltonian を単純な定数倍として直接同一視しない。式\eqref{eq:H11flow}は flow を本ノートの規約で比較するための functional である。現在安全に主張するのは、複素化した局所 field space 上の EOM / Poisson / $r$-matrix / $U,V$ hierarchy の対応であり、global/unitary theory の同一性ではない。
\end{cautionnote}

\begin{historymemo}
Class 6 の eleven-vertex identification 自体は早い段階で既知だと分かったため、Class 6 単独の「新 Lax」は主結果候補から外れた。ただし量子 time-Lax continuum limit の failure boundary としては Class 5 にはない情報を与えた。
\end{historymemo}


\part{Class 5/6 に共通する量子--古典補正と文献上の位置づけ}

\section{二次モーメントの symbol 差による共通整理}
Class 5/6 の規格化空間 Lax 作用素を一般に
\begin{equation}
\widehat L=\mathbf 1+\epsilon\ell_1+\epsilon^2\ell_2+\cdots
\end{equation}
と書く。空間 Lax 行列を $U:=\ell_1^\downarrow$ と定義し、二次モーメントの symbol 差を
\begin{equation}
\Xi:=(\ell_1^2)^\downarrow-(\ell_1^\downarrow)^2
\label{eq:XiCommon}
\end{equation}
と定義する。$\Xi$ は確率論的な正値分散を意味せず、作用素積と lower-symbol 行列積の差をまとめた補助空間行列である。

adjoint covariant derivative $D_x^{(U)}$ は式\eqref{eq:Dadj}で定義した。この二模型の独立計算は共通して
\begin{equation}
\mathcal C_{\rm site}=2\ii D_x^{(U)}\Xi
\label{eq:commonC}
\end{equation}
と書ける。選んだ規格化では局所時間補正も
\begin{equation}
\Delta V=2\ii\Xi-\frac{\ii}{4\lambda^2}\one2
\label{eq:commonDelta}
\end{equation}
という共通形を持つ。

Class 5 では $\Xi_5=\one2/(4\lambda^2)-U_5^2$ で $[\Xi_5,U_5]=0$ となり、式\eqref{eq:commonC}は全微分へ簡単化する。Class 6 では $\Xi_6$ に $2\ell_2^\downarrow$ が残り、$[\Xi_6,U_6]$ が一般に nonzero である。従って Class 6 は「Class 5 の全微分公式を一般公式だと誤認しないための failure boundary」になっている。

\begin{cautionnote}
式\eqref{eq:commonC}を任意の量子可積分鎖に対する一般定理としては主張しない。現時点で示したのは Class 5 と Class 6 の指定表示・規格化・scaling に対する二つの worked examples である。
\end{cautionnote}

\section{既知の一般論と今回の具体的結果を切り分ける}
今回の研究で一度過大評価しかけた点を明確にする。

第一に、coherent-state lower symbol が作用素積を通常の積へ写さないこと自体は既知である\cite{SuzukiTsuchiya2017,Schlichenmaier2010}。従って「star-product correction を発見した」という主張はしない。

第二に、可積分量子 spin chain から classical continuum Lax structure を取る systematic long-wavelength limit は既知である\cite{AvanDoikouSfetsos2010}。従って「量子鎖の coherent continuum limit を初めて正しく定式化した」という主張もしない。

第三に、古典空間 Lax 行列と classical $r$-matrix / monodromy から時間 Lax 成分を生成する Semenov--Tian--Shansky 型の方法も既知である\cite{DoikouKaraiskos2011}。Class 5 の $V$ を classical $r$-matrix から構成できたこと自体を新一般論とはしない。

現在の安全な中心は、次の限定された chain である：
\begin{equation}
\begin{aligned}
&\text{finite-lattice quantum time Lax operator}
\longrightarrow \text{shared-site lower-symbol correction}\\
&\hspace{42mm}\longrightarrow \text{continuum }V
\end{aligned}
\label{eq:safeClaim}
\end{equation}
これを Class 5/6 で明示し、Class 5 では独立な classical $r$-matrix / monodromy route と閉じたことが計算上の核である。

\section{現在の新規性・進歩性評価}
この表は priority の不存在証明ではなく、論文化方針を固定する内部評価である。

\begin{longtable}{L{42mm}L{48mm}L{48mm}}
\toprule
項目 & 新規性についての現在の評価 & 物理・方法論上の価値\\
\midrule\endhead
Class 5 Lax pair 自体 & Jordanian/null-like LL hierarchy に既知構造があるため新規とはしない。 & 単独では主結果にしない。\\
Class 5 quantum--classical bridge & recent Class 5 quantum chain から finite-lattice time operator、continuum $V$、保存量、既知 Jordanian hierarchy までの明示的接続が主候補。 & 原論文が残した continuum time-component の穴に直接関係する。\\
Class 6 Lax pair 自体 & quantum Class 6=eleven-vertex と LL hierarchy は既知。 & 単独では弱い。\\
Class 6 continuum dictionary & recent continuum variables と eleven-vertex EOM/Poisson/$r/U,V$ の明示的 dictionary は整理価値がある。 & Class 5 の bridge を補強し、既知 hierarchy との関係を明確化。\\
共有格子点 correction の共通整理 & star-product 非乗法性は既知。二模型の time-Lax continuum limit で $2\ii D_x^{(U)}\Xi$ に整理した点が具体的候補。 & Class 6 が Class 5 の単純全微分の failure boundary を与える。\\
Class 5 STS closure & STS method 自体は既知。量子 finite-lattice route と独立 route が同じ $V$ に閉じたことが検証上重要。 & quantum--classical bridge の循環性を除き、終了判定を強める。\\
ML & 別模型への適用自体は主新規性ではない。 & off-shell certificate、source audit、analytic ansatz への入口として有効。\\
\bottomrule
\end{longtable}


\part{研究上の分岐、修正、保留した問題}

\section{soft particle production の構造的起源をどこまで追ったか}
Class 5 原論文の物理的に目立つ現象は、integrability と soft $1\to2$ particle production の共存である\cite{DFN2026}。一時期、この現象の構造的起源まで説明しないと研究として弱いのではないかと考えた。しかし、現在の quantum--classical Lax bridge の論理的完成条件とは別問題であると判断し、follow-up branch として保留した。

それでも途中計算には興味深い構造があるため記録する。stereographic coordinates を $w,\widetilde w$ とし、場の再定義前後を bar の有無で区別すると、Class 5 の $x$ 依存 transformation は
\begin{equation}
w=\overline w-cx,
\qquad
\widetilde w=\frac{\overline{\widetilde w}}{1+cx\,\overline{\widetilde w}}
\label{eq:Mobius}
\end{equation}
という parabolic M\"obius transformation になる。元の periodic field を長さ $\ell$ の空間円に置くと、再定義後の holonomy は
\begin{equation}
W_5=e^{-c\ell M}=\one3-c\ell M+\frac{(c\ell)^2}{2}M^2
\label{eq:W5}
\end{equation}
で、generic な $c\ell\neq0$ では eigenvalue 1 の unipotent Jordan block を持つ。

別 frame の quadratic linearization では、Fourier momentum を $k$ と書くと
\begin{equation}
K_{\rm lin}(k)=\begin{pmatrix}-k^2&\alpha^2/2\\0&k^2\end{pmatrix},
\qquad
K_{\rm lin}(0)^2=0
\label{eq:Ksoft}
\end{equation}
となり、対角化可能でなくなる locus が $k=0$ に現れる。原論文の soft emission support と同じ momentum locus に nilpotent structure が現れることは示唆的である。

\begin{cautionnote}
この branch は未完成である。場の再定義は vacuum と boundary condition を変えるため、式\eqref{eq:Ksoft}をそのまま元 frame の S-matrix Hamiltonian と同一視できない。nilpotent charge / Drinfeld dressing の Ward identity から soft amplitude の係数を導くところまでは到達していない。従って現在の論文候補の必須主結果には含めない。
\end{cautionnote}

\section{途中で修正した主張と、最終的な安全な言い方}
\begin{longtable}{L{50mm}L{88mm}}
\toprule
途中で考えた言い方 & 修正後の安全な記録\\
\midrule\endhead
「Class 5/6 に新しい Lax pair を発見した」 & 両者とも既知 hierarchy と対応する。新しい Lax hierarchy 自体は主張しない。\\
「KLS は flatness だけ見ている」 & KLS も EOM equivalence を componentwise loss へ入れている。違いは本研究が off-shell $F=QE$ と rank を explicit certificate にした点。\\
「公開 Class 5 generic $L$ は誤植」 & 本ノートの literal implementation では direct twist/RLL と不一致。未記載規約の可能性を排除しない。\\
「coherent-state symbol の非乗法性が新規 mechanism」 & 非乗法性は既知。今回重要なのは finite-lattice time-Lax equation の共有 site で、その connected moment が continuum $V$ に有限寄与する具体的構造。\\
「Class 5 correction は $-2\ii\partial_xU^2$ が一般式」 & Class 6 で破れる。共通形は二模型について $2\ii D_x^{(U)}\Xi$。一般 theorem は未証明。\\
「LOZ Hamiltonian と Class 6 Hamiltonian が同じ」 & 公開 LOZ Hamiltonian の直接同一視はしない。局所 EOM / Poisson / $r/U,V$ dictionary と比較用 flow functional を区別する。\\
「soft particle production の起源まで必要」 & 興味深い follow-up だが、quantum--classical Lax bridge の完成条件とは別。\\
\bottomrule
\end{longtable}

\section{検証の階層：何をどの独立性で確認したか}
本ノートでは、同じ結論を複数経路で確認する場合、どの入力を共有しているかを区別する。
\begin{enumerate}[leftmargin=2.5em]
\item \textbf{純粋代数：} nilpotency、行列平方根、$RLL$、Sutherland relation、有限格子 zero-curvature equation。
\item \textbf{運動学的 off-shell：} spin constraint は使うが EOM を使わず、$F=QE$ を確認する。
\item \textbf{Hamiltonian route：} 量子 Hamiltonian の coherent limit または continuum Hamiltonian の変分から EOM を独立に導き、Lax 曲率と比較する。
\item \textbf{quantum time-Lax route：} finite-lattice $\mathcal A_n$ を構成し、作用素積を先に評価して continuum $V$ を得る。
\item \textbf{classical $r$-matrix route：} Class 5 では monodromy generating function から $V_r$ を独立に作り、quantum route の $V_q$ と照合する。
\item \textbf{known-hierarchy audit：} Class 5 は KY、Class 6 は eleven-vertex / LOZ と field、Poisson、時間、spectral parameter、Lax structure を比較する。
\end{enumerate}

\section{現在の計算上の終了点と、次の論文で圧縮すべきもの}
現在の計算上のゴールは達成したと判断する。Class 5 について、
\begin{equation}
\begin{aligned}
\text{quantum }R/L
&\longrightarrow \text{finite-lattice time operator}
\longrightarrow \text{shared-site correction}\\
&\longrightarrow \text{continuum }U,V
\end{aligned}
\label{eq:endchain1}
\end{equation}
を閉じ、さらに
\begin{equation}
\text{classical }r\text{-matrix / monodromy}
\longrightarrow V_r
\simeq V_q
\label{eq:endchain2}
\end{equation}
を確認した。Class 6 についても quantum finite-lattice route から continuum $U,V$ を独立に回収し、Class 5 の簡単化が一般ではないことを示した。

次の論文草稿では、本ノートのすべてを残す必要はない。圧縮時の候補は次である。
\begin{enumerate}[leftmargin=2.5em]
\item 研究開始時の ML 探索は motivation / method history として短くする。
\item continuum Lax pair の full derivation は必要な certificate だけ本文へ置き、長い algebra は appendix へ移す。
\item Class 5/6 の量子 time-Lax continuum limit と共通 $\Xi$ 構造を主部に置く。
\item Class 5 の STS closure を独立 cross-check として残す。
\item KY / eleven-vertex dictionary は novelty overclaim を防ぐため重要だが、論文の主線を切らない長さへ圧縮する。
\item soft particle branch は現段階では本文の主結果に入れない。
\end{enumerate}

ここで論文を書く前に必要なのは、新しい物理計算を無制限に増やすことではなく、今回の finite-lattice quantum time-Lax route が既存文献のどの具体式まで既知で、どこからが今回の明示的 bridge なのかを citation-level で最終精査することである。


\part{独立監査記録：模型別の詳細導出}
\label{part:audits}
以下の三章は、上の整理された本文とは別に、最終段階で実行した独立再導出を詳細な計算順のまま保存する。本文と式が重複する箇所があるが、ここでは「既知の古典時間成分を入力しなかったか」「どの数値検査を独立に行ったか」を監査可能にするため意図的に残す。

\section{独立監査 A：Class 5 量子作用素積からの時間 Lax 成分}
\subsection{問題設定と今回の到達点}

Class 5 は、対角化不可能な transfer matrix を持ち得る可積分スピン鎖の一つである。de Leeuw--Fontanella--Nieto Garc\'ia はその連続極限を構成し、古典空間 Lax 成分を議論したが、対応する時間成分を得られなかったことを明記している。原論文は代わりに保存量と boost の構成を用いる\cite{DFN}。本報告では、その模型の局所的な量子--古典 Lax 対応に問題を限定する。粒子生成、散乱状態、境界 dressing の解析は行わない。

先行する研究用計算ノートには、Class 5/6 の時間成分と lower-symbol 補正の明示式が収録されている。本報告の課題は、それらの式を再掲することではない。補正を $V_{\rm cont}-\vdown$ として定義する前に、量子側の積から独立に求められるかを検証する。既存ノートのプログラムや保存済み数値は今回の検証コードに読み込んでいない。

以下で行う計算の順序は
\begin{equation}
R\ \longrightarrow\ h,\ \mathcal A
\ \longrightarrow\ (\mathcal A L)^{\downarrow},\ (L\mathcal A)^{\downarrow}
\ \longrightarrow\ \Cstar
\ \longrightarrow\ \Delta V,\ V
\label{audit5q:eq:order}
\end{equation}
である。古典運動方程式と照合するのは最後であり、$\Cstar$ の導出に既知の $V$ は用いない。

量子スピン鎖の古典連続極限そのもの\cite{ADS}と、coherent-state star product\cite{KS}は既存の方法である。本報告で新たに検証する内容は、指定した Class 5 の表示における補正の独立導出と、その簡潔な行列表示である。この形が文献上初めてであること、あるいは原論文の計算が失敗した原因を特定したことは主張しない。

\subsection{量子格子の規約}
\subsubsection{基本 $R$ 行列と Hamiltonian}

補助空間を $\mathcal V=\mathbb C^2$、各格子点の量子空間を $\mathcal H_n=\mathbb C^2$ とする。$R$ 行列は $\mathcal V\otimes\mathcal H_n$ 上に作用する。添字を省いた二空間上で
\begin{equation}
P=\begin{pmatrix}1&0&0&0\\0&0&1&0\\0&1&0&0\\0&0&0&1\end{pmatrix},
\qquad
K=\begin{pmatrix}0&1&-1&0\\0&0&0&0\\0&0&0&0\\0&0&0&0\end{pmatrix}
\label{audit5q:eq:PK}
\end{equation}
を定義する。$P$ は二つの空間を交換する作用素、$K$ は本報告の表示を指定する二空間上の nilpotent 作用素である。両者は
\begin{equation}
P^2=\Id_4,\qquad K^2=0,\qquad PK=K,\qquad KP=-K
\label{audit5q:eq:PKalgebra}
\end{equation}
を満たす。

原論文の式 (1.1) に $a_1=1,\ a_2=2a,\ a_3=0$ を代入すると\cite{DFN}
\begin{equation}
R(u;a)=2u\Id_4+P+2auK
=\begin{pmatrix}
1+2u&2au&-2au&0\\
0&2u&1&0\\
0&1&2u&0\\
0&0&0&1+2u
\end{pmatrix}.
\label{audit5q:eq:R}
\end{equation}
$u$ は量子スペクトルパラメータ、$a$ は量子変形パラメータである。正則性 $R(0)=P$ に基づき、格子 Hamiltonian を
\begin{equation}
h_{n,n+1}(a)=P_{n,n+1}\partial_uR_{n,n+1}(u;a)\big|_{u=0}
=2P_{n,n+1}+2aK_{n,n+1},\qquad
H_{\rm lat}=\sum_n h_{n,n+1}
\label{audit5q:eq:h}
\end{equation}
と定める。周期鎖では格子添字を周期的に読む。以下の局所恒等式は隣接する二つの bond だけに依存する。

\subsubsection{有限格子の時間成分}

$L_n(u)=R_{0n}(u;a)$ と書き、$0$ を補助空間の添字に用いる。格子時間 $t_{\rm lat}$ に対する Heisenberg 方程式は
\begin{equation}
\partial_{t_{\rm lat}}L_n=\ii[H_{\rm lat},L_n]
\label{audit5q:eq:Heisenberg}
\end{equation}
とする。必要な Sutherland relation は
\begin{equation}
[h_{n-1,n},L_nL_{n-1}]
=L_n\partial_uL_{n-1}-(\partial_uL_n)L_{n-1}.
\label{audit5q:eq:Sutherland}
\end{equation}
これは式\eqref{audit5q:eq:R}を代入して作用素恒等式として確認できる。これにより
\begin{equation}
\mathcal A_n=\ii h_{n-1,n}
-\ii L_n^{-1}h_{n-1,n}L_n
-\ii L_n^{-1}\partial_uL_n
\label{audit5q:eq:rawA}
\end{equation}
を取れば
\begin{equation}
\ii[H_{\rm lat},L_n]=\mathcal A_{n+1}L_n-L_n\mathcal A_n
\label{audit5q:eq:discreteZC}
\end{equation}
が成立する。逆行列は正則なスペクトル領域で用い、得られる式は有理関数として継続する。

\subsubsection{格子間隔・時間・スカラー成分の規格化}

格子間隔 $\eps$ に対し
\begin{equation}
a=\eps\alpha,\qquad u=\frac{\lambda}{\eps},\qquad
x_n=n\eps,\qquad T=\eps^2t_{\rm lat},\qquad\lambda\ne0
\label{audit5q:eq:scaling}
\end{equation}
を用いる。$\alpha$ は連続理論の変形パラメータ、$\lambda$ は連続スペクトルパラメータである。規格化した空間演算子は厳密に
\begin{equation}
\widehat L_n=\frac{L_n}{2u}=\Id+\eps\ell_{0n},\qquad
\ell=\frac{P}{2\lambda}+\alpha K,\qquad
\ell^2=\frac{\Id_4}{4\lambda^2}
\label{audit5q:eq:Lhat}
\end{equation}
となる。最後の恒等式は式\eqref{audit5q:eq:PKalgebra}から従う。

時間成分の共通スカラーを
\begin{equation}
\widehat{\mathcal A}_n=\mathcal A_n+\frac{\ii}{u}\Id
=-\ii\widehat L_n^{-1}\bigl([h_{n-1,n},\widehat L_n]
+\partial_u\widehat L_n\bigr)
\label{audit5q:eq:Ahat}
\end{equation}
によって除く。この変更は式\eqref{audit5q:eq:discreteZC}を変えない。$h$ は式\eqref{audit5q:eq:h}で、規格化前の正則な $R$ から定義したままである。微分 $\partial_u$ は量子変形パラメータ $a$ を固定して行い、その後に式\eqref{audit5q:eq:scaling}を代入するので
\begin{equation}
\partial_u\widehat L_n=-\frac{P_{0n}}{2u^2}
=-\frac{\eps^2P_{0n}}{2\lambda^2}
\label{audit5q:eq:udiff}
\end{equation}
である。

\subsection{coherent-state lower symbol と共有格子点}
\subsubsection{場と行列の定義}

Pauli 行列を $\sigma_a$ とし、古典 spin 場 $\bs(T,x)$ は
\begin{equation}
\bs^{\mathsf T}\bs=1
\label{audit5q:eq:spinconstraint}
\end{equation}
を満たすとする。実球面上で spin-$\tfrac12$ coherent state の密度行列は
\begin{equation}
\rho(\bs)=\frac{\Id_2+\bs\cdot\bm\sigma}{2}.
\end{equation}
非 Hermitian 模型の局所代数を扱う際にはこれを複素化する。その場合も $\tr\rho=1,\ \rho^2=\rho$ は成り立つが、正値性は主張しない。

式を短くするため
\begin{equation}
p=S^1+\ii S^2,\qquad q=S^1-\ii S^2,\qquad z=S^3,
\qquad pq+z^2=1
\label{audit5q:eq:pqz}
\end{equation}
と書く。$p,q$ は spin の成分であって運動量ではなく、$z$ もスペクトルパラメータではない。複素化後は $p,q$ を独立に扱う。したがって
\begin{equation}
\rho=\frac12\begin{pmatrix}1+z&q\\p&1-z\end{pmatrix}.
\end{equation}

演算子 $B$ の lower symbol $B^{\downarrow}$ は物理空間だけの期待値である。例えば二格子点に依存する場合
\begin{equation}
B^{\downarrow}(\bs_1,\bs_2)
=\tr_{1,2}\bigl[(\Id_{\mathcal V}\otimes\rho(\bs_1)\otimes\rho(\bs_2))B\bigr]
\in\End(\mathcal V).
\label{audit5q:eq:lower}
\end{equation}
補助空間上の行列積の順序は維持する。空間成分は
\begin{equation}
U=\ell^{\downarrow}
=\frac1{4\lambda}
\begin{pmatrix}
1+z+2\alpha\lambda p&q-2\alpha\lambda(1+z)\\
p&1-z
\end{pmatrix}.
\label{audit5q:eq:U}
\end{equation}
特に
\begin{equation}
\tr U=\frac1{2\lambda}+\frac\alpha2p,\qquad
\det U=\frac{\alpha p}{4\lambda}
\label{audit5q:eq:traceU}
\end{equation}
であり、trace は一般に場に依存する。以下の簡潔な補正式ではこの $\mathfrak{gl}(2)$ 表示を用いる。trace を捨てた $U$ に同じ補正式を無条件で適用してはいけない。

後の線形写像を
\begin{equation}
\Qmap(\bm v)=\frac1{4\lambda}
\begin{pmatrix}
v^3+2\alpha\lambda(v^1+\ii v^2)&v^1-\ii v^2-2\alpha\lambda v^3\\
v^1+\ii v^2&-v^3
\end{pmatrix}
\label{audit5q:eq:Qmap}
\end{equation}
と定義する。これは $U$ の spin 依存部分の線形係数であり、$U_T=\Qmap(\bs_T)$、$U_x=\Qmap(\bs_x)$ である。

\subsubsection{通常積と star product の差}

一つの物理格子点に対して
\begin{equation}
A=A_0\otimes\Id_2+\sum_{a=1}^3A_a\otimes\sigma_a,
\qquad
B=B_0\otimes\Id_2+\sum_{b=1}^3B_b\otimes\sigma_b
\end{equation}
と展開する。$A_a,B_b$ は補助空間上の行列で、互いに可換とは仮定しない。Pauli 行列の積から厳密に
\begin{equation}
(AB)^{\downarrow}-A^{\downarrow}B^{\downarrow}
=\sum_{a,b}A_aB_b
\bigl(\delta_{ab}-S^aS^b+\ii\epsilon_{abc}S^c\bigr)
\label{audit5q:eq:star}
\end{equation}
を得る。ここで $\epsilon_{123}=1$ は Levi-Civita 記号であり、格子間隔 $\eps$ と区別する。一般に $(AB)^{\downarrow}=A^{\downarrow}\star B^{\downarrow}$ と書く star product の、spin-$\tfrac12$ における明示式である\cite{KS}。

$\widehat{\mathcal A}_{n+1}$ と $\widehat L_n$ は物理格子点 $n$ を共有する。そのため全状態を product coherent state にしても、式\eqref{audit5q:eq:star}の右辺は消えない。異なる格子点間の因子化と、同一格子点上の作用素積の因子化は別である。

\subsection{量子作用素積から補正を求める}
\subsubsection{必要な展開次数}

この節では三つの tensor factor を $(0,1,2)=($補助空間、左格子点、右格子点$)$ として固定する。略記
\begin{equation}
\ell_L=\ell_{01},\qquad\ell_R=\ell_{02},\qquad
h_0=2P_{12},\qquad h_1=2\alpha K_{12}
\end{equation}
を用いる。式\eqref{audit5q:eq:Ahat}の作用素展開は
\begin{align}
\widehat{\mathcal A}&=\eps A_1+\eps^2 A_2+O(\eps^3),\label{audit5q:eq:Aexpand}\\
A_1&=-\ii[h_0,\ell_R],\label{audit5q:eq:A1}\\
A_2&=\ii\ell_R[h_0,\ell_R]-\ii[h_1,\ell_R]
+\frac{\ii P_{02}}{2\lambda^2}.\label{audit5q:eq:A2}
\end{align}
$A_1,A_2$ は補助空間を含む $8\times8$ 行列である。

$\bm s_1,\bm s_2$ を二つの古典 spin とし
\begin{equation}
f_r(\bm s_1,\bm s_2)=A_r^{\downarrow}(\bm s_1,\bm s_2),\qquad r=1,2
\end{equation}
と置く。$\partial_1 f[\bm v]$ は第1引数の方向 $\bm v$ への微分、$\partial_2 f[\bm v]$ は第2引数についての同じ微分である。直接の Pauli contraction は
\begin{equation}
f_1(\bm s_1,\bm s_2)=-2\Qmap(\bm s_1\times\bm s_2),\qquad
f_1(\bs,\bs)=0
\label{audit5q:eq:f1}
\end{equation}
を与える。この零によって、作用素としては $O(\eps)$ の $\widehat{\mathcal A}$ の lower symbol は、滑らかな場上で $O(\eps^2)$ から始まる。

$x_n=x$ で $\widehat{\mathcal A}_n$ の支持は $(x-\eps,x)$ なので
\begin{equation}
\vdown=\lim_{\eps\to0}\frac{\widehat{\mathcal A}_n^{\downarrow}}{\eps^2}
=f_2(\bs,\bs)-\partial_1f_1(\bs,\bs)[\bs_x].
\label{audit5q:eq:Vdown}
\end{equation}
一方、離散零曲率式の左辺は $\partial_{t_{\rm lat}}\widehat L_n=\eps^3U_T$ に対応する。従って右辺の作用素積は $O(\eps^3)$ まで保つ必要がある。

\subsubsection{共有格子点の補正を分離する}

二格子点の関数として
\begin{align}
C_{R,2}&=(A_1\ell_L)^{\downarrow}-f_1\ell_L^{\downarrow},&
C_{L,2}&=(\ell_RA_1)^{\downarrow}-\ell_R^{\downarrow}f_1,\label{audit5q:eq:C2def}\\
C_{R,3}&=(A_2\ell_L)^{\downarrow}-f_2\ell_L^{\downarrow},&
C_{L,3}&=(\ell_RA_2)^{\downarrow}-\ell_R^{\downarrow}f_2
\label{audit5q:eq:C3def}
\end{align}
を定義する。添字 $R,L$ は離散式の右側・左側の時間成分に対応し、2,3 は $\eps$ の次数である。これら四行列は式\eqref{audit5q:eq:A1}、\eqref{audit5q:eq:A2}、\eqref{audit5q:eq:lower}だけから決まる。

連結部分は、右側の時間成分では $(\bs(x),\bs(x+\eps))$、左側では $(\bs(x-\eps),\bs(x))$ で評価される。$O(\eps^2)$ の差は
\begin{equation}
[C_{R,2}-C_{L,2}]_{(\bs,\bs)}=0
\end{equation}
であり、最初の有限の寄与は
\begin{equation}
\Cstar=\bigl[
\partial_2 C_{R,2}[\bs_x]+\partial_1 C_{L,2}[\bs_x]
+C_{R,3}-C_{L,3}
\bigr]_{(\bs,\bs)}.
\label{audit5q:eq:Cfromquantum}
\end{equation}
$A_3$ の通常積部分は同一点で左右の差が相殺され、$A_3$ と $\ell$ の連結部分は $O(\eps^4)$ なので、この次数の計算には必要ない。

式\eqref{audit5q:eq:Cfromquantum}を共有格子点の式\eqref{audit5q:eq:star}で評価すると
\begin{equation}
\Cstar=-2\ii\partial_x(U^2).
\label{audit5q:eq:Cresult}
\end{equation}
この段階では古典時間成分も運動方程式も入力していない。成分ごとの結果は付録\ref{audit5q:app:C}に掲げる。評価には $\bs^2=1$ とその $x$ 微分だけを用いる。

式\eqref{audit5q:eq:Lhat}に着目すると、同じ結果を一格子点の二次 moment の差で
\begin{align}
\Xi&=(\ell^2)^{\downarrow}-(\ell^{\downarrow})^2
=\frac{\Id_2}{4\lambda^2}-U^2,\label{audit5q:eq:Xi}\\
\Cstar&=2\ii\partial_x\Xi
\end{align}
と書ける。この簡約には Class 5 の式\eqref{audit5q:eq:PKalgebra}を用いており、任意の $R$ 行列に同じ式が成立するとは限らない。

\subsection{局所補正と時間 Lax 成分}
\subsubsection{通常の零曲率方程式への復元}

量子側から得た式は
\begin{equation}
U_T=\partial_x\vdown+[\vdown,U]+\Cstar
\label{audit5q:eq:projected}
\end{equation}
である。式\eqref{audit5q:eq:Cresult}に対して
\begin{equation}
\Delta V=-2\ii U^2+\frac{\ii}{4\lambda^2}\Id_2
\label{audit5q:eq:DeltaV}
\end{equation}
を取る。$\Delta V$ は $U$ の多項式なので
\begin{equation}
[\Delta V,U]=0,\qquad
\partial_x\Delta V+[\Delta V,U]=\Cstar
\label{audit5q:eq:absorption}
\end{equation}
が成り立つ。従って
\begin{equation}
V=\vdown+\Delta V,\qquad
F_{Tx}=U_T-\partial_xV+[U,V]=0
\label{audit5q:eq:ZC}
\end{equation}
を得る。式\eqref{audit5q:eq:DeltaV}のスカラー項は規約の選択であり、任意の場・$x$ に依存しない $f(\lambda)\Id_2$ をさらに足しても曲率は変わらない。ここで行っているのは時間成分の補正であり、それ自体を gauge transformation と呼ばない。

計算結果は次の短い形に整理できる：
\begin{align}
W&=-2\Qmap(\bs\times\bs_x),\label{audit5q:eq:Wdef}\\
Z&=2\ii U^2-\frac{\ii}{4\lambda^2}\Id_2,\label{audit5q:eq:Zdef}\\
\vdown&=W+2Z,\qquad
\Delta V=-Z,\qquad
V=W+Z.
\label{audit5q:eq:WZ}
\end{align}
従って
\begin{equation}
V=-2\Qmap(\bs\times\bs_x)+2\ii U^2-rac{\ii}{4\lambda^2}\Id_2.
\label{audit5q:eq:Vcompact}
\end{equation}
$W$ の成分を明示すると
\begin{align}
W_{11}&=-\frac{\ii}{4\lambda}
\bigl[pq_x-p_xq+4\alpha\lambda(zp_x-pz_x)\bigr],\label{audit5q:eq:W11}\\
W_{12}&=\frac{\ii}{2\lambda}
\bigl[zq_x-qz_x+\alpha\lambda(pq_x-p_xq)\bigr],\\
W_{21}&=\frac{\ii}{2\lambda}(pz_x-zp_x),\qquad
W_{22}=\frac{\ii}{4\lambda}(pq_x-p_xq).
\end{align}
式\eqref{audit5q:eq:U}、\eqref{audit5q:eq:WZ}、\eqref{audit5q:eq:Vcompact}だけで $U,V$ の全成分を復元できる。

\subsubsection{局所補正の限定的な一意性}

$\Delta V$ を、明示的な $x$ 依存も場の微分も持たない $\Delta V(\bs;\lambda)$ のクラスで探す。二つの解の差を $B$ とすると
\begin{equation}
\partial_x B+[B,U]=0
\label{audit5q:eq:homogeneous}
\end{equation}
が任意の滑らかな単位 spin 場に対して成り立たなければならない。同一点で接ベクトル $\bs_x$ を自由に変えると、$B$ は球面上で一定でなければならない。残る条件は、全ての $\bs$ に対して $[B,U(\bs)]=0$ である。

$\lambda\ne0$ では、$\Qmap(\bm v)$ の trace-free 部分は、$\bm v$ を変えると三次元の trace-free 行列空間全体を張る。従ってそれら全てと可換な定数 $2\times2$ 行列はスカラーだけである。このクラスでは式\eqref{audit5q:eq:DeltaV}は場に依存しないスカラーを除いて一意となる。より高い場の微分を許したクラスや別の gauge 表示について、一意性を主張してはいない。

\subsection{運動方程式による独立照合}
\subsubsection{Hamiltonian の連続極限}

$\widehat{\mathcal A}$ の計算とは独立に、式\eqref{audit5q:eq:h}の期待値を展開する。同じ規約で
\begin{equation}
h^{\downarrow}\bigl(\bs(x),\bs(x+\eps)\bigr)
=2+\eps^2\mathcal H+O(\eps^3),
\end{equation}
\begin{equation}
\mathcal H=-\frac12(p_xq_x+z_x^2)
+\frac\alpha2\bigl[(1+z)p_x-pz_x\bigr],\qquad
H_{\rm cl}=\int dx\,\mathcal H.
\label{audit5q:eq:Hcl}
\end{equation}
有限長を固定した滑らかな場では $(H_{\rm lat}-2N)/\eps$ の symbol が $H_{\rm cl}$ に対応する。$N$ は格子点数である。$\alpha p_x/2$ は全微分だが、表示の由来を明らかにするため式には残した。積分 Hamiltonian の同一視には、周期境界条件または対応する境界項の処理を仮定する。

spin-$\tfrac12$ の $[\sigma_a,\sigma_b]=2\ii\epsilon_{abc}\sigma_c$ と式\eqref{audit5q:eq:scaling}に対応する連続 Poisson bracket を
\begin{equation}
\{S^a(x),S^b(y)\}_{\rm cl}
=2\epsilon_{abc}S^c(x)\delta(x-y),\qquad
\bs_T=\{\bs,H_{\rm cl}\}_{\rm cl}
\label{audit5q:eq:Poisson}
\end{equation}
と取る。この係数2を落とすと、格子時間から導いた時間成分とは規格化が一致しない。複素成分では
\begin{equation}
\{p,q\}_{\rm cl}=-4\ii z\delta,\quad
\{p,z\}_{\rm cl}=2\ii p\delta,\quad
\{q,z\}_{\rm cl}=-2\ii q\delta
\end{equation}
である。式\eqref{audit5q:eq:Hcl}から
\begin{align}
p_T&=2\ii(pz_{xx}-zp_{xx}+\alpha pp_x),\label{audit5q:eq:Epflow}\\
q_T&=2\ii(zq_{xx}-qz_{xx}+\alpha pq_x),\label{audit5q:eq:Eqflow}\\
z_T&=\ii(qp_{xx}-pq_{xx}+2\alpha pz_x)\label{audit5q:eq:Ezflow}
\end{align}
を得る。これらは $\partial_T(pq+z^2)=0$ を満たす。

さらに、量子離散零曲率式の右辺を作用素積のまま $O(\eps^3)$ まで展開しても、式\eqref{audit5q:eq:Epflow}--\eqref{audit5q:eq:Ezflow}による $U_T$ が得られる。したがって Hamiltonian、量子作用素積、および補正後の古典曲率の三者を同じ規約で照合できる。

ここで議論しているのは coherent-state symbol 上の古典的なベクトル場である。厳密な量子時間発展が product coherent state の部分多様体を保つことや、任意の量子状態の相関関数がこの極限へ収束することは証明していない。

\subsubsection{off-shell 因数分解と逆向きの含意}

式\eqref{audit5q:eq:Epflow}--\eqref{audit5q:eq:Ezflow}の左辺から右辺を引いた残差を $E_p,E_q,E_z$ と定義する。$\bs^2=1$ と空間微分の拘束のみを使い、時間微分に運動方程式を課さず曲率を計算すると
\begin{equation}
F_{Tx}=\frac1{4\lambda}
\begin{pmatrix}
E_z+2\alpha\lambda E_p&E_q-2\alpha\lambda E_z\\
E_p&-E_z
\end{pmatrix}.
\label{audit5q:eq:offshell}
\end{equation}
$\lambda\ne0$ では逆に
\begin{equation}
E_p=4\lambda F_{21},\qquad E_z=-4\lambda F_{22},\qquad
E_q=4\lambda(F_{12}-2\alpha\lambda F_{22})
\end{equation}
と復元できる。従って $F_{Tx}=0$ と古典運動方程式は同値である。単に運動方程式の上で消える行列を作っただけではない。

\subsubsection{原論文の場・時間との局所的な対応}

原論文の式 (2.22) は、場 $\bm S_{\rm ref}$、変形パラメータ $\alpha_{\rm ref}$、時間 $t_{\rm ref}$ により\cite{DFN}
\begin{equation}
\partial_{t_{\rm ref}}\bm S_{\rm ref}
=\bm S_{\rm ref}\times\bigl(\bm S_{{\rm ref},xx}
-\alpha_{\rm ref}M_{\rm ref}\bm S_{{\rm ref},x}\bigr),\qquad
M_{\rm ref}=\begin{pmatrix}0&0&-1\\0&0&\ii\\1&-\ii&0\end{pmatrix}
\label{audit5q:eq:refEOM}
\end{equation}
と書ける。本報告の式との辞書は
\begin{equation}
\bm S_{\rm ref}=\operatorname{diag}(1,-1,1)\bs,\qquad
\alpha_{\rm ref}=-\alpha,\qquad t_{\rm ref}=2T.
\label{audit5q:eq:dictionary}
\end{equation}
反射により外積の符号が変わることを含めると、式\eqref{audit5q:eq:Epflow}--\eqref{audit5q:eq:Ezflow}は式\eqref{audit5q:eq:refEOM}に一致する。この辞書は局所方程式と時間規格化の照合であり、量子 Hilbert 空間上の unitary equivalence を主張するものではない。

\subsection{XXX 極限と有限格子間隔での検証}
\subsubsection{変形を消しても補正は必要である}

$\alpha=0$ とし $\mathsf S=\bs\cdot\bm\sigma$ と書くと
\begin{align}
U&=\frac{\Id_2+\mathsf S}{4\lambda},\qquad
W=-\frac{(\bs\times\bs_x)\cdot\bm\sigma}{2\lambda},\\
\vdown&=W+\frac{\ii\mathsf S}{2\lambda^2},\qquad
\Delta V=-\frac{\ii\mathsf S}{4\lambda^2},\qquad
V=W+\frac{\ii\mathsf S}{4\lambda^2},\label{audit5q:eq:XXX}\\
\Cstar&=-\frac{\ii\,\bs_x\cdot\bm\sigma}{4\lambda^2}.
\end{align}
正しい運動方程式は $\bs_T=-2\bs\times\bs_{xx}$ である。式\eqref{audit5q:eq:XXX}は、今回の差が Jordanian 変形だけに由来する問題ではなく、固定 spin-$\tfrac12$ での symbol 積の扱いに由来することを示す。

\subsubsection{数値検証の定義}

全ての以下の数値は本報告のコードを実行して得たものであり、既存ノートからの転記ではない。まず有限格子の局所恒等式\eqref{audit5q:eq:discreteZC}を、補助空間と三つの物理格子点からなる $16$ 次元空間で検査した。ランダムな複素 $a,u$ の24例で、相対 Frobenius residual の最大値は
\begin{equation}
4.74\times10^{-16}.
\end{equation}
相対 residual は $\|L-R\|_{\rm F}/\max(1,\|L\|_{\rm F},\|R\|_{\rm F})$ と定義する。また、非可換な補助行列係数を含む式\eqref{audit5q:eq:star}を32例で直接 trace と比較し、絶対 Frobenius residual の最大値は $6.78\times10^{-15}$ であった。

格子間隔の検査には
\begin{align}
\theta(x)&=0.9+0.2\sin x,\qquad
\phi(x)=0.3+0.4x+0.1\cos(2x),\\
\bs(x)&=(\sin\theta\cos\phi,\ \sin\theta\sin\phi,\ \cos\theta),\qquad
x=0.37,\quad\lambda=0.8+0.2\ii
\end{align}
を用いた。場の $x$ 微分はこの式から解析的に求める。$\alpha=0$ と $\alpha=0.7-0.2\ii$ を比較する。

Taylor 展開を使わず有限 $\eps$ の行列逆と積を直接計算し
\begin{align}
J_{\rm exact}^{(\eps)}&=\eps^{-3}
\bigl[(\widehat{\mathcal A}_{n+1}\widehat L_n)^{\downarrow}
-(\widehat L_n\widehat{\mathcal A}_n)^{\downarrow}\bigr],\\
J_{\rm fact}^{(\eps)}&=\eps^{-3}
\bigl[\widehat{\mathcal A}_{n+1}^{\downarrow}\widehat L_n^{\downarrow}
-\widehat L_n^{\downarrow}\widehat{\mathcal A}_n^{\downarrow}\bigr],\\
C^{(\eps)}&=J_{\rm exact}^{(\eps)}-J_{\rm fact}^{(\eps)}
\end{align}
を得る。表\ref{audit5q:tab:eps}では
\begin{equation}
e_{\rm exact}=\|J_{\rm exact}^{(\eps)}-U_T\|_{\rm F},\qquad
e_C=\|C^{(\eps)}-\Cstar\|_{\rm F},\qquad
e_{\rm fact}=\|J_{\rm fact}^{(\eps)}-U_T\|_{\rm F}
\end{equation}
を示す。$U_T$ は式\eqref{audit5q:eq:Epflow}--\eqref{audit5q:eq:Ezflow}から評価した。

\begin{table}[htbp]
\centering
\caption{有限格子間隔での直接行列計算。全て絶対 Frobenius norm。}
\label{audit5q:tab:eps}
\small
\begin{tabular}{llrrr}
\toprule
模型 & $\eps$ & $e_{\rm exact}$ & $e_C$ & $e_{\rm fact}$\\
\midrule
XXX & $0.08$ &$6.213\times10^{-4}$ &$8.835\times10^{-3}$ &$0.15721$\\
    & $0.02$ &$3.888\times10^{-5}$ &$2.116\times10^{-3}$ &$0.15137$\\
    & $0.005$ &$2.430\times10^{-6}$ &$5.235\times10^{-4}$ &$0.15001$\\
\addlinespace
Class 5 & $0.08$ &$9.875\times10^{-4}$ &$3.090\times10^{-2}$ &$0.39991$\\
        & $0.02$ &$6.179\times10^{-5}$ &$7.545\times10^{-3}$ &$0.37683$\\
        & $0.005$ &$3.862\times10^{-6}$ &$1.876\times10^{-3}$ &$0.37130$\\
\bottomrule
\end{tabular}
\end{table}

$\eps=0.02,0.01,0.005$ の3点に対する対数回帰では、$e_{\rm exact}$ の収束次数は両模型とも約 $2.000$、$e_C$ は XXX で約 $1.008$、Class 5 で約 $1.004$ である。一方、因子化した処方の誤差は零にならず
\begin{equation}
\lim_{\eps\to0} e_{\rm fact}=\|\Cstar\|_{\rm F}
=\begin{cases}
0.149571& (\alpha=0),\\
0.369473& (\alpha=0.7-0.2\ii)
\end{cases}
\end{equation}
という非零の値に近づく。単に格子を細かくしても、この処方の差は取り除けない。示した収束次数は選んだ滑らかな場での数値的確認であり、全ての初期条件に対する収束定理ではない。

\subsection{結果の範囲と研究の終了条件}

今回の段階で、Class 5 の指定した spin-$\tfrac12$ 表示について、量子 $R$ 行列から $h,\mathcal A$ を構成し、共有格子点上の作用素積から $\Cstar$ を独立に求め、局所補正により時間 Lax 成分を復元した。補正後の曲率は運動方程式と同値であり、XXX 極限でも同じ処方が必要であることを確認した。これは、既知の古典 $V$ を使って差を定義することと、量子側からその差を導くことを区別した検証である。

この結果を述べる際には、次の範囲を保つ必要がある。第一に、$\Cstar=-2\ii\partial_x(U^2)$ は式\eqref{audit5q:eq:U}の規格化と表示に関する式であり、任意の gauge 表示で不変な公式ではない。第二に、coherent-state 期待値の非乗法性を、量子--古典対応そのものの一般的な障害と呼ぶ必要はない。正しい積を保つと、この模型では局所補正に吸収できる。

第三に、spin-$j$ の正規化演算子 $J_a/j$ の二次 moment が $1/(2j)$ の補正を持つことは、一般の時間 Lax 演算子の補正全体が同じ係数で縮小することを直ちには意味しない。時間成分の逆行列や Hamiltonian は表現に依存する。本報告は高 spin の対応を実行しておらず、その一般化を今回の結論に含めない。

第四に、原論文が companion を得られなかったことは文献上の事実であるが\cite{DFN}、その著者らが実際に式\eqref{audit5q:eq:star}の補正を落としたことは示されていない。本報告が特定した失敗は、明示的に定義した $J_{\rm fact}^{(\eps)}$ という処方についてのものである。

Class 5 のこの局所的な量子--古典対応は、今回の段階で閉じた計算として提示できる。Class 5/6 を合わせた報告を完成させる際の次の限定された確認は、Class 6 でも同じ「作用素積から先に補正を求める」手順を実行し、模型依存の部分を区別することである。Class 6 の以前の結果は今回再計算していない。任意の模型の一般定理や soft particle production の説明を、現在の結果の完成条件へ追加する必要はない。投稿上の新規性の評価は、この計算の正しさとは別に文献との式レベルの比較を要する。

% embedded standalone report appendix boundary removed
\subsection{共有格子点の補正の成分表示}
\label{audit5q:app:C}

式\eqref{audit5q:eq:Cfromquantum}を直接評価した成分は、$pq+z^2=1$ と $qp_x+pq_x+2zz_x=0$ を用いて
\begin{align}
(\Cstar)_{11}&=-\frac{\ii}{4\lambda^2}
\bigl[4\alpha^2\lambda^2pp_x+
\alpha\lambda(pz_x+p_xz+p_x)+z_x\bigr],\\
(\Cstar)_{12}&=\frac{\ii}{4\lambda^2}
\bigl[2\alpha^2\lambda^2(pz_x+p_xz+p_x)
+2\alpha\lambda(z+1)z_x-q_x\bigr],\\
(\Cstar)_{21}&=-\frac{\ii}{4\lambda^2}(2\alpha\lambda p+1)p_x,\\
(\Cstar)_{22}&=\frac{\ii}{4\lambda^2}
\bigl[\alpha\lambda(pz_x+p_xz+p_x)+z_x\bigr].
\end{align}
式\eqref{audit5q:eq:U}を二乗して微分すると、これら四成分が $-2\ii\partial_x(U^2)$ と一致する。

この一致は既知の $V$ から逆算したものではない。再現コードでは式\eqref{audit5q:eq:A1}、\eqref{audit5q:eq:A2}の $8\times8$ 行列を作り、四つの連結積を物理空間で trace し、式\eqref{audit5q:eq:Cfromquantum}に代入した結果と比較している。

\subsection{再現手順と検証項目}

報告に添付したコードは Python 3、NumPy、SymPy だけを使用する。今回実行した版は NumPy 2.3.5、SymPy 1.14.0 である。ファイル構成は
\begin{quote}\small\ttfamily
report/class5\_lax\_coherent\_symbol\_ja.tex\\
report/class5\_lax\_coherent\_symbol\_ja.pdf\\
code/verify\_symbolic.py\\
code/verify\_numeric.py\\
results/symbolic\_verification.json\\
results/symbolic\_expressions.json\\
results/numeric\_verification.json\\
results/finite\_spacing.csv
\end{quote}
である。パッケージの root directory から
\begin{quote}\small\ttfamily
python code/verify\_symbolic.py\\
python code/verify\_numeric.py
\end{quote}
を実行する。\texttt{verify\_symbolic.py} は25個の厳密な等式検査を行い、いずれかに失敗すると停止する。検査には Sutherland relation、共有格子点の補正、$V^{\downarrow}=W+2Z$、補正後曲率、Hamiltonian の変分、原論文の運動方程式との辞書が含まれる。

拘束を用いる記号計算では
\begin{align}
pq+z^2-1&=0,\\
qp_x+pq_x+2zz_x&=0,\\
qp_{xx}+pq_{xx}+2p_xq_x+2zz_{xx}+2z_x^2&=0
\end{align}
が生成する多項式 ideal で剰余を取る。係数体は $\mathbb Q(\ii,\lambda,\alpha)$、変数順序は
\[
(p_{xx},q_{xx},z_{xx},p_x,q_x,z_x,p,q,z)
\]
とした lexicographic order である。今回の25検査は全て厳密に零となった。これは数値的な低損失判定ではない。

数値コードは Taylor 展開の係数を再利用せず、有限 $\eps$ の $R$、逆行列、作用素積を直接作る。乱数 seed は20260909に固定した。詳細なパラメータ、誤差、収束次数は JSON と CSV に保存される。LaTeX 原稿は LuaLaTeX で2回コンパイルする。

\clearpage


% Split only for repository maintainability; the three files below form one continuous audit section.
\section{独立監査 B：Class 6 量子作用素積からの時間 Lax 成分}
\subsection{問題設定と今回の計算の範囲}
Class 6は、最近接相互作用を持つ可積分なスピン二分の一の量子鎖の一つである。本報告の文献入力は、de Leeuw--Fontanella--Nieto Garcíaの論文\cite{DFN}の付録Aにある量子R行列と、比較用の連続運動方程式である。Class 6がeleven-vertex模型に対応することは既知であり\cite{CdL}、対応するLandau--Lifshitz型場の理論も既に研究されている\cite{LOZ}。従って「Class 6にLax表示がある」こと自体を新規な主張にはしない。

調べるのは、\emph{有限格子の量子Lax方程式の各作用素積をコヒーレント状態で評価したとき、通常の古典零曲率方程式にどう整理されるか}である。量子鎖から古典連続模型を得る方法論\cite{ADS}や、コヒーレント状態のsymbolの積を扱う方法\cite{KS}も既知である。本報告では、Class 6で必要な項を省かずに計算し、古典時間成分を先に知っていなくても局所補正を構成できるかを検査する。

計算の順序は、量子R行列、量子Hamiltonian、有限格子の時間Lax作用素、作用素積の期待値、連続極限、時間Lax成分の局所補正、独立な運動方程式との照合、の順とする。最後に得た古典時間成分を既知の式から差し引いて補正を\emph{定義}する方法は採らない。

本報告の式はすべてこの文書内で定義する。Class 5との比較は第\ref{audit6q:sec:comparison}節に限り、その比較に必要な規約もそこで再掲する。

\subsection{量子鎖、R行列、連続極限}
\subsubsection{空間と作用素の規約}
各格子点の量子空間を$\mathcal H_n=\mathbb C^2$、Lax作用素が付加的に作用する補助空間を$\mathcal V_0=\mathbb C^2$とする。添字$0$は補助空間、整数$n$は物理的な格子点を表す。二つの空間に作用する行列には、その空間の添字を付ける。例えば$R_{0n}$は$\mathcal V_0\otimes\mathcal H_n$に作用する。

$d$次単位行列を$\one d$、虚数単位を$\ii$、交換子を$[A,B]=AB-BA$と書く。二つの二次元空間を交換する置換作用素$P$と、Class 6の変形を表す行列$K$を、基底$(\uparrow\uparrow,\uparrow\downarrow,\downarrow\uparrow,\downarrow\downarrow)$で
\begin{equation}
P=\begin{pmatrix}1&0&0&0\\0&0&1&0\\0&1&0&0\\0&0&0&1\end{pmatrix},\qquad
K=\begin{pmatrix}0&1&1&0\\0&0&0&-1\\0&0&0&-1\\0&0&0&0\end{pmatrix}
\label{audit6q:eq:PK}
\end{equation}
と定義する。これらは
\begin{equation}
P^2=\one4,\qquad PK=KP=K,\qquad K^3=0
\label{audit6q:eq:PKalgebra}
\end{equation}
を満たす。ここでの$K$は量子空間上の四次行列であり、後に導入する古典スピンの異方性行列とは別物である。

\subsubsection{量子R行列とHamiltonian}
量子スペクトルパラメータを$u$、量子鎖の変形パラメータを$a$とする。文献\cite{DFN}の式(A.1)で交換相互作用の係数$a_1=1$としたR行列は
\begin{align}
R(u;a)&=2u\one4+P+a u(1+2u)K
 +\frac{a^2u^2(1+2u)^2}{2}K^2
\label{audit6q:eq:R}\\
&=\begin{pmatrix}
1+2u&au(1+2u)&au(1+2u)&-a^2u^2(1+2u)^2\\
0&2u&1&-au(1+2u)\\
0&1&2u&-au(1+2u)\\
0&0&0&1+2u
\end{pmatrix}.
\end{align}
以下、微分$\partial_u$では$a$を固定する。最近接の局所Hamiltonian$h$は
\begin{equation}
h=P\,\partial_uR(0;a)=2P+aK,
\qquad H_{\rm lat}=\sum_n h_{n,n+1}
\label{audit6q:eq:hlat}
\end{equation}
である。有限周期鎖を取る場合は添字を周期的に解釈する。時間を$t_{\rm lat}$とし、$\hbar=1$、Heisenberg方程式$\partial_{t_{\rm lat}}O=\ii[H_{\rm lat},O]$を用いる。

式\eqref{audit6q:eq:R}から直接、
\begin{equation}
R(0;a)=P,\qquad
R(u;a)R(-u;a)=(1-4u^2)\one4
\label{audit6q:eq:inverse}
\end{equation}
が得られる。後者は逆行列を与える代数的関係であり、Hermitian共役を含む物理的なユニタリティ条件ではない。逆行列を使う有限格子計算では$u\ne\pm\tfrac12$とする。

\subsubsection{格子間隔に関する展開}
格子間隔を$\epsilon$、位置を$x_n=n\epsilon$とする。連続理論の変形パラメータ$\alpha$、連続スペクトルパラメータ$\lambda\ne0$、連続時間$T$を、
\begin{equation}
a=\epsilon^2\alpha,\qquad u=\frac{\lambda}{\epsilon},\qquad
T=\epsilon^2t_{\rm lat}
\label{audit6q:eq:scaling}
\end{equation}
によって定義する。長さの単位を戻すなら、$u,a$を無次元として$\lambda$は長さ、$\alpha$は長さの逆二乗、$T$は長さの二乗の次元を持つ。

量子空間Lax作用素を$L_n=R_{0n}(u;a)$とする。連続極限で単位行列に近づく規格化を選び、
\begin{equation}
\widehat L_n=\frac{1}{2u}L_n
\label{audit6q:eq:Lnorm}
\end{equation}
と定義する。格子間隔の一次、二次、三次の作用素係数を、それぞれ$\ell_1,\ell_2,\ell_3$と書く。式\eqref{audit6q:eq:R}の代入だけで、厳密な有限多項式
\begin{align}
\widehat L&=\one4+\epsilon\ell_1+\epsilon^2\ell_2+\epsilon^3\ell_3,
\label{audit6q:eq:Lseries}\\
\ell_1&=\frac{P}{2\lambda}+\alpha\lambda K+\alpha^2\lambda^3K^2,
\label{audit6q:eq:l1}\\
\ell_2&=\frac{\alpha}{2}K+\alpha^2\lambda^2K^2,
\qquad
\ell_3=\frac{\alpha^2\lambda}{4}K^2
\label{audit6q:eq:l23}
\end{align}
を得る。Class 6では二次係数$\ell_2$を落とせない。特に
\begin{equation}
\ell_1^2=\frac{1}{4\lambda^2}\one4+2\ell_2
\label{audit6q:eq:l1square}
\end{equation}
が成り立ち、一次係数の作用素としての二乗はスカラー行列ではない。

有限格子の時間成分で必要なスペクトル微分は、規格化した作用素を$u$で微分してから式\eqref{audit6q:eq:scaling}を入れると
\begin{equation}
\left.\partial_u\widehat L\right|_{a\ {\rm fixed}}
=\epsilon^2\partial_\lambda\ell_1
 +\epsilon^3\partial_\lambda\ell_2
 +\epsilon^4\partial_\lambda\ell_3
\label{audit6q:eq:spectralderivative}
\end{equation}
となる。右辺の$\lambda$微分では$\alpha$を固定する。

\subsection{コヒーレント状態と空間Lax行列}
\subsubsection{期待値で何を残すか}
古典スピン場を$\Svec(T,x)=(S^1,S^2,S^3)^{\mathsf T}$とし、単位長の拘束$\Svec^{\mathsf T}\Svec=1$を課す。内積は複素共役を取らない双線形形式である。Pauli行列は
\begin{equation}
\sigma_1=\begin{pmatrix}0&1\\1&0\end{pmatrix},\quad
\sigma_2=\begin{pmatrix}0&-\ii\\\ii&0\end{pmatrix},\quad
\sigma_3=\begin{pmatrix}1&0\\0&-1\end{pmatrix}
\label{audit6q:eq:Pauli}
\end{equation}
とする。スピン・コヒーレント状態の密度行列を
\begin{equation}
\rho(\Svec)=\frac12\left(\one2+\sum_{a=1}^3S^a\sigma_a\right)
\label{audit6q:eq:rho}
\end{equation}
と書く。実単位球面上では通常の純粋状態である。非Hermitianな模型の式を複素化した場へ拡張した場合には、$\rho^2=\rho$と$\tr\rho=1$は保たれるが、正値な密度行列であるとは主張しない。

作用素$B$の上付き矢印$B^\downarrow$は、\emph{物理的な量子空間だけで期待値を取り、補助空間の行列を残す}操作、すなわちlower symbolを意味する。一格子点では
\begin{equation}
B^\downarrow(\Svec)=\tr_n[(\one2\otimes\rho(\Svec))B_{0n}],
\label{audit6q:eq:lower}
\end{equation}
二格子点では
\begin{equation}
B^\downarrow(\Svec_1,\Svec_2)
=\tr_{1,2}[(\one2\otimes\rho(\Svec_1)\otimes\rho(\Svec_2))B]
\label{audit6q:eq:lower2}
\end{equation}
である。$\tr_{1,2}$は物理空間1と2の部分トレースを表す。

異なる物理空間の状態が直積でも、同じ格子点を共有する作用素積について、期待値を先に取った行列の積へ置き換えてはならない。例えば一格子点の作用素を
$A=A_0\otimes\one2+\sum_a A_a\otimes\sigma_a$、
$B=B_0\otimes\one2+\sum_b B_b\otimes\sigma_b$と展開する。ここで$A_a,B_b$は補助空間上の行列である。Kronecker記号を$\delta_{ab}$、$\varepsilon_{123}=1$の完全反対称記号を$\varepsilon_{abc}$とすると、
\begin{equation}
(AB)^\downarrow-A^\downarrow B^\downarrow
=\sum_{a,b=1}^3 A_a B_b
\left(\delta_{ab}-S^aS^b+\ii\sum_{c=1}^3\varepsilon_{abc}S^c\right).
\label{audit6q:eq:star}
\end{equation}
この式はスピン二分の一に対して厳密であり、補助行列の積の順序も保っている。反対称記号$\varepsilon_{abc}$と格子間隔$\epsilon$は別の記号である。

\subsubsection{空間Lax行列の明示式}
複素スピン成分を$S^+=S^1+\ii S^2$、$S^-=S^1-\ii S^2$と定義する。これらは一般の複素化した場では互いの複素共役とは限らず、拘束は$S^+S^-+(S^3)^2=1$となる。

連続空間Lax行列$U(T,x;\lambda)$を、規格化した作用素の一次係数のlower symbolとして
\begin{equation}
U=\ell_1^\downarrow
=\frac1{4\lambda}
\begin{pmatrix}
1+S^3+2\alpha\lambda^2S^+&
S^-+4\alpha\lambda^2S^3-4\alpha^2\lambda^4S^+\\
S^+&1-S^3-2\alpha\lambda^2S^+
\end{pmatrix}
\label{audit6q:eq:U}
\end{equation}
と定める。従って$\widehat L_n^\downarrow=\one2+\epsilon U+O(\epsilon^2)$である。トレースを除いた空間行列を$U_0$と書き、
\begin{equation}
U_0=U-\frac1{4\lambda}\one2
\label{audit6q:eq:U0}
\end{equation}
と定義する。$U$と$U_0$は以下で区別する。

後の式を短くするため、三成分ベクトルをトレース零の二次行列へ写す線形写像$\Qmap$を導入する。任意の$\vvec=(v^1,v^2,v^3)^{\mathsf T}$に対し、$v^\pm=v^1\pm\ii v^2$として
\begin{equation}
\Qmap(\vvec)=\frac1{4\lambda}
\begin{pmatrix}
v^3+2\alpha\lambda^2v^+&v^-+4\alpha\lambda^2v^3-4\alpha^2\lambda^4v^+\\
v^+&-v^3-2\alpha\lambda^2v^+
\end{pmatrix}.
\label{audit6q:eq:Qmap}
\end{equation}
この写像は保存量ではない。特に$U_0=\Qmap(\Svec)$、$\partial_TU=\Qmap(\partial_T\Svec)$である。

\subsection{有限格子の時間成分と必要な展開次数}
\subsubsection{Sutherland relationからの構成}
有限格子の時間Lax作用素は、隣接する二つの物理格子点と補助空間に作用する。まず$L_n=R_{0n}$に対して、式\eqref{audit6q:eq:R}を代入するとSutherland relation
\begin{equation}
[h_{n-1,n},L_nL_{n-1}]
=L_n\partial_uL_{n-1}-(\partial_uL_n)L_{n-1}
\label{audit6q:eq:Sutherland}
\end{equation}
が厳密に成立する。これに基づいて時間作用素$\mathcal A_n$を
\begin{equation}
\mathcal A_n=\ii h_{n-1,n}
-\ii L_n^{-1}h_{n-1,n}L_n-\ii L_n^{-1}\partial_uL_n
\label{audit6q:eq:Atime}
\end{equation}
と取ると、
\begin{equation}
\partial_{t_{\rm lat}}L_n
=\ii[H_{\rm lat},L_n]
=\mathcal A_{n+1}L_n-L_n\mathcal A_n
\label{audit6q:eq:discreteZC}
\end{equation}
を満たす。ここでは有限格子の作用素等式そのものを検査する。

規格化$\widehat L_n=L_n/(2u)$を使うと、共通のスカラー時間成分を除くために
\begin{align}
\widehat{\mathcal A}_n&=\mathcal A_n+\frac{\ii}{u}\one8
\nonumber\\
&=-\ii\widehat L_n^{-1}
\left([h_{n-1,n},\widehat L_n]+\partial_u\widehat L_n\right)
\label{audit6q:eq:Anorm}
\end{align}
と定義できる。添字に対応する恒等作用は省略している。スカラーは両隣で相殺されるため、式\eqref{audit6q:eq:discreteZC}はハット付きの作用素でも成立する。

\subsubsection{格子間隔二次までの時間作用素}
この小節ではテンソル積の順序を$(0,1,2)=($補助空間、左格子点、右格子点$)$に固定する。変形のない局所Hamiltonianを$h_0=2P_{12}$、時間作用素の一次・二次係数を$A_1,A_2$として
\begin{equation}
\widehat{\mathcal A}=\epsilon A_1+\epsilon^2 A_2+O(\epsilon^3)
\label{audit6q:eq:Aseries}
\end{equation}
と展開する。式\eqref{audit6q:eq:Anorm}の直接展開から
\begin{align}
A_1&=-\ii[h_0,\ell_{1,02}],
\label{audit6q:eq:A1}\\
A_2&=\ii\ell_{1,02}[h_0,\ell_{1,02}]
-\ii[h_0,\ell_{2,02}]-\ii\partial_\lambda\ell_{1,02}
\label{audit6q:eq:A2}
\end{align}
となる。$h=h_0+\epsilon^2\alpha K_{12}$なので、Hamiltonianの変形はこの展開では三次以降の時間作用素に入る。その寄与を物理的に無視しているわけではない。


時間の再規格化$T=\epsilon^2t_{\rm lat}$と、空間作用素の一次項$\epsilon U$のため、格子零曲率方程式の三次係数が連続の時間発展を与える。従って、積の三次項には
\begin{equation}
A_2\ell_1\quad\hbox{に加えて}\quad A_1\ell_2
\label{audit6q:eq:neededproducts}
\end{equation}
が必要である。同一の場に評価した時間作用素の三次係数は、隣接する二項の差で相殺される。その係数と空間作用素の一次項との積は四次であり、ここで求める三次係数には不要である。詳細な抽出式を付録\ref{audit6q:app:contraction}に与える。

\subsection{量子作用素積から得られる局所補正}
\subsubsection{時間作用素の期待値と作用素積の差}
連続時間行列の候補を、時間作用素のlower symbolの極限
\begin{equation}
V^\downarrow(T,x;\lambda)
=\lim_{\epsilon\to0}\epsilon^{-2}
\widehat{\mathcal A}_n^\downarrow,
\qquad x=x_n
\label{audit6q:eq:Vdown}
\end{equation}
で定義する。左格子点の場は$\Svec(x-\epsilon)$、右格子点は$\Svec(x)$である。時間作用素の一次項は、同じスピンを両格子点に置くと期待値が消える。このため式\eqref{audit6q:eq:Vdown}は有限になる。

次に、格子零曲率方程式の右辺について、\emph{作用素積のlower symbol}と\emph{lower symbolの行列積}との差を取り、その$\epsilon^3$係数を$\mathcal C$と呼ぶ。すなわち
\begin{align}
\mathcal C=\lim_{\epsilon\to0}\epsilon^{-3}\Big[&
(\widehat{\mathcal A}_{n+1}\widehat L_n)^\downarrow
-(\widehat L_n\widehat{\mathcal A}_{n})^\downarrow
\nonumber\\
&-\widehat{\mathcal A}_{n+1}^{\downarrow}\widehat L_n^\downarrow
+\widehat L_n^\downarrow\widehat{\mathcal A}_{n}^{\downarrow}\Big].
\label{audit6q:eq:Cdefinition}
\end{align}
これは、単に期待値を因子化したために連続方程式から失われる局所行列である。正しい作用素積を使った連続の右辺は
\begin{equation}
\partial_x V^\downarrow+[V^\downarrow,U]+\mathcal C
\label{audit6q:eq:rawcontinuum}
\end{equation}
となる。

この差を式\eqref{audit6q:eq:A1}--\eqref{audit6q:eq:A2}と部分トレースから計算した。結果を表示するために、一次の空間作用素係数の二乗に関する差を$\Xi$と定義する：
\begin{align}
\Xi&=(\ell_1^2)^\downarrow-(\ell_1^\downarrow)^2
\label{audit6q:eq:Xi}\\
&=\frac1{4\lambda^2}\one2+2\ell_2^\downarrow-U^2.
\label{audit6q:eq:Xi2}
\end{align}
$\Xi$は補助空間上の行列であって、正値な分散を意味しない。式\eqref{audit6q:eq:Xi2}は作用素恒等式\eqref{audit6q:eq:l1square}による。

\emph{独立に計算した式\eqref{audit6q:eq:Cdefinition}の結果}が、次の形に整理されることを確認した：
\begin{equation}
\mathcal C=2\ii\left(\partial_x\Xi+[\Xi,U]\right).
\label{audit6q:eq:mainC}
\end{equation}
この恒等式の検査が、既知の古典時間成分を使わない導出の中心である。$\Xi$から右辺を作って正しいと仮定したのではなく、8行8列の行列の作用素積を縮約して得た左辺と、拘束面上で比較した。

\subsubsection{時間Lax成分の復元}
時間Lax成分に加える局所行列を$\Delta V$とし、
\begin{align}
\Delta V&=2\ii\Xi-\frac{\ii}{4\lambda^2}\one2
\label{audit6q:eq:DeltaXi}\\
&=4\ii\ell_2^\downarrow-2\ii U^2+\frac{\ii}{4\lambda^2}\one2
\label{audit6q:eq:Delta}
\end{align}
と取る。最後のスカラー項は$T,x$に依存せず、最終的な時間行列のトレースを零にする規格化である。式\eqref{audit6q:eq:mainC}から
\begin{equation}
\mathcal C=\partial_x\Delta V+[\Delta V,U]
\label{audit6q:eq:absorb}
\end{equation}
が従う。従って、補正した連続時間Lax行列を
\begin{equation}
V=V^\downarrow+\Delta V
\label{audit6q:eq:Vcorrected}
\end{equation}
とすれば、式\eqref{audit6q:eq:rawcontinuum}は$\partial_xV+[V,U]$にまとまる。この操作は、指定した空間行列$U$に対する時間成分の修正である。場のゲージ変換と同一視してはいない。

補助空間の二成分ベクトルを$\psi(T,x;\lambda)$とすると、得られた二つの行列は補助線形問題
\begin{equation}
\partial_x\psi=U\psi,\qquad \partial_T\psi=V\psi
\label{audit6q:eq:linearproblem}
\end{equation}
に現れる。その両微分の両立条件が$\partial_TU-\partial_xV+[U,V]=0$であり、次節で運動方程式との同値性を確認する。

スピンの外積を通常の三次元外積とし、$\Svec_x=\partial_x\Svec$と書く。部分トレースの結果を明示的に整理すると、
\begin{align}
V^\downarrow&=-2\Qmap(\Svec\times\Svec_x)
-2\ii\partial_\lambda U_0+\ii\alpha S^+S^3\one2,
\label{audit6q:eq:Vdownexplicit}\\
\Delta V&=\ii\partial_\lambda U_0-\ii\alpha S^+S^3\one2,
\label{audit6q:eq:Deltaexplicit}\\
V&=-2\Qmap(\Svec\times\Svec_x)-\ii\partial_\lambda U_0,
\qquad \tr V=0.
\label{audit6q:eq:Vfinal}
\end{align}
ここで$\partial_\lambda$はスピン場と$\alpha$を固定した\emph{スペクトルパラメータ微分}であり、時空微分ではない。線形写像$\Qmap$は式\eqref{audit6q:eq:Qmap}で既に明示したので、式\eqref{audit6q:eq:Vfinal}は任意のスピンとその一階空間微分から評価できる完全な式である。

\subsubsection{空間的に一様な場にも残る補正}
式\eqref{audit6q:eq:mainC}を空間行列の係数に戻すと、
\begin{equation}
\mathcal C=-2\ii\partial_x(U^2)
+4\ii\left(\partial_x\ell_2^\downarrow+[\ell_2^\downarrow,U]\right)
\label{audit6q:eq:sourceexpanded}
\end{equation}
となる。Class 6では第二項があり、特に交換子が零とは限らない。

具体的に、ある時刻の一様な配置$\Svec=(0,0,1)^{\mathsf T}$、$\Svec_x=0$を取る。このとき
\begin{equation}
U=\begin{pmatrix}1/(2\lambda)&\alpha\lambda\\0&0\end{pmatrix},\qquad
\Delta V=\begin{pmatrix}-\ii/(4\lambda^2)&\ii\alpha\\0&\ii/(4\lambda^2)\end{pmatrix},
\label{audit6q:eq:uniformmatrices}
\end{equation}
であり、
\begin{equation}
\partial_x\Delta V=0,\qquad
\mathcal C=[\Delta V,U]
=\begin{pmatrix}0&-\ii\alpha/\lambda\\0&0\end{pmatrix}.
\label{audit6q:eq:uniformC}
\end{equation}
$\alpha\ne0$なら補正は有限に残る。これは\emph{一様な瞬間的配置}による検査であって、この場が定常解や真空であるとは述べていない。

もし$4\ii\ell_2^\downarrow$を式\eqref{audit6q:eq:Delta}から落とせば、残りは$U$の多項式となる。一様な配置では空間微分も$U$との交換子も零になるため、式\eqref{audit6q:eq:uniformC}を復元できない。この例は、Class 5の簡略化した補正式をClass 6へそのまま適用できないことを直接示す。

\subsection{Hamiltonianと運動方程式による独立検証}
\subsubsection{古典Hamiltonianと規約}
量子Hamiltonianとは別の記号として、連続Hamiltonianを$H_{\rm cl}$、その密度を$\mathcal H$とする。隣接場$\Svec(x)$と$\Svec(x+\epsilon)$を式\eqref{audit6q:eq:hlat}に入れると、単位スピンの拘束のもとで
\begin{equation}
h_{n,n+1}^\downarrow=2+\epsilon^2\mathcal H+O(\epsilon^3),\qquad
\mathcal H=-\frac12\Svec_x^{\mathsf T}\Svec_x+\alpha S^+S^3.
\label{audit6q:eq:Hsymbol}
\end{equation}
格子点数を$N_{\rm site}$とすれば、一定の空間長で
$H_{\rm cl}=\lim_{\epsilon\to0}(H_{\rm lat}^\downarrow-2N_{\rm site})/\epsilon
=\int dx\,\mathcal H$となる。変分では周期境界、または変分による境界項が消える条件を取る。

古典スピンの異方性を表す三次行列を$N$と定義する：
\begin{equation}
N=\begin{pmatrix}0&0&-1\\0&0&-\ii\\-1&-\ii&0\end{pmatrix},\qquad N^3=0.
\label{audit6q:eq:N}
\end{equation}
すると
\begin{equation}
H_{\rm cl}=\int dx\left[-\frac12\Svec_x^{\mathsf T}\Svec_x
-\frac\alpha2\Svec^{\mathsf T}N\Svec\right].
\label{audit6q:eq:Hcl}
\end{equation}
この符号は量子Hamiltonianと時間の規格化\eqref{audit6q:eq:hlat}、\eqref{audit6q:eq:scaling}から固定している。時間Lax成分との照合では、このHamiltonianの符号と時間の規格化を同時に固定する。

Poisson括弧を
\begin{equation}
\{S^a(x),S^b(y)\}=2\sum_{c=1}^3\varepsilon_{abc}S^c(x)\delta(x-y)
\label{audit6q:eq:PB}
\end{equation}
とする。ここで$\delta(x-y)$はDiracのデルタ関数である。式\eqref{audit6q:eq:Hcl}を変分し、$\partial_T\Svec=\{\Svec,H_{\rm cl}\}$を計算すると、
\begin{equation}
\partial_T\Svec=-2\Svec\times(\Svec_{xx}-\alpha N\Svec)
\label{audit6q:eq:EOM}
\end{equation}
となる。一方、量子作用素積から得た連続右辺\eqref{audit6q:eq:rawcontinuum}も、
$\Qmap[-2\Svec\times(\Svec_{xx}-\alpha N\Svec)]$
と一致する。これは補正後の時間Lax成分を求めた後に行う独立な照合である。

\subsubsection{零曲率方程式と運動方程式の同値性}
運動方程式の残差ベクトルを$\boldsymbol{\mathcal E}$、Lax曲率を$F_{Tx}$と定義する：
\begin{align}
\boldsymbol{\mathcal E}&=\partial_T\Svec+2\Svec\times(\Svec_{xx}-\alpha N\Svec),
\label{audit6q:eq:Eres}\\
F_{Tx}&=\partial_TU-\partial_xV+[U,V].
\label{audit6q:eq:curvature}
\end{align}
運動方程式を場へ代入することなく、
\begin{equation}
F_{Tx}=\Qmap(\boldsymbol{\mathcal E})
\label{audit6q:eq:offShell}
\end{equation}
が成立する。ここで使う空間的な拘束は
\begin{equation}
\Svec^{\mathsf T}\Svec=1,\qquad
\Svec^{\mathsf T}\Svec_x=0,\qquad
\Svec^{\mathsf T}\Svec_{xx}=-\Svec_x^{\mathsf T}\Svec_x
\label{audit6q:eq:kinematic}
\end{equation}
だけである。この意味で式\eqref{audit6q:eq:offShell}はoff shellの恒等式である。

曲率の行列要素を$F_{ij}=(F_{Tx})_{ij}$、残差の複素成分を$\mathcal E^\pm=\mathcal E^1\pm\ii\mathcal E^2$と書くと、残差は
\begin{align}
\mathcal E^+&=4\lambda F_{21},\nonumber\\
\mathcal E^3&=4\lambda(F_{11}-2\alpha\lambda^2F_{21}),\nonumber\\
\mathcal E^-&=4\lambda(F_{12}-4\alpha\lambda^2F_{11}+12\alpha^2\lambda^4F_{21})
\label{audit6q:eq:inverseQ}
\end{align}
と復元できる。成分の順序を$(\mathcal E^+,\mathcal E^-,\mathcal E^3)$から$(F_{21},F_{11},F_{12})$へ固定した係数行列の行列式は$-1/(64\lambda^3)$である。従って$\lambda\ne0$において
\begin{equation}
F_{Tx}=0\quad\Longleftrightarrow\quad\boldsymbol{\mathcal E}=0.
\label{audit6q:eq:equivalence}
\end{equation}
単位スピンに独立な運動方程式は接平面の二成分で足りるが、式\eqref{audit6q:eq:inverseQ}は三成分表示でも情報を失わないことを示している。これだけを、保存量の相互のPoisson可換性や量子可積分性の新たな証明と言い換えることはしない。

\subsubsection{文献の連続運動方程式との対応}
文献\cite{DFN}の式(A.11)--(A.12)のスピン場、時間、異方性行列を、この小節に限り$\Svec_{\rm ref}$、$t_{\rm ref}$、$N_{\rm ref}$と表す。具体的には
\begin{equation}
N_{\rm ref}=\begin{pmatrix}0&0&1\\0&0&-\ii\\1&-\ii&0\end{pmatrix},\qquad
\partial_{t_{\rm ref}}\Svec_{\rm ref}
=\Svec_{\rm ref}\times(\partial_x^2\Svec_{\rm ref}-\alpha N_{\rm ref}\Svec_{\rm ref}).
\label{audit6q:eq:refEOM}
\end{equation}
定数の反射行列$O=\operatorname{diag}(-1,1,1)$によって
\begin{equation}
\Svec_{\rm ref}=O\Svec,\qquad t_{\rm ref}=2T,\qquad
N=O N_{\rm ref} O
\label{audit6q:eq:dictionary}
\end{equation}
と対応させると式\eqref{audit6q:eq:EOM}と一致する。$\det O=-1$のため外積の符号も変わる。Hamiltonianの符号、Poisson括弧の規格化、時間の倍率を一緒に扱う必要がある。これは局所的な運動方程式の対応であり、実条件・境界条件を含めた大域的な量子理論の同値性ではない。

\subsection{記号計算と有限行列による検算}
\subsubsection{厳密な等式検査}
再現用プログラム\path{code/verify_symbolic.py}は、R行列から計算を開始し、保存済みの古典時間行列を読み込まない。R行列の正則性、Hamiltonianの微分表示、逆行列関係、Sutherland relation、空間・時間作用素の展開、共有格子点の積の寄与、局所補正、Hamiltonian変分、off-shell曲率と残差の相互変換など、\textbf{40項目}の等式を厳密に検査し、すべて零になった。

拘束面上の等式は、式\eqref{audit6q:eq:kinematic}の三多項式で還元して判定する。コード中だけの略記として$p=S^+$、$q=S^-$、$z=S^3$を用いる。これらは運動量でもスペクトルパラメータでもない。多項式は
\begin{align}
pq+z^2-1,&\nonumber\\
qp_x+pq_x+2zz_x,&\nonumber\\
qp_{xx}+pq_{xx}+2p_xq_x+2zz_{xx}+2z_x^2&
\label{audit6q:eq:ideal}
\end{align}

であり、係数体を$\mathbb Q(\ii,\lambda,\alpha)$とするGröbner基底で還元する。変数の順序とすべての検査名は\path{results/symbolic_verification.json}に保存した。式\eqref{audit6q:eq:uniformC}を回収できない誤った補正も、零でない対照として記録した。

\subsubsection{独立した数値実装}
\path{code/verify_numeric.py}では、NumPyの有限行列積と逆行列を用い、記号計算の縮約結果を読み込まずに量子格子の等式を検査する。乱数seedは20260909である。行列ノルムにはFrobeniusノルム$\|M\|_{\rm F}=(\sum_{ij}|M_{ij}|^2)^{1/2}$を用いる。二つの行列$A,B$の相対誤差は
\begin{equation}
e_{\rm rel}(A,B)=\frac{\|A-B\|_{\rm F}}{\max(1,\|A\|_{\rm F},\|B\|_{\rm F})}
\label{audit6q:eq:relerr}
\end{equation}
と定義した。

\begin{table}[htbp]
\centering\small
\begin{tabular}{L{82mm}rr}
\toprule
検査&例数&最大相対誤差\\\midrule
有限格子の作用素零曲率方程式&32&$5.02\times10^{-15}$\\
量子Yang--Baxter方程式&32&$4.46\times10^{-16}$\\
off-shell曲率の残差への因数分解&128&$1.18\times10^{-15}$\\
有限格子Hamiltonianのsymbolとの一致&12&$7.68\times10^{-12}$\\
\bottomrule
\end{tabular}
\caption{今回の実行で得た検算結果。最下段は連続極限の誤差ではなく、同じ有限格子間隔で二つの評価法を比べた誤差である。}
\label{audit6q:tab:checks}
\end{table}

最下段の照合内容を明示する。中央と左右のスピンを$\boldsymbol s_0=\Svec(x)$、$\boldsymbol s_{\rm L}=\Svec(x-\epsilon)$、$\boldsymbol s_{\rm R}=\Svec(x+\epsilon)$とする。量子Hamiltonianの可換子を直積状態で評価して得る、連続時間単位での瞬間的なスピンの変化率は
\begin{equation}
\boldsymbol f_\epsilon=-2\boldsymbol s_0\times\left[
\frac{\boldsymbol s_{\rm R}+\boldsymbol s_{\rm L}-2\boldsymbol s_0}{\epsilon^2}
-\frac\alpha2 N(\boldsymbol s_{\rm R}+\boldsymbol s_{\rm L})\right].
\label{audit6q:eq:finiteflow}
\end{equation}
式\eqref{audit6q:eq:Lseries}のlower symbolをこの変化率で微分した値と、有限格子の作用素積を直接部分トレースした値が一致する。$\epsilon^{-3}$で規格化した差を計算するため、最小の格子間隔では丸め誤差が増幅される。

重要な限定として、式\eqref{audit6q:eq:finiteflow}は直積コヒーレント状態で評価した\emph{瞬間的なベクトル場}である。厳密な量子時間発展が直積コヒーレント状態の集合を保つことは仮定も証明もしていない。

\subsubsection{有限格子間隔での収束と因子化誤差}
連続極限の数値検査には、滑らかな実単位スピン
\begin{align}
\Svec(x)&=(\sin\theta(x)\cos\phi(x),\sin\theta(x)\sin\phi(x),\cos\theta(x))^{\mathsf T},\nonumber\\
\theta(x)&=0.9+0.2\sin x,\qquad
\phi(x)=0.3+0.4x+0.1\cos(2x)
\label{audit6q:eq:profile}
\end{align}
を使う。評価位置は$x=0.37$、スペクトルパラメータは$\lambda=0.8+0.2\ii$とし、Class 6では$\alpha=0.7-0.2\ii$、変形のない等方的Heisenberg（XXX）鎖の対照では$\alpha=0$とする。

次の二つの右辺を、同じ有限格子間隔$\epsilon$で計算する：
\begin{align}
J_{\rm exact}(\epsilon)&=\epsilon^{-3}\left[
(\widehat{\mathcal A}_{n+1}\widehat L_n)^\downarrow
-(\widehat L_n\widehat{\mathcal A}_n)^\downarrow\right],\nonumber\\
J_{\rm fact}(\epsilon)&=\epsilon^{-3}\left[
\widehat{\mathcal A}_{n+1}^\downarrow\widehat L_n^\downarrow
-\widehat L_n^\downarrow\widehat{\mathcal A}_n^\downarrow\right].
\label{audit6q:eq:Jnumerical}
\end{align}
前者は作用素積を保つ処方、後者は期待値を因子化した処方である。比較対象$\partial_TU$は式\eqref{audit6q:eq:EOM}によって計算する。三つの\emph{絶対誤差}を
\begin{align}
e_0&=\|J_{\rm exact}-\partial_TU\|_{\rm F},\nonumber\\
e_{\mathcal C}&=\|J_{\rm exact}-J_{\rm fact}-\mathcal C\|_{\rm F},\nonumber\\
e_{\rm fact}&=\|J_{\rm fact}-\partial_TU\|_{\rm F}
\label{audit6q:eq:spacingerrors}
\end{align}
と定義する。さらに、空間作用素の二次・三次係数を残した比較対象
\begin{equation}
J_{\rm target}=\partial_TU+\epsilon\partial_T(\ell_2^\downarrow)
+\epsilon^2\partial_T(\ell_3^\downarrow),\qquad
e_{\rm full}=\|J_{\rm exact}-J_{\rm target}\|_{\rm F}
\label{audit6q:eq:fullerror}
\end{equation}
も評価する。式\eqref{audit6q:eq:fullerror}の各時間微分には同じ連続運動方程式を用いる。これは$e_0$と\emph{異なる比較対象}である。

\begin{table}[htbp]
\centering\small
\begin{tabular}{lrrrrr}
\toprule
模型&$\epsilon$&$e_0$&$e_{\rm full}$&$e_{\mathcal C}$&$e_{\rm fact}$\\\midrule
Class 6&0.080&$6.43\times10^{-2}$&$1.31\times10^{-3}$&$4.65\times10^{-2}$&0.7522\\
&0.020&$1.60\times10^{-2}$&$7.86\times10^{-5}$&$1.15\times10^{-2}$&0.7260\\
&0.005&$3.99\times10^{-3}$&$4.86\times10^{-6}$&$2.85\times10^{-3}$&0.7196\\\midrule
XXX&0.080&$6.21\times10^{-4}$&$6.21\times10^{-4}$&$8.84\times10^{-3}$&0.1572\\
&0.020&$3.89\times10^{-5}$&$3.89\times10^{-5}$&$2.12\times10^{-3}$&0.1514\\
&0.005&$2.43\times10^{-6}$&$2.43\times10^{-6}$&$5.24\times10^{-4}$&0.1500\\\bottomrule
\end{tabular}
\caption{有限格子間隔での絶対誤差。全六段階の格子間隔のデータは添付CSVにある。}
\label{audit6q:tab:spacing}
\end{table}

Class 6では小さい方の三つの格子間隔から求めたべき指数が、$e_0$で1.001、$e_{\rm full}$で2.008、$e_{\mathcal C}$で1.002となった。$e_0$の一次収束は、式\eqref{audit6q:eq:Lseries}の二次係数によって、$\epsilon^{-1}\partial_T\widehat L^\downarrow$自体に$O(\epsilon)$の項があるためである。Class 6の$e_0$が二次収束するとは報告しない。

一方、因子化した処方の誤差は零へ向かわず、
\begin{equation}
\lim_{\epsilon\to0}e_{\rm fact}=\|\mathcal C\|_{\rm F}
=\begin{cases}0.7174723&\text{Class 6},\\0.1495710&\text{XXX}\end{cases}
\label{audit6q:eq:limits}
\end{equation}
となる。これらは指定した場・パラメータに対する値であり、模型全体で一律の誤差ではない。$e_{\mathcal C}$の収束は、式\eqref{audit6q:eq:mainC}が実際に作用素積と期待値の積との差を表していることを数値的にも支持する。

\subsection{Class 5との比較と到達点}
\label{audit6q:sec:comparison}
Class 5についての計算報告では、そこで指定した量子R行列の規格化において、空間作用素の一次係数の二乗が$\ell_1^2=\one4/(4\lambda^2)$となった。その結果、局所補正は$-2\ii U^2+\ii\one2/(4\lambda^2)$で足りた。ここでClass 5の$U$は\emph{同報告で定義した}空間行列を指し、Class 6の式\eqref{audit6q:eq:U}とは別の関数である。

両模型の指定した規格化では、作用素の二乗とsymbolの二乗の差$\Xi$を使うと、同じ形$\Delta V=2\ii\Xi-\ii\one2/(4\lambda^2)$に整理できる。しかしClass 6では式\eqref{audit6q:eq:l1square}の$2\ell_2$が残り、式\eqref{audit6q:eq:Delta}の$4\ii\ell_2^\downarrow$が必要になる。二例で共通の式が得られたことから、任意の量子鎖にも成り立つと推論することはできない。スカラー規格化や局所表示を変えた場合にも係数を再計算する必要がある。

今回で、Class 6の量子R行列から有限格子の時間作用素、共有格子点の作用素積、局所補正、古典時間Lax成分、運動方程式との同値性までが、既知の古典時間成分を入力しない順序でつながった。Class 5に対して既に行った計算と合わせて、\emph{指定した二模型の局所的な量子--古典Lax対応を独立に閉じて検算する}という作業範囲は満たした。

論文化に際して残るのは、既存の連続極限の構成法やeleven-vertexの表示との比較を、今回の補正の導出に対して正確に位置づけることである。計算が成立したことと、導出・解釈が文献上新しいことは別に判断しなければならない。原著者が同じ因子化を行ったために時間Lax成分を得られなかった、といった研究過程についての推測はしない。

また、粒子生成の機構、散乱状態のdressing、任意の量子鎖での一般命題は、この計算を完結させるための追加条件ではない。本報告はそれらを扱わず、指定した表示と連続極限における局所Lax構成を結論とする。

% embedded standalone report appendix boundary removed
\subsection{共有格子点の作用素積を縮約する手順}
\label{audit6q:app:contraction}
この付録は、式\eqref{audit6q:eq:mainC}を既知の時間成分から逆算せずに再計算するための手順を与える。物理空間1と2のスピンを$\boldsymbol s_1,\boldsymbol s_2$とする。時間作用素の係数のsymbolを$f_j(\boldsymbol s_1,\boldsymbol s_2)=A_j^\downarrow$（$j=1,2$）と定義する。直接の部分トレースから
\begin{equation}
f_1(\boldsymbol s_1,\boldsymbol s_2)=-2\Qmap(\boldsymbol s_1\times\boldsymbol s_2),
\qquad f_1(\Svec,\Svec)=0
\label{audit6q:eq:f1}
\end{equation}
を得る。

二変数行列関数$G(\boldsymbol s_1,\boldsymbol s_2)$の、第$j$引数について方向$\boldsymbol v$に取る微分を
\begin{equation}
\partial_jG[\boldsymbol v]
=\sum_{a=1}^3v^a\frac{\partial G}{\partial s_j^a}
\label{audit6q:eq:directional}
\end{equation}
と定義する。また$[G]_{\rm same}$は、\emph{微分と行列積を実行した後に}$\boldsymbol s_1=\boldsymbol s_2=\Svec$を代入することを意味する。すると
\begin{equation}
V^\downarrow=[f_2-\partial_1f_1[\Svec_x]]_{\rm same}
\label{audit6q:eq:Vdownextraction}
\end{equation}
である。

この付録内だけで、左右に作用する空間係数を
\begin{equation}
\ell_{\rm L}=\ell_{1,01},\quad\ell_{\rm R}=\ell_{1,02},\qquad
m_{\rm L}=\ell_{2,01},\quad m_{\rm R}=\ell_{2,02}
\label{audit6q:eq:appendixaliases}
\end{equation}
と略記する。$m$は質量ではなく空間作用素の二次係数の別名である。右側の時間作用素を含む積から出る差を$C_{{\rm R},2},C_{{\rm R},3}$、左側の時間作用素を含む積から出る差を$C_{{\rm L},2},C_{{\rm L},3}$とする。末尾の2,3は格子間隔の次数であり、行列要素の添字ではない。それぞれ
\begin{align}
C_{{\rm R},2}&=(A_1\ell_{\rm L})^\downarrow-f_1\ell_{\rm L}^\downarrow,
\nonumber\\
C_{{\rm L},2}&=(\ell_{\rm R}A_1)^\downarrow-\ell_{\rm R}^\downarrow f_1,
\label{audit6q:eq:C2}\\
C_{{\rm R},3}&=(A_2\ell_{\rm L}+A_1m_{\rm L})^\downarrow
-f_2\ell_{\rm L}^\downarrow-f_1m_{\rm L}^\downarrow,
\nonumber\\
C_{{\rm L},3}&=(\ell_{\rm R}A_2+m_{\rm R}A_1)^\downarrow
-\ell_{\rm R}^\downarrow f_2-m_{\rm R}^\downarrow f_1
\label{audit6q:eq:C3}
\end{align}
で計算する。同じ場に評価した二次の差は$[C_{{\rm R},2}-C_{{\rm L},2}]_{\rm same}=0$である。左右の場のTaylor展開を取ると、三次の差は
\begin{equation}
\mathcal C=
[\partial_2C_{{\rm R},2}[\Svec_x]
+\partial_1C_{{\rm L},2}[\Svec_x]
+C_{{\rm R},3}-C_{{\rm L},3}]_{\rm same}
\label{audit6q:eq:Cextraction}
\end{equation}
となる。右側の組では$(\boldsymbol s_1,\boldsymbol s_2)=(\Svec(x),\Svec(x+\epsilon))$、左側の組では$(\Svec(x-\epsilon),\Svec(x))$なので、式\eqref{audit6q:eq:Cextraction}の二つの方向微分は同符号になる。

式\eqref{audit6q:eq:Cextraction}を、式\eqref{audit6q:eq:A1}--\eqref{audit6q:eq:A2}、\eqref{audit6q:eq:lower2}、\eqref{audit6q:eq:l1}--\eqref{audit6q:eq:l23}から評価し、その後に拘束\eqref{audit6q:eq:kinematic}で還元する。得られた行列と$2\ii(\partial_x\Xi+[\Xi,U])$との差が厳密に零であることが、記号計算の\path{source_equals_covariant_second_moment}検査である。

\subsection{再現方法と記号一覧}
\subsubsection{同梱ファイルと実行手順}
検証コードはPython 3.13.5、SymPy 1.14.0、NumPy 2.3.5で実行した。パッケージのルートで、以下を順に実行する。
\begin{verbatim}
python -m pip install -r requirements.txt
python code/verify_symbolic.py
python code/verify_numeric.py
cd report
lualatex -interaction=nonstopmode -halt-on-error \
  class6_lax_coherent_symbol_ja.tex
lualatex -interaction=nonstopmode -halt-on-error \
  class6_lax_coherent_symbol_ja.tex
\end{verbatim}
最後の二回は相互参照と目次を確定するためのコンパイルである。日本語組版にはLuaLaTeX、LuaTeX-ja、Harano Ajiフォントを備えたTeX環境を使う。フォント本体は添付しない。

\path{results/symbolic_verification.json}は厳密検査の結果、
\path{results/symbolic_expressions.json}は導出した各行列の式、
\path{results/numeric_verification.json}は数値条件と最大誤差、
\path{results/finite_spacing.csv}は各格子間隔での絶対誤差を保存する。コード内の$\ell_1,\ell_2,\ell_3$は、それぞれ\texttt{ell,m,n}と表記している。コード中の\texttt{Q}は式\eqref{audit6q:eq:Qmap}の写像であり、\texttt{source}は式\eqref{audit6q:eq:Cextraction}の直接縮約結果である。

\subsubsection{主な記号}
\small
\begin{longtable}{L{26mm}L{103mm}}
\toprule
記号&意味と定義箇所\\\midrule\endfirsthead
\toprule 記号&意味と定義箇所\\\midrule\endhead
\bottomrule\endfoot
$u,a$&量子スペクトルパラメータ、量子鎖の変形パラメータ。式\eqref{audit6q:eq:R}。\\
$\epsilon,x_n$&格子間隔、格子点の位置$x_n=n\epsilon$。\\
$\lambda,\alpha$&連続スペクトルパラメータと変形パラメータ。式\eqref{audit6q:eq:scaling}。\\
$t_{\rm lat},T$&量子格子の時間と連続時間。$T=\epsilon^2t_{\rm lat}$。\\
$P,K$&置換作用素と変形の四次行列。式\eqref{audit6q:eq:PK}。\\
$h,H_{\rm lat}$&量子鎖の局所Hamiltonianと全Hamiltonian。式\eqref{audit6q:eq:hlat}。\\
$\widehat L_n$&単位行列に近づく規格化の量子空間Lax作用素。式\eqref{audit6q:eq:Lnorm}。\\
$\ell_1,\ell_2,\ell_3$&空間作用素の格子間隔一次・二次・三次の係数。式\eqref{audit6q:eq:Lseries}。\\
$\widehat{\mathcal A}_n$&有限格子の規格化した時間Lax作用素。式\eqref{audit6q:eq:Anorm}。\\
$A_1,A_2$&時間作用素の格子間隔一次・二次の係数。式\eqref{audit6q:eq:Aseries}。\\
$\Svec,S^\pm$&古典スピン場と$S^1\pm\ii S^2$。$\Svec^{\mathsf T}\Svec=1$。\\
$\rho,B^\downarrow$&コヒーレント状態の密度行列と、物理空間だけの期待値。式\eqref{audit6q:eq:rho}--\eqref{audit6q:eq:lower2}。\\
$U,U_0$&連続空間Lax行列とそのトレース零の部分。式\eqref{audit6q:eq:U}--\eqref{audit6q:eq:U0}。\\
$\Qmap$&三成分ベクトルをトレース零行列へ写す線形写像。式\eqref{audit6q:eq:Qmap}。\\
$V^\downarrow$&量子時間作用素のlower symbolから直接得る連続行列。式\eqref{audit6q:eq:Vdown}。\\
$\mathcal C$&作用素積を期待値の積へ置き換えると失われる連続の項。式\eqref{audit6q:eq:Cdefinition}。\\
$\Xi$&$(\ell_1^2)^\downarrow-(\ell_1^\downarrow)^2$。式\eqref{audit6q:eq:Xi}。\\
$\Delta V,V$&局所的な時間成分の補正と、補正した時間Lax行列。式\eqref{audit6q:eq:Delta}、\eqref{audit6q:eq:Vcorrected}。\\
$N$&古典スピンの三次の異方性行列。式\eqref{audit6q:eq:N}。$K$とは別物。\\
$H_{\rm cl},\mathcal H$&連続Hamiltonianとその密度。式\eqref{audit6q:eq:Hsymbol}--\eqref{audit6q:eq:Hcl}。\\
$\boldsymbol{\mathcal E},F_{Tx}$&運動方程式の残差ベクトルとLax曲率。式\eqref{audit6q:eq:Eres}--\eqref{audit6q:eq:curvature}。\\
$J_{\rm exact},J_{\rm fact}$&有限格子の右辺を、作用素積を保って／期待値を因子化して評価した行列。式\eqref{audit6q:eq:Jnumerical}。\\
\end{longtable}
\normalsize

\clearpage


% Split only for repository maintainability; the two files below form one continuous audit section.
\section{独立監査 C：Class 5 古典 $r$-matrix / monodromy route}
\subsection{目的と今回の結論}

2026年のClass 5原論文は、量子Lax作用素から古典空間Lax行列を構成し、Kostant--Kirillov型のPoisson bracketと古典$r$-matrixを得た一方、運動方程式を再現するcompanion matrixを見つけられなかったと記している\cite{DFN}。原論文はその代わりにboost法から保存量のtowerを構成して古典可積分性を示した。

先行する我々の計算報告では、原論文とは別の方向から、有限格子の量子零曲率方程式に現れる時間Lax作用素を直接連続極限へ移した。その際、同じ物理格子点を共有する作用素積のlower symbolを因子化すると有限の項を失うことを確認し、その補正を保持することで時間Lax行列を得た。さらに零曲率方程式とHamiltonianからの運動方程式がoff-shellで同値であることを確認した。

残っていた確認は、その時間Lax行列が、古典可積分系で標準的な$r$-matrix／monodromy routeから得る時間成分と一致するかどうかであった。Doikou--Karaiskosの記法では、transfer matrixを$t(\zeta)$、空間Lax行列を$\mathbb U(x,\mu)$とすると、生成Hamiltonian flowは
\begin{equation}
\{\ln t(\zeta),\mathbb U(x,\mu)\}
=
\partial_x\mathbb V(x;\zeta,\mu)
+[\mathbb V(x;\zeta,\mu),\mathbb U(x,\mu)]
\label{audit5r:eq:STSflow}
\end{equation}
と書け、時間側の生成行列$\mathbb V$はmonodromyと古典$r$-matrixから構成される\cite{DK}。この種の関係を以下ではSemenov--Tian--Shansky型の構成と呼ぶ。

今回の結論は、Class 5でこの照合が通った、というものである。二つの時間行列は場に依存しないスカラー行列を除いて一致する。従って、Class 5/6の量子--古典Lax bridgeについて、追加の模型計算を必須条件として残す必要はない。

\subsection{記号と物理的意味}

本報告では、monodromyを展開するためのスペクトルパラメータと、最終的なLax pairのスペクトルパラメータを区別する。この区別をしないと、式の意味が分かりにくくなるためである。

\begin{longtable}{L{39mm}L{105mm}}
\toprule
記号 & 物理的意味と定義\\
\midrule
\endhead
古典spin場 $\Svec(x)$ & 単位球面上の古典spin密度。成分を$S^1,S^2,S^3$とする。\\
複素spin成分 $p,q,z$ & $p:=S^1+\ii S^2$, $q:=S^1-\ii S^2$, $z:=S^3$。従って$\Svec^2=1$は$pq+z^2=1$となる。$p$は運動量ではない。\\
Laxスペクトルパラメータ $\mu$ & 最終的な空間Lax行列$U(x,\mu)$と時間Lax行列$V(x,\mu)$に残るスペクトルパラメータ。\\
生成スペクトルパラメータ $\zeta$ & transfer matrix $t(\zeta)$を展開し、可積分hierarchyの各Hamiltonian flowを取り出すための別のスペクトルパラメータ。\\
coherent-state projector $\rho$ & spin-$\tfrac12$ coherent stateに対応するrank-one行列。式\eqref{audit5r:eq:rho}で定義する。\\
Jordanian有限項 $K$ & Class 5空間Lax行列のうち、$\mu\to0$で$1/\mu$極を持たない変形項。式\eqref{audit5r:eq:UK}で定義する。\\
permutation作用素 $P_{12}$ & 二つの$2$次元補助空間を交換する$4\times4$行列。\\
上向き基底へのprojector $\Pup$ & $\Pup:=\operatorname{diag}(1,0)$。\\
raising matrix $\Eplus$ & $\Eplus:=\begin{psmallmatrix}0&1\\0&0\end{psmallmatrix}$。\\
Jordanian tensor $C_{12}$ & $C_{12}:=\Pup\otimes\Eplus-\Eplus\otimes\Pup$。\\
古典$r$-matrix $r_{12}(\zeta,\mu)$ & 空間Lax行列の線形Poisson algebraを与える$4\times4$行列。式\eqref{audit5r:eq:rcont}で定義する。\\
monodromy $T(x,y;\zeta)$ & $\partial_xT(x,y;\zeta)=U(x,\zeta)T(x,y;\zeta)$を満たすpath-ordered exponential。\\
局所monodromy projector $\Pi(x;\zeta)$ & $\zeta\to0$の局所漸近展開で、$\rho$へ近づくrank-one projector。式\eqref{audit5r:eq:Piexp}で定義する。\\
$r$-matrix routeの時間行列 $V_r$ & transfer matrixの一次のlocal chargeが生成する時間Lax行列。\\
量子極限の時間行列 $V_q$ & 前報告で有限格子の量子time operatorから直接得た連続時間Lax行列。\\
\bottomrule
\end{longtable}

\subsection{Class 5の連続空間Lax行列}

複素spin成分を
\begin{equation}
p=S^1+\ii S^2,\qquad q=S^1-\ii S^2,\qquad z=S^3,
\qquad pq+z^2=1
\label{audit5r:eq:spinconstraint}
\end{equation}
とする。spin-$\tfrac12$ coherent stateのrank-one projectorを
\begin{equation}
\rho(x)=\frac12
\begin{pmatrix}
1+z&q\\
p&1-z
\end{pmatrix},
\qquad
\rho^2=\rho,
\qquad
\tr\rho=1
\label{audit5r:eq:rho}
\end{equation}
と定義する。

前報告で量子$R$行列の長波長極限から得たClass 5の空間Lax行列は
\begin{equation}
U(x,\mu)
=
\frac{1}{4\mu}
\begin{pmatrix}
1+z+2\alpha\mu p&q-2\alpha\mu(1+z)\\
p&1-z
\end{pmatrix}
\label{audit5r:eq:Uexplicit}
\end{equation}
である。ここで連続変形パラメータを$\alpha$とした。

式\eqref{audit5r:eq:Uexplicit}を、小さい$\mu$で発散する部分と有限なJordanian変形に分けるため、有限項$K(x)$を
\begin{equation}
K(x)
:=\frac12
\begin{pmatrix}
p&-(1+z)\\0&0\end{pmatrix}
\label{audit5r:eq:K}
\end{equation}
と定義する。すると
\begin{equation}
U(x,\mu)=\frac{\rho(x)}{2\mu}+\alpha K(x)
\label{audit5r:eq:UK}
\end{equation}
と書ける。この分解が以下の小さい生成スペクトルパラメータ$\zeta$の展開を簡単にする。

Poisson bracketの規格化は前報告と同じく
\begin{equation}
\{p(x),q(y)\}=-4\ii z(x)\delta(x-y),\quad
\{p(x),z(y)\}=2\ii p(x)\delta(x-y),\quad
\{q(x),z(y)\}=-2\ii q(x)\delta(x-y)
\label{audit5r:eq:PBpqz}
\end{equation}
とする。

\subsection{連続古典$r$-matrixの固定}

Class 5原論文では、古典Lax行列がquadraticなKostant--Kirillov型Poisson bracketを満たし、その$r$-matrixを量子$R$行列から与えている\cite{DFN}。本報告では長波長規格化後の空間Lax行列$U$に直接Poisson bracketを取り、連続線形algebraの$r$-matrixを固定する。

まず、二つの補助空間を交換するpermutation作用素を
\begin{equation}
P_{12}=
\begin{pmatrix}
1&0&0&0\\
0&0&1&0\\
0&1&0&0\\
0&0&0&1
\end{pmatrix}
\end{equation}
とする。さらに、上向き基底へのprojectorとraising matrixを
\begin{equation}
\Pup=\begin{pmatrix}1&0\\0&0\end{pmatrix},
\qquad
\Eplus=\begin{pmatrix}0&1\\0&0\end{pmatrix}
\end{equation}
とし、Class 5のJordanian tensorを
\begin{equation}
C_{12}:=\Pup\otimes\Eplus-\Eplus\otimes\Pup
\label{audit5r:eq:Cj}
\end{equation}
と定義する。

このとき、前節のPoisson bracket規約に対応する連続古典$r$-matrixは
\begin{equation}
r_{12}(\zeta,\mu)
=
\frac{\ii}{2(\zeta-\mu)}P_{12}
+\ii\alpha C_{12}
\label{audit5r:eq:rcont}
\end{equation}
である。実際、補助空間1と2に作用する空間Lax行列を
\[
U_1=U\otimes\Id_2,
\qquad
U_2=\Id_2\otimes U
\]
と書くと、成分を直接計算して
\begin{equation}
\{U_1(x,\zeta),U_2(y,\mu)\}
=
[r_{12}(\zeta,\mu),U_1(x,\zeta)+U_2(y,\mu)]\,\delta(x-y)
\label{audit5r:eq:linearPB}
\end{equation}
が成立する。

式\eqref{audit5r:eq:rcont}の第一項は通常のisotropic Landau--Lifshitz模型のrational $r$-matrixに対応し、第二項がClass 5のJordanian変形である。原論文の式(2.36)にも、量子$R$行列のClass 5変形に由来する同じrational部分とJordanian部分が現れる\cite{DFN}。ここでは係数と$\ii$の位置を原論文から読み替えるのではなく、式\eqref{audit5r:eq:PBpqz}と式\eqref{audit5r:eq:Uexplicit}から直接固定した。

\subsection{monodromyから時間成分を生成する標準式}

空間Lax行列$U(x,\zeta)$から、区間$[y,x]$のmonodromyを
\begin{equation}
T(x,y;\zeta)
=
\mathcal P\exp\!\left(\int_y^x U(s,\zeta)\,ds\right)
\label{audit5r:eq:monodromy}
\end{equation}
と定義する。周期系ではtransfer matrixは
\begin{equation}
t(\zeta):=\tr T(L,0;\zeta)
\end{equation}
である。

線形Poisson algebra\eqref{audit5r:eq:linearPB}に対する標準的なmonodromy計算から、生成時間行列は
\begin{equation}
\mathbb V_b(x;\zeta,\mu)
=
t(\zeta)^{-1}
\tr_a\!\left[
T_a(L,x;\zeta)
\,r_{ab}(\zeta,\mu)\,
T_a(x,0;\zeta)
\right]
\label{audit5r:eq:VSTS}
\end{equation}
となる\cite{DK}。添字$a$はtraceを取る補助空間、添字$b$は最終的に$2\times2$行列として残る補助空間を表す。

ここで重要なのは、$\zeta$と$\mu$を区別することである。生成スペクトルパラメータ$\zeta$について$\ln t(\zeta)$と$\mathbb V(x;\zeta,\mu)$を展開し、その各係数が別々のHamiltonian flowを与える。一方、Laxスペクトルパラメータ$\mu$はそのflowのLax pairに残る。

\subsection{小さい生成スペクトルパラメータの局所projector}

Class 5の空間Lax行列は
\[
U(x,\zeta)=\frac{\rho(x)}{2\zeta}+\alpha K(x)
\]
なので、$\zeta\to0$ではrank-one projector $\rho$に比例する$1/\zeta$極が支配的である。monodromyのlocal chargeを得る標準的な漸近展開では、このleading eigenspaceに対応する局所rank-one projectorを用いる。

そのprojectorを$\Pi(x;\zeta)$と書き、
\begin{equation}
\Pi^2=\Pi,
\qquad
\tr\Pi=1,
\qquad
\partial_x\Pi=[U(x,\zeta),\Pi]
\label{audit5r:eq:Piconstraints}
\end{equation}
を満たし、$\zeta\to0$で$\rho$へ近づくbranchを選ぶ。展開を
\begin{equation}
\Pi(x;\zeta)=\rho(x)+\zeta\Pi_1(x)+O(\zeta^2)
\label{audit5r:eq:Piexp}
\end{equation}
と定義する。

式\eqref{audit5r:eq:Piconstraints}のprojector条件を$O(\zeta)$まで展開すると
\begin{equation}
\rho\Pi_1+\Pi_1\rho=\Pi_1
\label{audit5r:eq:Pi1projector}
\end{equation}
となる。一方、輸送方程式の$O(\zeta^0)$は
\begin{equation}
\rho_x
=
\frac12[\rho,\Pi_1]
+\alpha[K,\rho]
\label{audit5r:eq:Pi1transport}
\end{equation}
を与える。

$\rho$がrank-one projectorであることを使うと、これら二式の解は
\begin{equation}
\Pi_1
=
2[\rho,\rho_x]
-2\alpha[\rho,[K,\rho]]
\label{audit5r:eq:Pi1}

\end{equation}
である。これは空間微分の寄与とJordanian変形の寄与を明確に分けている。

ここで用いている$\Pi$は、transfer matrixの局所漸近展開を記述するformal projectorである。monodromyの二つの固有値の比に由来する$\exp(-c/\zeta)$型の項は、$\zeta$のpower seriesの全次数を越えるため、ここで取り出すlocal hierarchyの係数には入らない。本報告はその非摂動的な漸近項を主張の対象にしない。

\subsection{$r$-matrix routeから時間Lax行列を取り出す}

式\eqref{audit5r:eq:VSTS}の局所power seriesでは、monodromy挿入を前節の$\Pi$で表せる。permutation作用素について
\begin{equation}
\tr_a(\Pi_a P_{ab})=\Pi_b
\end{equation}
が成立する。また、Jordanian tensor\eqref{audit5r:eq:Cj}から
\begin{equation}
\tr_a(\Pi_a C_{ab})
=
\tr(\Pi\Pup)\Eplus
-\tr(\Pi\Eplus)\Pup
\end{equation}
を得る。

従って、生成時間行列の局所展開は
\begin{equation}
\mathbb V(x;\zeta,\mu)
=
\mathbb V^{(0)}(x;\mu)
+\zeta\mathbb V^{(1)}(x;\mu)
+O(\zeta^2)
\end{equation}
と書け、その一次係数は
\begin{equation}
\mathbb V^{(1)}
=
-\frac{\ii}{2\mu}
\left(\Pi_1+\frac{\rho}{\mu}\right)
+\ii\alpha
\left[
\tr(\Pi_1\Pup)\Eplus
-\tr(\Pi_1\Eplus)\Pup
\right]
\label{audit5r:eq:Vgen1}
\end{equation}
となる。

符号に注意する。標準式\eqref{audit5r:eq:STSflow}はHamiltonianをPoisson bracketの第一引数に置き、
\[
\{H,U\}=\partial_x\mathbb V+[\mathbb V,U]
\]
と書いている。一方、前報告の時間発展は
\[
U_T=\{U,H\}
\]
という規約である。Poisson bracketの反対称性により、同じ物理時間に対応する時間Lax行列は
\begin{equation}
V_r(x;\mu):=-\mathbb V^{(1)}(x;\mu)
\label{audit5r:eq:Vrdef}
\end{equation}
である。ここで添字$r$は古典$r$-matrix routeから得たことを示す。

この段階では、$\ln t(\zeta)$の一次係数をあらかじめClass 5 Hamiltonianと同一視していない。式\eqref{audit5r:eq:Vrdef}が生成するvector fieldを、量子有限格子極限から独立に得たClass 5の時間Lax行列と比較して同定する。

\subsection{量子有限格子極限から得た時間行列との比較}

前報告では、Class 5の時間Lax行列を、空間微分を含む行列$W$と、空間微分を含まない行列$Z$に分けて
\begin{equation}
V_q=W+Z
\label{audit5r:eq:Vq}
\end{equation}
と書いた。ここで$q$はquantum finite-lattice routeを表す添字であり、spin成分$q$とは無関係である。混同を避けるため、本文では必要な箇所で「量子極限の時間行列$V_q$」と併記する。

微分部分$W$は
\begin{align}
W_{11}&=-\frac{\ii}{4\mu}
\left[pq_x-p_xq+4\alpha\mu(zp_x-pz_x)\right],\\
W_{12}&=\frac{\ii}{2\mu}
\left[zq_x-qz_x+\alpha\mu(pq_x-p_xq)\right],\\
W_{21}&=\frac{\ii}{2\mu}(pz_x-zp_x),\\
W_{22}&=\frac{\ii}{4\mu}(pq_x-p_xq)
\label{audit5r:eq:Wcomponents}
\end{align}
であり、代数部分$Z$は
\begin{equation}
Z=2\ii U^2-\frac{\ii}{4\mu^2}\Id_2
\label{audit5r:eq:Z}
\end{equation}
である。

式\eqref{audit5r:eq:Pi1}を式\eqref{audit5r:eq:Vgen1}へ代入し、spin-length constraint
\[
pq+z^2=1
\]
だけを用いて整理すると、今回の中心結果
\begin{equation}
V_r(x;\mu)-V_q(x;\mu)
=
\frac{\ii}{4\mu^2}\Id_2
\label{audit5r:eq:central}
\end{equation}
を得る。

右辺は場$\Svec$にも空間座標$x$にも依存しない。従って
\begin{equation}
\partial_x\left(\frac{\ii}{4\mu^2}\Id_2\right)=0,
\qquad
\left[\frac{\ii}{4\mu^2}\Id_2,U\right]=0
\end{equation}
であり、二つの時間Lax行列は同じ零曲率方程式を与える。

この一致は、Class 5のHamiltonian flowを生成するlocal chargeが、transfer matrixの小$\zeta$展開で一次の係数として現れることも意味する。厳密には、その一次係数と前報告のHamiltonianの差として、固定spin-length leaf上でvector fieldを持たないCasimir、場に依存しない定数、または境界条件で消える全微分を加える自由度が残る。しかし運動方程式と時間Lax成分の照合には影響しない。

\subsection{undeformed極限での確認}

変形パラメータ$\alpha$を0にすると
\[
U=\frac{\rho}{2\mu}
\]
となる。このとき式\eqref{audit5r:eq:Pi1}は
\begin{equation}
\Pi_1=2[\rho,\rho_x]
\end{equation}
へ簡約される。Pauli行列を用い、$\rho=(\Id_2+\Svec\cdot\bm\sigma)/2$と書けば
\begin{equation}
2[\rho,\rho_x]
=
\ii(\Svec\times\Svec_x)\cdot\bm\sigma
\end{equation}
である。

この式を$r$-matrix routeの時間行列へ代入すると、標準的なisotropic Landau--Lifshitzの時間成分が、やはり場に依存しない単位行列成分を除いて得られる。Doikou--Karaiskosの標準Lax pair\cite{DK}とも、Poisson bracketと時間の規格化を合わせれば同じ構造になる。従ってClass 5の一致は、変形項だけの偶然ではない。

\subsection{記号計算による独立検証}

本報告の主要恒等式はSymPyで独立に再計算した。検証コードは本報告と同じディレクトリに保存した。使用した代数的拘束は
\begin{equation}
pq+z^2=1,
\qquad
p_xq+pq_x+2zz_x=0
\label{audit5r:eq:constraints}
\end{equation}
である。第二式は第一式を空間微分した接ベクトル条件である。

検証した項目は次の通りである。
\begin{longtable}{L{69mm}L{70mm}}
\toprule
検査 & 結果\\
\midrule
\endhead
連続$r$-matrixによる線形Poisson algebra、式\eqref{audit5r:eq:linearPB} & 全$4\times4$成分で厳密に一致\\
projector条件$\rho\Pi_1+\Pi_1\rho=\Pi_1$ & 拘束\eqref{audit5r:eq:constraints}の下で厳密に成立\\
projector輸送方程式\eqref{audit5r:eq:Pi1transport} & 拘束\eqref{audit5r:eq:constraints}の下で厳密に成立\\
$r$-matrix routeと量子有限格子routeの差、式\eqref{audit5r:eq:central} & $pq+z^2=1$の下で厳密に成立\\
\bottomrule
\end{longtable}

最後の差は実際に
\begin{equation}
V_r-V_q
=
\begin{pmatrix}
\ii/(4\mu^2)&0\\
0&\ii/(4\mu^2)
\end{pmatrix}
\end{equation}
となる。数値フィッティングではなく、symbolic reductionによる恒等式である。

\subsection{今回の一致が意味すること}

\subsubsection{計算上の意味}

これまでのClass 5計算には、時間Lax行列の正しさを支える三つの独立な経路がある。
\begin{enumerate}
\item 量子有限格子のSutherland relationから時間Lax作用素を構成し、その作用素積を保持して連続極限を取る経路。
\item 連続HamiltonianとPoisson bracketから運動方程式を導き、補正後の零曲率方程式とoff-shellで照合する経路。
\item 本報告で行った、古典$r$-matrixとmonodromy hierarchyから時間Lax成分を生成する経路。
\end{enumerate}

三つ目は一つ目の量子time operatorを入力していない。従って
\begin{equation}
V_{\rm quantum\,limit}
\equiv
V_{r\text{-matrix}}
\quad
\text{mod }f(\mu)\Id_2
\label{audit5r:eq:bridgeclosed}
\end{equation}
が得られたことは、前報告の共有格子点補正と時間成分の独立な検算になっている。

\subsubsection{新規性についての慎重な読み方}

この一致によって、「Class 5の時間Lax成分は古典$r$-matrixから原理的に構成できない」という主張にはならない。むしろ逆で、標準的なmonodromy/$r$-matrix hierarchyを適切な小スペクトルパラメータ展開で実行すれば、同じ時間成分に到達することを今回確認した。

従って投稿時に強調すべき点は、
\begin{quote}
Class 5の古典$r$-matrix一般論そのものではなく、finite-lattice quantum zero-curvature equationのtime operatorから、共有格子点のcoherent-state symbol補正を保持して、同じ古典時間Lax成分へ直接降りる経路を明示したこと
\end{quote}
である。

Class 5原論文がcompanion matrixを得られなかったことは明記している\cite{DFN}が、原著者がどの具体的ansatzやmonodromy展開を試したかは公開情報から分からない。そのため、「原著者がsymbol correctionを落としたため失敗した」とは主張しない。

\subsection{現在の研究の終了判定}

前回の位置づけ報告では、このClass 5の$r$-matrix routeとの式レベルの照合を、Class 5/6 quantum--classical Lax bridgeを閉じるための最後の必須計算とした。本報告でその一致\eqref{audit5r:eq:central}が得られたので、現在の計算課題は一度終了としてよいと判断する。

現時点で完了している内容は次の通りである。
\begin{enumerate}
\item Class 5：量子$R/L$から有限格子time operatorを構成し、coherent-state積補正を含めて連続$U,V$を独立導出した。
\item Class 5：補正後の零曲率方程式とHamiltonian EOMのoff-shell同値性を確認した。
\item Class 6：同じ手順を独立に行い、Class 5の単純な全微分補正をそのまま流用できないことを確認した。
\item Class 5/6：共有格子点補正を二次モーメントsymbol差のadjoint covariant derivativeとして共通記法に整理した。
\item Class 5：古典$r$-matrix／monodromy routeから得る時間Lax行列が、量子有限格子極限の時間Lax行列とスカラー項を除いて一致することを本報告で確認した。
\end{enumerate}

従って、soft particle productionの機構解明、任意spin-$j$への一般化、任意の量子可積分鎖に対する一般定理は、この研究を完了させるための必須条件ではない。それらは独立したfollow-upの候補である。

次に必要なのは新しい計算を増やすことではなく、論文として
\begin{enumerate}
\item どの結果を主定理・主命題に置くか、
\item Class 5とClass 6をどの順番で提示するか、
\item standardな$r$-matrix routeとの差をどの程度本文で説明するか、
\item 「何を初めて示した」と安全に書けるか、
\end{enumerate}
を決めることである。

\subsection{まとめ}

Class 5 Landau--Lifshitz模型について、古典$r$-matrixとmonodromyから時間成分を作る標準的な構成を、前報告の規約に合わせて実行した。連続古典$r$-matrixはrationalなpermutation項とJordanianな定数項からなり、空間Lax行列の線形Poisson algebraを厳密に再現する。

小さい生成スペクトルパラメータ$\zeta$でmonodromyの局所rank-one projectorを展開すると、一次補正$\Pi_1$は式\eqref{audit5r:eq:Pi1}で決まる。transfer matrix hierarchyの一次係数から作る時間Lax行列$V_r$は、量子有限格子のtime operatorを連続極限へ移して得た時間Lax行列$V_q$と
\[
V_r-V_q=\frac{\ii}{4\mu^2}\Id_2
\]
だけ異なる。この差は零曲率方程式に影響しない。

従って、Class 5の時間Lax成分は、量子finite-lattice routeと古典$r$-matrix routeという独立な二方向から同じものとして回収された。前回設定したClass 5/6 quantum--classical Lax bridgeの計算上のゴールは、ここで達成したと判断する。



\section{局所時間Lax方程式における共有サイト補正の一般形}
\label{sec:shared-site-general-proof}

本節では、スピン$1/2$の積状態を用いた長波長極限で、有限格子の時間Lax演算子から
連続時間Lax行列を得る際の補正を解析的に導く。出発点は規格化された有限格子の
零曲率方程式であり、既知の連続時間行列を入力しない。まず成立条件を明示し、
二次係数の空間Taylor展開と三次係数の同一点での値を別々に計算する。
その後、Sutherland relationと空間行列のスピン依存性の階数から、条件の一部が従うことを示す。
ここで得るのは指定した次数の局所恒等式であり、任意の量子鎖の可積分性や有限時間の量子発展の収束を主張するものではない。

\subsection{対象とする極限と成立条件}
\label{subsec:shared-conditions}

局所物理空間を$\mathcal H=\mathbb C^2$、補助空間を有限次元の$\mathcal V_a$とする。
$P=P_{12}$は二つの物理サイトを交換する演算子である。以下、補助空間の添字$a$は省略する。
$X_1$、$X_2$などの添字は物理サイトを指定し、補助空間は共有する。
したがって異なる物理サイトの演算子でも、補助空間での行列積の順序は交換できない。

格子間隔を$\epsilon$、量子および連続理論のスペクトルパラメータを$u,\lambda$とし、
$u=\lambda/\epsilon$と置く。空間Lax演算子、二サイトHamiltonian、スペクトル微分を
\begin{align}
\widehat L&=1+\epsilon X+\epsilon^2Y+O(\epsilon^3),
&h&=2P+\epsilon h^{(1)}+O(\epsilon^2),\label{eq:ss-expansion}\\
\partial_u\widehat L&=\epsilon^2Z+O(\epsilon^3)
\label{eq:ss-spectral-derivative}
\end{align}
と展開する。$X,Y,Z\in\End(\mathcal V_a\otimes\mathcal H)$は位置に明示的に依存しない。
微分は格子結合を固定して取る。連続結合を固定する通常のスケーリングでは$Z=\partial_\lambda X$となるが、
以下の演算子恒等式では、Sutherland relationを課す前は$X,Y,Z$を独立に扱える。

$\widehat L$に対するSutherland relationと、その二次係数の残差を
\begin{align}
[h,\widehat L_2\widehat L_1]
&=\widehat L_2\partial_u\widehat L_1-(\partial_u\widehat L_2)\widehat L_1,
\label{eq:ss-sutherland}\\
\mathcal B&:=[2P,X_2X_1+Y_2+Y_1]+[h^{(1)},X_2+X_1]-Z_1+Z_2
\label{eq:ss-residual}
\end{align}
と定義する。式\eqref{eq:ss-sutherland}が成り立てば$\mathcal B=0$である。
また$P(Y_1+Y_2)P=Y_1+Y_2$より$[P,Y_1+Y_2]=0$である。

各サイトのcoherent-state lower symbolは、rank-one射影
\begin{equation}
\rho(x)=\frac{1+S^i(x)\sigma^i}{2},\qquad
\rho^2=\rho,\quad\tr\rho=1,\quad S^iS^i=1
\label{eq:ss-rho}
\end{equation}
によって取る。$\sigma^i$はPauli行列で、$i=1,2,3$について和を取る。
$\boldsymbol S=(S^1,S^2,S^3)^{\mathsf T}$は単位スピン場である。
まず実スピンで計算し、得られた代数式を$\boldsymbol S^{\mathsf T}\boldsymbol S=1$上へ複素化してもよい。
隣接サイトには滑らかな$\rho(x)$を割り当てる。

一サイト演算子と二サイト演算子の同一点でのsymbolをそれぞれ
\begin{equation}
X^\downarrow:=\tr_{\mathcal H}(\rho X),\qquad
\langle Q\rangle:=\tr_{12}\bigl[(\rho\otimes\rho)Q\bigr]
\label{eq:ss-symbols}
\end{equation}
とする。物理空間のみをトレースし、補助空間の行列を残す。
連続空間行列$U$と第二モーメントの差$\Xi$を
\begin{equation}
U:=X^\downarrow,\qquad \Xi:=(X^2)^\downarrow-U^2
\label{eq:ss-Xi}
\end{equation}
と定義する。演算子積をsymbolの積に置き換えられないこと自体は既知の性質である。
本節で調べるのは、その差が局所時間Lax方程式に現れる次数と形である。

十分条件として、三重項部分空間への射影$P_+=(1+P)/2$に対し
\begin{equation}
P_+h^{(1)}P_+=cP_+
\label{eq:ss-triplet}
\end{equation}
を仮定する。$c$はスピンに依存しない複素スカラーである。
これは同じ向きのスピンの積状態について、一次Hamiltonian補正の期待値が方向に依存しないことを保証する。
本節後半では、$X$のスピン依存性が十分な階数を持つ場合、この条件が$\mathcal B=0$から従うことを示す。

\subsection{時間演算子の展開と二次項の空間変化}

スカラー規格化を固定した時間演算子を
\begin{equation}
\widehat A=-i\widehat L_2^{-1}
\bigl([h,\widehat L_2]+\partial_u\widehat L_2\bigr)
=\epsilon A^{(1)}+\epsilon^2A^{(2)}+O(\epsilon^3)
\label{eq:ss-time}
\end{equation}
とする。その係数は直接展開により
\begin{align}
A^{(1)}&=-i[2P,X_2]=-2i(X_1-X_2)P,\label{eq:ss-A1}\\
A^{(2)}&=iX_2[2P,X_2]-i[2P,Y_2]-i[h^{(1)},X_2]-iZ_2
\label{eq:ss-A2}
\end{align}
となる。$PX_1=X_2P$、$PX_2=X_1P$を用いた。

格子零曲率式の二つの積について、symbolの非乗法性を
\begin{align}
D_n^+&:=(\widehat A_{n+1}\widehat L_n)^\downarrow
       -\widehat A_{n+1}^\downarrow\widehat L_n^\downarrow,\nonumber\\
D_n^-&:=(\widehat L_n\widehat A_n)^\downarrow
       -\widehat L_n^\downarrow\widehat A_n^\downarrow
\label{eq:ss-D}
\end{align}
で表す。それぞれ二次、三次の係数を$D_2^\pm,D_3^\pm$と書く。
$D^+$の物理サイトは$n,n+1$、$D^-$では$n-1,n$である。

同一点では$P(\rho\otimes\rho)=(\rho\otimes\rho)P=\rho\otimes\rho$であるため、
\begin{equation}
\langle A^{(1)}\rangle=0,\qquad
D_2^+(\rho,\rho)=\langle A^{(1)}X_1\rangle=2i\Xi,
\qquad D_2^-(\rho,\rho)=\langle X_2A^{(1)}\rangle=2i\Xi.
\label{eq:ss-D2}
\end{equation}
例えば最初の積は$-2i\langle(X_1-X_2)X_2P\rangle$となり、
$\langle X_1X_2\rangle=U^2$と$\langle X_2^2\rangle=(X^2)^\downarrow$から従う。
これにより$D_n^+-D_n^-$の二次項が消える。

次に空間Taylor係数を求める。評価点$x_0$でのみ$\rho(x_0)=\operatorname{diag}(1,0)$となる
\emph{定数の}物理基底を選ぶ。これは位置依存のゲージ変換ではなく、微分時に基底を動かさない。
$X$の物理行列要素を補助空間の行列$X_{ij}$で表し、
\begin{equation}
X=\sum_{i,j=0}^1X_{ij}\otimes|i\rangle\langle j|,
\qquad
\rho_x:=\partial_x\rho=\begin{pmatrix}0&p\\q&0\end{pmatrix}
\label{eq:ss-blocks}
\end{equation}
と置く。$\rho\rho_x+\rho_x\rho=\rho_x$により対角成分はゼロである。
$p,q$は射影の接ベクトルの成分で、複素化した計算では独立に扱える。
この基底では
\begin{align}
U&=X_{00},\qquad \Xi=X_{01}X_{10},\nonumber\\
\partial_x\Xi&=p(X_{11}-X_{00})X_{10}
             +qX_{01}(X_{11}-X_{00}).
\label{eq:ss-Xi-tangent}
\end{align}
二次項の各引数を別々に微分すると
\begin{align}
\delta_R D_2^+(\rho,\rho)[\rho_x]&=2ip(X_{11}-X_{00})X_{10},\nonumber\\
\delta_L D_2^-(\rho,\rho)[\rho_x]&=2iqX_{01}(X_{11}-X_{00}).
\label{eq:ss-D2-tangent}
\end{align}
$\delta_L,\delta_R$は左または右の物理状態に関する方向微分である。
$D_n^-$を引く負号と$\rho(x-\epsilon)$のTaylor展開の負号により、両者は加算される。
その和は$2i\partial_x\Xi$となる。ここまでは$h^{(1)}$への条件もSutherland relationも必要ない。

\subsection{三次項の導出とSutherland残差}

同一点での三次係数の差は
\begin{equation}
D_3^+-D_3^-=
\langle A^{(2)}X_1-X_2A^{(2)}+A^{(1)}Y_1-Y_2A^{(1)}\rangle
-[\langle A^{(2)}\rangle,U].
\label{eq:ss-D3-start}
\end{equation}
$A^{(1)}$の同一点でのsymbolがゼロであることを使った。
式\eqref{eq:ss-A2}の四つの項に従い、交換相互作用、$Y$、$h^{(1)}$、$Z$の寄与へ分ける。

交換相互作用に由来する寄与は
\begin{align}
(D_3^+-D_3^-)_{\mathrm{ex}}
&=2i\bigl(\langle X_2X_1X_2\rangle-(X^2)^\downarrow U\bigr)
  +2i[\Xi,U]\nonumber\\
&=2i[\Xi,U]-i\langle X_1[2P,X_2X_1]\rangle.
\label{eq:ss-cubic-exchange}
\end{align}
第二行では、物理サイト交換の下で期待値が不変であることを用いた。
$\langle X_2X_1X_2\rangle$は一般に$U^3$ではなく、ここでも積の順序と物理サイトの重なりを保つ。

$Y$に由来する項をすべて合わせると
\begin{equation}
(D_3^+-D_3^-)_{Y}
=-2i\langle([Y_1,X_2]+[X_1,Y_2])P\rangle
=-2i\bigl([Y^\downarrow,U]+[U,Y^\downarrow]\bigr)=0.
\label{eq:ss-Y-cancel}
\end{equation}
時間演算子中の$-i[2P,Y_2]$と、積中の$A^{(1)}Y$をともに保持した後の相殺である。
最初から$Y$を省略してよいという意味ではない。

スペクトル微分による寄与は
\begin{equation}
(D_3^+-D_3^-)_{Z}
=i\bigl((XZ)^\downarrow-UZ^\downarrow\bigr)
=-i\langle X_1(-Z_1+Z_2)\rangle.
\label{eq:ss-Z-part}
\end{equation}

残る$h^{(1)}$項の整理に、式\eqref{eq:ss-triplet}を使う。
一次Hamiltonian補正の交換子のsymbolを
$K_h:=\langle[h^{(1)},X_2]\rangle$と定義する。
この補助空間行列は時間行列にも現れる。同条件の下では
\begin{align}
\langle(X_1+X_2)[h^{(1)},X_1+X_2]\rangle&=0,
\label{eq:ss-h-collective}\\
\langle[h^{(1)},X_2X_1]\rangle&=[K_h,U].
\label{eq:ss-h-product}
\end{align}
式\eqref{eq:ss-h-collective}では、$X_1+X_2$が$P$と可換であり三重項を保存する。
したがって期待値の中で$h^{(1)}$を$c$へ置き換えられる。

式\eqref{eq:ss-h-product}も明示的に示せる。式\eqref{eq:ss-blocks}と同じ基底で、
$h^{(1)}$の物理行列要素を$h_{ij,kl}$とし、
$a_h:=h_{00,01}$、$b_h:=h_{01,00}$と置く。
三重項条件は$h_{00,10}=-a_h$、$h_{10,00}=-b_h$、
$h_{00,11}=h_{11,00}=0$、$h_{00,00}=c$を与える。
従って
\begin{equation}
K_h=a_hX_{10}-b_hX_{01},\qquad
\langle[h^{(1)},X_2X_1]\rangle
=a_h[X_{10},X_{00}]-b_h[X_{01},X_{00}]=[K_h,U].
\label{eq:ss-h-block-proof}
\end{equation}
$a_h,b_h$は物理Hamiltonianのスカラー成分であり、補助空間行列とは可換である。

これら二式と$[h^{(1)},X_2X_1]=[h^{(1)},X_2]X_1+X_2[h^{(1)},X_1]$から
\begin{align}
(D_3^+-D_3^-)_{h}
&=-i\langle[h^{(1)},X_2]X_1-X_2[h^{(1)},X_2]\rangle+i[K_h,U]
\nonumber\\
&=-i\langle X_1[h^{(1)},X_1+X_2]\rangle
\label{eq:ss-h-part}
\end{align}
となる。式\eqref{eq:ss-cubic-exchange}--\eqref{eq:ss-h-part}を加えると
\begin{equation}
D_3^+(\rho,\rho)-D_3^-(\rho,\rho)
=2i[\Xi,U]-i\langle X_1\mathcal B\rangle
\label{eq:ss-cubic-final}
\end{equation}
を得る。これは数値的な係数合わせではなく、交換演算子とrank-one射影の恒等式による導出である。
補助空間に関するCayley--Hamilton恒等式は使用していないため、補助空間の次元を二に限る必要もない。
一方、式\eqref{eq:ss-h-block-proof}では物理空間が二次元であることを使用した。

以上から、連続極限の補正を
$\mathcal C:=\lim_{\epsilon\to0}\epsilon^{-3}(D_n^+-D_n^-)$と定義すると
\begin{equation}
\mathcal C=2i\bigl(\partial_x\Xi+[\Xi,U]\bigr)-i\langle X_1\mathcal B\rangle.
\label{eq:ss-source-final}
\end{equation}
Sutherland relationの該当次数が成立する場合、$\mathcal B=0$であるため
\begin{equation}
\Delta V=2i\Xi+f(\lambda)1,
\qquad \mathcal C=\partial_x\Delta V+[\Delta V,U]
\label{eq:ss-DeltaV}
\end{equation}
となる。$f(\lambda)$は位置とスピンに依存しない任意のスカラーである。

なお三重項条件を課さない場合にも、式\eqref{eq:ss-source-final}に追加する残差は具体的に書ける。
その行列を$\mathcal R_h$とすれば
\begin{equation}
\mathcal R_h=i\left\{
\langle(X_1+X_2)[h^{(1)},X_1+X_2]\rangle
-\langle[h^{(1)},X_2X_1]\rangle+[K_h,U]\right\}.
\label{eq:ss-h-defect}
\end{equation}
任意の$h^{(1)}$について、補正は式\eqref{eq:ss-source-final}の右辺に$\mathcal R_h$を加えたものとなる。
式\eqref{eq:ss-h-collective}と\eqref{eq:ss-h-product}が、この残差を消す条件の働きを明示している。

\subsection{階数条件の下では三重項条件が従う}
\label{subsec:ss-rank}

$h^{(1)}$のサイト交換に対して偶な部分を
$h^{(1)}_+:=(h^{(1)}+Ph^{(1)}P)/2$と定義する。
式\eqref{eq:ss-residual}の交換相互作用項と$-Z_1+Z_2$はサイト交換に対して奇である。
従って
\begin{equation}
\frac{\mathcal B+P\mathcal BP}{2}=[h^{(1)}_+,X_1+X_2].
\label{eq:ss-even-residual}
\end{equation}

$X$を物理空間のPauli基底で展開し、補助空間行列$\mathsf X_0,\mathsf X_i$を
\begin{equation}
X=\mathsf X_0\otimes1+\sum_{i=1}^3\mathsf X_i\otimes\sigma^i,
\qquad U=\mathsf X_0+\sum_{i=1}^3\mathsf X_i S^i
\label{eq:ss-pauli-expansion}
\end{equation}
で定義する。三つの$\mathsf X_i$が線形独立、すなわち
$\boldsymbol v\mapsto\sum_i\mathsf X_i v^i$の階数が三と仮定する。
$\mathcal B=0$と式\eqref{eq:ss-even-residual}から
\begin{equation}
[h^{(1)}_+,\sigma^i_1+\sigma^i_2]=0\qquad(i=1,2,3)
\label{eq:ss-collective-commute}
\end{equation}
が従う。二つのスピン$1/2$の積は、重複度一の三重項と一重項に分解される。
従って全スピンと可換な行列は各部分空間でスカラーとなり、
$P_-=(1-P)/2$を用いて
\begin{equation}
h^{(1)}_+=c_+P_++c_-P_-.
\label{eq:ss-h-even-form}
\end{equation}
よって$P_+h^{(1)}P_+=c_+P_+$であり、式\eqref{eq:ss-triplet}は独立の仮定でなくなる。

この議論は階数が三に満たない場合には使えない。例えば補助空間の非零行列$T_a$を用いた
$X=T_a\otimes\sigma^3$、$Y=Z=0$、$h^{(1)}=\sigma^3\otimes\sigma^3$では
$\mathcal B=0$だが、三重項での$h^{(1)}$はスカラーではない。
これは階数条件なしに三重項条件を導くことへの反例であり、補正式自体への反例とは主張しない。

\subsection{補正後の時間行列と二模型への適用}

同一点の時間symbolも式\eqref{eq:ss-A2}から直接評価できる。
$[2P,Y_2]$の期待値はゼロであるため
\begin{equation}
\langle A^{(2)}\rangle=-2i\Xi-iK_h-iZ^\downarrow.
\label{eq:ss-A2-symbol}
\end{equation}
$\widehat A^\downarrow=\epsilon^2V^\downarrow+O(\epsilon^3)$に含まれる
$A^{(1)}$の空間変化と合わせると、式\eqref{eq:ss-DeltaV}による補正後の時間行列は
\begin{equation}
V=-2(\boldsymbol S\times\partial_x\boldsymbol S)^i
        \frac{\partial U}{\partial S^i}
-iK_h-iZ^\downarrow+f(\lambda)1.
\label{eq:ss-corrected-time}
\end{equation}
$x$微分の次数を数えると、$t=\epsilon^2t_{\rm lat}$の時間規格化では
$\partial_{t_{\rm lat}}\widehat L$も$\epsilon^3$から始まる。
従って共有サイト補正は連続時間方程式と同じ次数で残る。
ただし、左辺の量子交換子のsymbolを古典Hamiltonian流の$\partial_tU$と同定することは別の力学的な確認である。
式\eqref{eq:ss-corrected-time}だけからその確認や、曲率と全運動方程式の同値性までは主張しない。

具体例の名前と演算子を本節内で定義しておく。
Class 5、Class 6は文献\cite{DPR2019,DFN2026}の二つの非Hermitianスピン鎖である。
順序付き基底$\{|00\rangle,|01\rangle,|10\rangle,|11\rangle\}$で
\begin{equation}
K_5=\begin{pmatrix}0&1&-1&0\\0&0&0&0\\0&0&0&0\\0&0&0&0\end{pmatrix},
\qquad
K_6=\begin{pmatrix}0&1&1&0\\0&0&0&-1\\0&0&0&-1\\0&0&0&0\end{pmatrix}
\label{eq:ss-model-K}
\end{equation}
と置く。連続変形結合$\alpha_5,\alpha_6$を用いて、その第一空間係数は
\begin{equation}
X_5=\frac{P}{2\lambda}+\alpha_5K_5,
\qquad
X_6=\frac{P}{2\lambda}+\alpha_6\lambda K_6+\alpha_6^2\lambda^3K_6^2.
\label{eq:ss-model-X}
\end{equation}
Class 5では$h^{(1)}=2\alpha_5K_5$、Class 6では$h^{(1)}=0$であり、
両者とも$Z=\partial_\lambda X$である。
$U_{12},U_{21},U_{11}$の順にスピン3成分への係数を取り出した$3\times3$小行列の行列式は、いずれも
\begin{equation}
\det\frac{\partial(U_{12},U_{21},U_{11})}{\partial(S^1,S^2,S^3)}
=\frac{i}{32\lambda^3}.
\label{eq:ss-model-rank}
\end{equation}
従って$\lambda\ne0$では式\eqref{eq:ss-collective-commute}の議論を適用できる。

Class 6では$X_6^2=1/(4\lambda^2)+2Y_6$という模型固有の恒等式があり、
$\Xi$を$Y_6$で書き換えると第二空間係数に由来する交換子が見える。
一方、一般導出では明示的な$Y$は式\eqref{eq:ss-Y-cancel}で相殺する。
したがって第二係数は途中の計算で保持すべきだが、最終補正に必要な独立の追加データとは限らない。
実際、式\eqref{eq:ss-pauli-expansion}より
\begin{equation}
\Xi=\sum_{i,j=1}^3\mathsf X_i\mathsf X_j
\bigl(\delta_{ij}-S^iS^j+i\varepsilon_{ijk}S^k\bigr),
\qquad\varepsilon_{123}=1
\label{eq:ss-Xi-from-U}
\end{equation}
と復元できる。$\delta_{ij}$はKroneckerのデルタ、$\varepsilon_{ijk}$は完全反対称テンソルであり、$k$についても和を取る。
$U$のスピン依存性全体を知れば、スピン$1/2$ではこの第二モーメントも決まる。

独立な検算として、既知XXZ鎖の弱異方性極限も使える。
格子異方性を$\eta=\epsilon\gamma$とし、連続異方性$\gamma$を固定する。
標準six-vertex $R$行列の規格化により
\begin{equation}
X_{\rm XXZ}=\begin{pmatrix}a&0&0&0\\0&0&b&0\\0&b&0&0\\0&0&0&a\end{pmatrix},
\quad a=\gamma\coth(2\gamma\lambda),\quad b=\frac{\gamma}{\sinh(2\gamma\lambda)}
\label{eq:ss-xxz-X}
\end{equation}
となる。この場合の同じ小行列式は$iab^2/4$であり、$ab\ne0$の領域で階数条件を満たす。
特別なスペクトル点で階数が落ちても、$h^{(1)}=0$なので三重項条件を直接確認できる。
補正式の正しさと、あるスペクトル点で運動方程式を曲率から逆算できるかは区別する必要がある。

\subsection{文献との差、検証範囲、原稿への反映}

Avan--Doikou--Sfetsos\cite{ADS2010}の第2節はcoherent-state期待値から
空間Lax行列とmonodromyを構成し、格子和の数と格子間隔の次数を対応させて連続極限を論じる。
期待値と非線形演算が一般には交換しない点も第2.4節で明記されている。
Doikou--Findlay\cite{DoikouFindlay2017}は量子時間Lax成分の階層を構成し、
第5節でcoherent statesを用いた時間発展にも触れている。
本節の計算は、これらの枠組みやsymbolの非乗法性を新規とするものではない。
比較すべき対象は、局所零曲率式の$\epsilon^3$係数と、その局所的な時間行列への吸収である。
当該文献を確認した範囲では本節と同一の局所補正式を確認できていないが、文献上の優先性は未確定である。

独立検証には二つの実装を用いた。既存の一般成分を用いるプログラムは50項目を再実行した。
追加プログラムは補助空間の成分を非可換な記号として保持し、本節の各相殺、任意の$h^{(1)}$での
式\eqref{eq:ss-h-defect}、Sutherland残差の交換偶成分、全スピンと可換な行列の空間、
二模型とXXZの階数を34項目で確認した。合計84項目のうち、80項目は厳密にゼロ、
4項目は仮定を外した検算で予期した非零結果となった。
この数は独立の物理結果の数ではなく、実装中のassertion数である。

本節の結果に基づけば、原稿には共通補正式の条件と導出を先に置き、二模型をその具体化として扱う余地がある。
ただし、原稿本文の変更は文献比較と独立レビューの後に行う。
高次Hamiltonian、より大きい物理スピン、有限時間の量子発展の誤差制御は本節では扱っていない。
既存のClass 6とeleven-vertex模型の直接対応も、この局所恒等式とは別に照合する必要がある。

\section{Class 6 と既知 eleven-vertex Landau--Lifshitz Lax 対の直接照合}
\label{sec:class6-eleven-vertex-comparison}

Class 6 の量子模型が eleven-vertex model に対応すること、および eleven-vertex $R$-matrix から rational Landau--Lifshitz 型場の理論が構成されること自体は既知である\cite{CdL2024,LOZ2014}。したがって本節の目的は新しい可積分 hierarchy を主張することではない。前節までに有限格子の量子 time-Lax 演算子から得た Class 6 の連続 Lax 対が、既知の eleven-vertex Landau--Lifshitz Lax 対とどの規約変換で一致するかを直接示す。

Levin--Olshanetsky--Zotov の eleven-vertex 構成\cite{LOZ2014}に基づく rational Landau--Lifshitz Lax 対は、Atalikov--Zotov\cite{AtalikovZotov2021}の Sec.~4, Eqs.~(4.1)--(4.5) に標準基底で明示されている。以下ではその式を比較の出発点とする。

\subsection{既知の rational eleven-vertex Landau--Lifshitz Lax 対}

既知模型のスピン行列を
\begin{equation}
\Sigma=
\begin{pmatrix}
\Sigma_{11}&\Sigma_{12}\\
\Sigma_{21}&-\Sigma_{11}
\end{pmatrix}
\label{eq:eleven-sigma-generic}
\end{equation}
とし、既知模型のスペクトルパラメータを$z$とする。空間 Lax 行列は
\begin{equation}
U_{\rm 11v}(z,\Sigma)=
\begin{pmatrix}
\dfrac{\Sigma_{11}}{z}-z\Sigma_{12}
&\dfrac{\Sigma_{12}}{z}\\[2mm]
\dfrac{\Sigma_{21}}{z}-2z\Sigma_{11}-z^3\Sigma_{12}
&-\dfrac{\Sigma_{11}}{z}+z\Sigma_{12}
\end{pmatrix}.
\label{eq:eleven-U}
\end{equation}
これは eleven-vertex classical $r$-matrix から
$U_{\rm 11v}=\operatorname{tr}_2(r_{12}(z)\Sigma_2)$
として得られる行列である\cite{LOZ2014,AtalikovZotov2021}。

同じ規約で classical mechanics の companion matrix は
\begin{equation}
M_{\rm top}(z,\Sigma)=
\begin{pmatrix}
\Sigma_{12}&0\\
2\Sigma_{11}+2z^2\Sigma_{12}&-\Sigma_{12}
\end{pmatrix}.
\label{eq:eleven-Mtop}
\end{equation}
スピン行列の固有値を$\pm\Lambda_{\Sigma}$とし、
$\Sigma^2=\Lambda_{\Sigma}^2\Id_2$とする。空間微分の係数を$k$とすると、
\begin{equation}
\mathsf v
=-\frac{k}{4\Lambda_{\Sigma}^2}[\Sigma,\partial_x\Sigma]
\label{eq:eleven-vaux}
\end{equation}
を用いて time Lax 行列は
\begin{equation}
V_{\rm 11v}(z,\Sigma)
=\frac12\left[
\frac1zU_{\rm 11v}(z,\Sigma)
+2M_{\rm top}(z,\Sigma)
-U_{\rm 11v}(z,\mathsf v)
\right].
\label{eq:eleven-V}
\end{equation}
既知模型の時間を$\tau$とすると、この組は
\begin{equation}
\partial_{\tau}U_{\rm 11v}
+k\partial_xV_{\rm 11v}
-[U_{\rm 11v},V_{\rm 11v}]=0
\label{eq:eleven-zc}
\end{equation}
を満たす\cite{AtalikovZotov2021}。

対応する anisotropy map は
\begin{equation}
J_{\rm 11v}(\Sigma)
=-\begin{pmatrix}
\Sigma_{12}&0\\
2\Sigma_{11}&-\Sigma_{12}
\end{pmatrix},
\label{eq:eleven-J}
\end{equation}
運動方程式は
\begin{equation}
\partial_{\tau}\Sigma
=[J_{\rm 11v}(\Sigma),\Sigma]
-\frac{k^2}{8\Lambda_{\Sigma}^2}
[\Sigma,\partial_x^2\Sigma].
\label{eq:eleven-eom}
\end{equation}

\subsection{Class 6 との規約辞書}

本ノートの Class 6 の古典スピン成分$S^+,S^-,S^3$から
\begin{equation}
\Sigma:=
\begin{pmatrix}
S^3&S^+\\
S^-&-S^3
\end{pmatrix}
\label{eq:class6-eleven-spin}
\end{equation}
を定義する。Class 6 の単位スピン条件
\begin{equation}
S^+S^-+(S^3)^2=1
\label{eq:class6-eleven-constraint}
\end{equation}
より
\begin{equation}
\Sigma^2=\Id_2,
\qquad \Lambda_{\Sigma}^2=1.
\label{eq:class6-eleven-characteristic}
\end{equation}

Class 6 の連続変形結合を$\alpha_6$、本ノートの連続スペクトルパラメータを$\lambda$とする。補助的な非零複素数$\beta$を
\begin{equation}
\beta^2=-2\alpha_6
\label{eq:class6-eleven-beta}
\end{equation}
で定義する。$\beta$の二つの平方根の選択は以下の全式を同時に変換するだけであり、対応する Lax 方程式は同じである。既知模型の規約との辞書は
\begin{equation}
z=\beta\lambda,
\qquad
k=-\frac4\beta,
\qquad
\tau=\frac{i\beta^2}{2}t=-i\alpha_6 t
\label{eq:class6-eleven-dictionary}
\end{equation}
と取れる。

本ノートの Class 6 空間行列$U_6$の traceless part を
\begin{equation}
U_{6,0}:=U_6-\frac{1}{4\lambda}\Id_2
\label{eq:class6-eleven-U0}
\end{equation}
とする。式\eqref{eq:eleven-U}へ
\eqref{eq:class6-eleven-spin}--\eqref{eq:class6-eleven-dictionary}
を代入すると、成分ごとに
\begin{equation}
U_{6,0}(\lambda)
=\frac{\beta}{4}
U_{\rm 11v}(\beta\lambda,\Sigma)^{\mathsf T}
\label{eq:class6-eleven-U-map}
\end{equation}
が成り立つ。右辺の転置$\mathsf T$は補助空間の基底規約の差を表す。本ノートの$U_6$に含まれる$\Id_2/(4\lambda)$は$x,t$に依存せず、交換子にも寄与しないため零曲率式を変えない。

同じ辞書を time matrix に適用する。式\eqref{eq:class6-eleven-characteristic}と$k=-4/\beta$から
\begin{equation}
\mathsf v
=\frac1\beta[\Sigma,\partial_x\Sigma].
\end{equation}
これを式\eqref{eq:eleven-V}へ代入し、本ノートの Class 6 time matrix の各成分と比較すると
\begin{equation}
V_6(\lambda)
=\frac{i\beta^2}{2}
V_{\rm 11v}(\beta\lambda,\Sigma)^{\mathsf T}
=-i\alpha_6
V_{\rm 11v}(\beta\lambda,\Sigma)^{\mathsf T}
\label{eq:class6-eleven-V-map}
\end{equation}
を得る。これはon-shellで零曲率が一致するというだけでなく、$S^\pm,S^3$とそれらの一階空間微分を独立な記号として扱った行列恒等式である。

転置によって交換子の順序が反転することに注意する。式\eqref{eq:eleven-zc}を転置し、式\eqref{eq:class6-eleven-dictionary}--\eqref{eq:class6-eleven-V-map}を使うと
\begin{equation}
\partial_tU_{6,0}
-\partial_xV_6+[U_{6,0},V_6]=0
\end{equation}
となる。ここで
\begin{equation}
\frac{\beta}{4}k=-1
\end{equation}
が空間微分項の符号を合わせる。$U_6-U_{6,0}$はスカラーなので、同じ式は$U_6$にも成立する。

\subsection{運動方程式の直接照合}

Lax 行列だけでなく、spin equation も同じ辞書で一致する。式\eqref{eq:class6-eleven-spin}を式\eqref{eq:eleven-J}へ代入すると
\begin{equation}
J_{\rm 11v}(\Sigma)
=-\begin{pmatrix}
S^+&0\\
2S^3&-S^+
\end{pmatrix}.
\end{equation}
また、式\eqref{eq:class6-eleven-dictionary}から
\begin{equation}
\frac{k^2}{8}=-\frac1{\alpha_6}.
\end{equation}
既知模型の時間$\tau=-i\alpha_6 t$を用いて
式\eqref{eq:eleven-eom}を本ノートの時間$t$へ変換すると
\begin{equation}
\partial_t\Sigma
=-i[\Sigma,\partial_x^2\Sigma]
-i\alpha_6[J_{\rm 11v}(\Sigma),\Sigma].
\label{eq:class6-eleven-matrix-eom}
\end{equation}
その三つの独立成分は
\begin{align}
\partial_t S^+
&=2i\left(
S^+\partial_x^2S^3-S^3\partial_x^2S^+
+\alpha_6(S^+)^2
\right),\\
\partial_t S^-
&=2i\left(
S^3\partial_x^2S^--S^-\partial_x^2S^3
+2\alpha_6(S^3)^2-\alpha_6S^+S^-
\right),\\
\partial_t S^3
&=i\left(
S^-\partial_x^2S^+-S^+\partial_x^2S^-
-2\alpha_6S^+S^3
\right).
\end{align}
これは本ノートで量子 Class 6 Hamiltonian の連続極限から独立に得た運動方程式と一致する。

\subsection{位置づけ}

この直接照合により、Class 6 についても有限格子の量子 time-Lax 演算子から得た$U_6,V_6$が、既知の eleven-vertex Landau--Lifshitz hierarchy の$U,V$と具体的な規約辞書のもとで一致することが分かった。Class 5 で行った classical $r$-matrix / monodromy route との比較と役割が対応している。

一方、この一致は Class 6 の Lax hierarchy の新規性を支持するものではない。むしろ逆で、最終的な classical pair が既知構造に着地することを明示する独立検算である。本研究で別に検討している新しい点は、既知の最終 pair の存在ではなく、finite-lattice quantum time equation から fixed-spin long-wavelength limit を取る際に共有物理サイトの operator product correction がどのように残り、その time component を再構成するかという経路である。

検証プログラム `scripts/class6/verify_eleven_vertex_comparison.py` は、式\eqref{eq:class6-eleven-U-map}、式\eqref{eq:class6-eleven-V-map}、運動方程式、および零曲率規約の変換係数を記号計算で独立に確認する。比較では$S^+,S^-,S^3$と一階・二階空間微分を独立な記号として扱い、最終行列の成分一致に運動方程式を代入していない。

\part{主要記号索引}
\label{part:index}
この part は本文を読み直すときの lookup 用である。同じ文字を複数用途に使わないよう努めたが、文献比較専用の添字・上付きも含むため、最終的な意味をここに再掲する。

\small
\begin{longtable}{L{31mm}L{111mm}}
\toprule
記号 & 意味\\
\midrule\endhead
\bottomrule\endfoot
$\boldsymbol S$ & 文脈で明示した classical spin field。基本拘束は $\boldsymbol S^{\mathsf T}\boldsymbol S=1$。\\
$M$ & Class 5 の反対称 nilpotent matrix。$M\boldsymbol x=\boldsymbol m_5\times\boldsymbol x$。\\
$\boldsymbol m_5$ & Class 5 の complex null vector $(-\ii,-1,0)^{\mathsf T}$。$\boldsymbol m_5^{\mathsf T}\boldsymbol m_5=0$。\\
$K_5$ & Class 5 の $M^2=\boldsymbol m_5\boldsymbol m_5^{\mathsf T}$。rank-one nilpotent で、projector ではない。\\
$\alpha_5$ & Class 5 の continuum coupling。\\
$c$ & $c=\alpha_5/2$。Class 5 field redefinition の係数。\\
$\boldsymbol S_5^{\rm fr}$ & $\boldsymbol S=e^{cxM}\boldsymbol S_5^{\rm fr}$ で定義する field-redefined spin。\\
$z_5$ & Class 5 continuum Lax pair の spectral parameter。\\
$Q_{{\rm EOM},5}$ & Class 5 の $F=QE$ に現れる EOM coefficient matrix。保存量 $Q_n$ とは別物。\\
$u$ & 量子格子の spectral parameter。continuum の $\lambda$、$z_5$、$\zeta_6$ とは式で対応を定義するまで同一視しない。\\
$\epsilon$ & 格子間隔。Levi--Civita symbol は $\varepsilon_{abc}$ と書く。\\
$L_{a,n}(u)$ & 量子格子の空間 Lax 作用素。添字 $a$ は auxiliary space、$n$ は物理格子 site。\\
$A_{a,n}(u)$ & 量子格子の時間 Lax 作用素。独立監査の旧式記録では $\mathcal A_{a,n}$ と書いた箇所もある。\\
$U,V$ & continuum の空間・時間 Lax 行列。\\
$X^\downarrow$ & 量子作用素 $X$ の coherent-state lower symbol。物理空間だけ期待値を取り、auxiliary matrix を残す。\\
$\star$ & $(XY)^\downarrow=X^\downarrow\star Y^\downarrow$ で定義する star product。\\
$\Xi$ & $(\ell_1^2)^\downarrow-(\ell_1^\downarrow)^2$。作用素二乗と lower-symbol 行列積の差。\\
$D_x^{(U)}$ & $D_x^{(U)}X=\partial_xX+[X,U]$。本ノートで用いる adjoint covariant derivative。\\
$N$ & Class 6 の計算基底における三次 nilpotent anisotropy matrix。$N^3=0$。\\
$\zeta_6$ & Class 6 の square-root Lax pair の spectral parameter。\\
$A_{\zeta_6}$ & Class 6 の $A_{\zeta_6}^2=\zeta_6^2\mathbf 1_3+\alpha N$ を満たす三次元行列。量子格子の時間 Lax 作用素 $A_{a,n}$ とは scope が異なる。\\
$\mathcal Q_\lambda$ & Class 6 独立監査で三成分 residual を traceless $2\times2$ 行列へ写す線形写像。\\
$\rho$ & Class 5 $r$-matrix audit で使う rank-one coherent-state projector。\\
$\Pi$ & Class 5 monodromy の小生成スペクトルパラメータ展開に現れる局所 rank-one projector。\\
$\zeta$ & Class 5 STS audit の生成スペクトルパラメータ。\\
$\mu$ & Class 5 STS audit の最終 Lax spectral parameter。\\
$V_r,V_q$ & Class 5 STS audit の $r$-matrix route と quantum finite-lattice route の時間 Lax 行列。\\
\end{longtable}
\normalsize


\clearpage
\begin{thebibliography}{99}

\bibitem{DPR2019}
M.~de Leeuw, A.~Pribytok and P.~Ryan,
``Classifying two-dimensional integrable spin chains,''
J. Phys. A \textbf{52} (2019) 505201,
\href{https://arxiv.org/abs/1904.12005}{arXiv:1904.12005}.

\bibitem{DFN2026}
M.~de Leeuw, A.~Fontanella and J.~M.~Nieto Garc\'ia,
``An integrable deformed Landau--Lifshitz model with particle production?,''
SciPost Phys. \textbf{20} (2026) 041,
\href{https://arxiv.org/abs/2506.13598}{arXiv:2506.13598}.

\bibitem{CdL2024}
M.~Corcoran and M.~de Leeuw,
``All regular $4\times4$ solutions of the Yang--Baxter equation,''
SciPost Phys. Core \textbf{7} (2024) 045,
\href{https://arxiv.org/abs/2306.10423}{arXiv:2306.10423}.

\bibitem{LOZ2014}
A.~Levin, M.~Olshanetsky and A.~Zotov,
``Yang--Baxter equations with two Planck constants,''
Nucl. Phys. B \textbf{887} (2014) 400--434,
\href{https://arxiv.org/abs/1406.2995}{arXiv:1406.2995}.

\bibitem{KY2014}
T.~Kameyama and K.~Yoshida,
``Anisotropic Landau--Lifshitz sigma models from $q$-deformed AdS$_5\times S^5$ superstrings,''
JHEP \textbf{08} (2014) 110,
\href{https://arxiv.org/abs/1405.4467}{arXiv:1405.4467}.

\bibitem{KLS2021}
S.~Krippendorf, D.~L\"ust and M.~Syvaeri,
``Integrability Ex Machina,''
Fortschr. Phys. \textbf{69} (2021) 2100057,
\href{https://arxiv.org/abs/2103.07475}{arXiv:2103.07475}.

\bibitem{ADS2010}
J.~Avan, A.~Doikou and K.~Sfetsos,
``Systematic classical continuum limits of integrable spin chains and emerging novel dualities,''
Nucl. Phys. B \textbf{840} (2010) 469--490,
\href{https://arxiv.org/abs/1005.4605}{arXiv:1005.4605}.

\bibitem{DoikouFindlay2017}
A.~Doikou and I.~Findlay,
``The quantum auxiliary linear problem and the time-like lattice,''
Nucl. Phys. B \textbf{919} (2017) 591--616,
\href{https://arxiv.org/abs/1706.06052}{arXiv:1706.06052}.

\bibitem{SuzukiTsuchiya2017}
H.~Suzuki and Y.~Tsuchiya,
``Coherent-state path integral and Berezin symbols,''
Prog. Theor. Exp. Phys. \textbf{2017} (2017) 043B07.

\bibitem{Schlichenmaier2010}
M.~Schlichenmaier,
``Berezin--Toeplitz quantization for compact K\"ahler manifolds. A review of results,''
Adv. Math. Phys. \textbf{2010} (2010) 927280.

\bibitem{DoikouKaraiskos2011}
A.~Doikou and N.~Karaiskos,
``Generalized Landau--Lifshitz models on the interval,''
Nucl. Phys. B \textbf{853} (2011) 436--460,
\href{https://arxiv.org/abs/1105.5042}{arXiv:1105.5042}.

\bibitem{DFN}
M.~de Leeuw, A.~Fontanella and J.~M.~Nieto Garc\'ia,
``An integrable deformed Landau--Lifshitz model with particle production?,''
SciPost Phys. \textbf{20} (2026) 041,
\href{https://arxiv.org/abs/2506.13598}{arXiv:2506.13598}.

\bibitem{CdL}
M.~Corcoran and M.~de Leeuw,
``All regular $4\times4$ solutions of the Yang--Baxter equation,''
SciPost Phys. Core \textbf{7} (2024) 045,
\href{https://arxiv.org/abs/2306.10423}{arXiv:2306.10423}.

\bibitem{ADS}
J.~Avan, A.~Doikou and K.~Sfetsos,
``Systematic classical continuum limits of integrable spin chains and emerging novel dualities,''
Nucl. Phys. B \textbf{840} (2010) 469--490,
\href{https://arxiv.org/abs/1005.4605}{arXiv:1005.4605}.

\bibitem{KS}
H.~Suzuki and Y.~Tsuchiya,
``Coherent-state path integral and Berezin symbols,''
Prog. Theor. Exp. Phys. \textbf{2017} (2017) 043B07.

\bibitem{LOZ}
A.~Levin, M.~Olshanetsky and A.~Zotov,
``Yang--Baxter equations with two Planck constants,''
Nucl. Phys. B \textbf{887} (2014) 400--434,
\href{https://arxiv.org/abs/1406.2995}{arXiv:1406.2995}.

\bibitem{DK}
A.~Doikou and N.~Karaiskos,
``Generalized Landau--Lifshitz models on the interval,''
Nucl. Phys. B \textbf{853} (2011) 436--460,
\href{https://arxiv.org/abs/1105.5042}{arXiv:1105.5042}.

\end{thebibliography}

\end{document}
